\documentclass[11pt,a4paper]{extbook}

\usepackage{fontspec}
\usepackage{polyglossia}
\setdefaultlanguage{russian}
\setotherlanguage{english}
\usepackage{geometry}
\usepackage{graphicx}
\usepackage{caption}
\usepackage{hyperref}
\usepackage{tocloft}
\usepackage{fancyhdr}
\usepackage{xcolor}
\usepackage{listings}
\usepackage{array}
\usepackage{longtable}
\usepackage{booktabs}
\usepackage{microtype}
\usepackage{amsmath}
\usepackage{amssymb}
\usepackage{amsthm}

% Шрифты с полной кириллицей, доступные практически в любом дистрибутиве
% (Liberation — метрики Times/Helvetica/Courier). Если на машине есть
% Linux Libertine (как в проектах Ринциклопедий), можно вернуть его здесь.
\defaultfontfeatures{Ligatures=TeX}
\setmainfont{Liberation Serif}
\setsansfont{Liberation Sans}
\setmonofont{DejaVu Sans Mono}[Scale=0.85]

\tolerance=2000
\emergencystretch=3em

\geometry{left=2.4cm, right=2.4cm, top=2.5cm, bottom=2.5cm}

\hypersetup{
    colorlinks=true,
    linkcolor=black,
    urlcolor=blue,
    citecolor=black,
    linktoc=all,
}

\usepackage{titlesec}
\titleformat{\chapter}[hang]{\Large\bfseries}{\thechapter.}{18pt}{\Large\bfseries}
\titleformat{\section}[hang]{\large\bfseries}{\thesection.}{12pt}{\large\bfseries}
\titleformat{\subsection}[hang]{\normalsize\bfseries}{\thesubsection.}{10pt}{\normalsize\bfseries}

\setlength{\headheight}{17.0pt}
\pagestyle{fancy}
\fancyhf{}
\lhead{\small\nouppercase{\leftmark}}
\cfoot{\thepage}

% Листинги: обычный код и псевдокод.
\lstdefinestyle{code}{
    basicstyle=\small\ttfamily,
    breaklines=true,
    frame=single,
    numbers=left,
    numberstyle=\tiny\color{gray},
    keywordstyle=\color{blue},
    commentstyle=\color{gray},
    stringstyle=\color{red!70!black},
    showstringspaces=false,
    extendedchars=true,
    columns=fullflexible,
    keepspaces=true,
}
\lstdefinestyle{pseudo}{
    basicstyle=\small\ttfamily,
    breaklines=true,
    frame=leftline,
    framerule=1pt,
    rulecolor=\color{gray!50},
    xleftmargin=1em,
    mathescape=true,
    showstringspaces=false,
    extendedchars=true,
    columns=fullflexible,
    keepspaces=true,
}
\lstset{style=code}

% Диаграммы.
\usepackage{tikz}
\usetikzlibrary{arrows.meta,positioning,fit,backgrounds,calc,shapes.geometric}

% Врезки. Заголовки заключены в {} — иначе запятые ломают ключи tcolorbox.
\usepackage[most]{tcolorbox}
\newtcolorbox{intuition}[1][Интуиция]{colback=gray!5,colframe=gray!45,title={#1},fonttitle=\bfseries,breakable}
\newtcolorbox{pitfall}[1][Ловушка]{colback=red!4,colframe=red!35,title={#1},fonttitle=\bfseries,breakable}
\newtcolorbox{workedbox}[1][Разбор примера]{colback=blue!3,colframe=blue!30,title={#1},fonttitle=\bfseries,breakable}
\newtcolorbox{openq}[1][Открытый вопрос]{colback=orange!6,colframe=orange!60,title={#1},fonttitle=\bfseries,breakable}

% Теоремоподобные среды.
\theoremstyle{definition}
\newtheorem{definition}{Определение}[chapter]
\newtheorem{example}{Пример}[chapter]
\newtheorem{principle}{Принцип}[chapter]

% Сокращения.
\newcommand{\Neigh}{\mathcal{N}}
\newcommand{\conf}{\mathrm{conf}}
\newcommand{\mood}{\mathrm{mood}}
\newcommand{\Reff}{R_{\mathrm{eff}}}

\title{Конструирование графовой памяти Рин\\[0.4em]
\large Монография о субстрате, динамике и эксперименте:\\
как построить, оживить и проверить память локального AI-компаньона}
\author{Рабочая монография проекта RIN}
\date{\today}

\begin{document}

\frontmatter
\maketitle

\vspace{1.5em}
\begin{center}
\begin{minipage}{0.82\textwidth}\itshape\small
Рабочий документ проекта RIN. Составлен как конспект серии проектных обсуждений и
приведён в форму монографии. Изложение намеренно избыточно: лучше сохранить лишнюю
деталь, псевдокод или разобранный пример, чем оставить механизм недосказанным.
Формулы служат точности, врезки «Интуиция», «Ловушка» и «Разбор примера» — смыслу,
предупреждению и проверке. Приложения содержат нотацию, глоссарий внешних идей и
сводку алгоритмов.
\end{minipage}
\end{center}
\vspace{1.5em}

\tableofcontents

\mainmatter

\chapter{Введение: память как единый субстрат}

\section{Зачем вообще граф}

У большинства современных диалоговых агентов память устроена плоско: история
сообщений, иногда векторная база, куда складываются «воспоминания» как
независимые записи, и retrieval по косинусной близости к текущему запросу. Это
работает для ассистента, который обслуживает разрозненные запросы, но плохо
описывает существо, у которого есть биография, характер, меняющееся настроение
и убеждения, способные оказаться ошибочными. У человека память не список; она
сеть. Одно воспоминание тянет за собой другое; факт связан с тем, откуда он
узнан; недавняя новость медленно перекрашивает старые выводы; забытое не
исчезает мгновенно, а тускнеет и иногда всплывает.

Поэтому центральная ставка проекта Рин — \emph{память как граф}. Но сам по себе
выбор структуры данных ничего не решает. Интересна не «база в виде графа», а то,
какие процессы по этому графу текут и каким законам они подчиняются. Эта
монография — про процессы.

\section{Главный тезис}

\begin{intuition}[Одна фраза, которую стоит держать в голове до конца]
Память Рин — это \textbf{один типизированный темпоральный граф с метаданными на
узлах и рёбрах}, а всё остальное — \emph{операторы над этим субстратом}:
retrieval, пропагация уверенности, забывание и консолидация, ночные обходы,
настроение, слой личности. Если довести до чистоты \emph{схему субстрата}, то
все прочие механизмы становятся модулями, читающими и пишущими одни и те же
поля. Если схему не довести — каждый модуль потащит свою полусовместимую модель
памяти, и они не срастутся.
\end{intuition}

Эта установка определяет и порядок изложения, и порядок разработки. Сначала мы
описываем субстрат (глава~\ref{ch:substrate}) и его временную многослойность
(глава~\ref{ch:timescales}). Затем — динамические процессы: настроение
(глава~\ref{ch:mood}), обход графа (глава~\ref{ch:walker}), работу с
уверенностью и ревизией убеждений (главы~\ref{ch:epistemics}
и~\ref{ch:propagation}). Отдельная глава~\ref{ch:topology} посвящена тому, как
\emph{топология} самого графа незаметно управляет всеми этими потоками. Затем мы
переходим к слою личности и голоса (главы~\ref{ch:personality}
и~\ref{ch:voice}), к вопросу о том, в каком смысле всё это является
\emph{обучением без обучения} (глава~\ref{ch:learning}), и, наконец, к программе
экспериментов на суррогатном графе (глава~\ref{ch:experiment}).

\section{Четыре кита и почему их нельзя тюнить порознь}

Полезно заранее назвать четыре процесса, вокруг которых всё вращается:

\begin{enumerate}
\item \textbf{Правильный retrieval} — как из графа достаётся релевантное, с
учётом запроса и текущего состояния.
\item \textbf{Правильная пропагация уверенности} — как новый факт меняет
уверенность соседних узлов, не устраивая ни ступора, ни лавины.
\item \textbf{Правильное забывание и консолидация} — что тускнеет, что
укрепляется, что сливается в обобщения.
\item \textbf{Правильные ночные обходы} — фоновая реорганизация и починка,
которую агент делает «во сне».
\end{enumerate}

Важнейшее наблюдение: эти четыре кита \emph{не независимы}. Они читают и пишут
одни и те же поля активации и уверенности на одном и том же графе. Забывание
меняет топологию; изменённая топология меняет retrieval; изменённый retrieval
меняет то, что пропагируется; пропагация меняет то, что найдёт ночной обходчик.
Поэтому их нельзя настраивать по отдельности — их надо со-настраивать, и именно
поэтому (глава~\ref{ch:experiment}) нужен управляемый стенд с заранее известной
истиной: без него взаимодействие операторов просто не измерить.

\section{Родство с математикой DTI: одна спина, два применения}

Эта монография выросла рядом с другим исследованием автора — магистерской
работой о предсказании аффинности связывания белок–лиганд средствами
графового и геометрического глубокого обучения. На первый взгляд это далёкие
темы. На деле у них общая математическая спина.

Консолидация памяти (переход от подробных эпизодов к сжатым смыслам) — это
\emph{огрубление графа} (\textit{coarse-to-local}): какой сигнал сохраняется при
переходе от мелкого масштаба к крупному, а какой теряется. Это в точности
центральный вопрос DTI-работы, только там он про interface-релевантный сигнал
взаимодействия, а здесь — про то, что выживает при консолидации воспоминаний.
Дальше по тексту мы увидим, что over-squashing (удушение дальней связи узким
местом графа), спектральные представления, персистентный лапласиан и
эффективное сопротивление — инструменты, которые работают в обе стороны.
Осознание этого родства — не поэтическая вольность: те же операторы, только
приложенные к разным графам. Мы будем возвращаться к этой связке, когда она
даёт конкретную пользу.

\section{Для кого этот текст}

У монографии два адресата, и оба важны.

Первый — \emph{автор через две недели}, вернувшийся после напряжённой работы над
другими задачами (магистерская, беттинг) и забывший половину деталей. Для него
текст должен восстанавливать замысел \emph{с нуля}, а не намекать на него.

Второй — \emph{коллеги по команде}, которые в проектировании этой памяти пока не
участвовали и «не в теме». Для них текст должен вводить каждое понятие явно, до
первого употребления, и объяснять не только «как», но и «зачем».

Отсюда жёсткое правило изложения: \textbf{ни один объект или механизм не
используется раньше, чем определён}. Глава~\ref{ch:architecture} даёт систему
целиком одним взглядом; глава~\ref{ch:definitions} — словарь всех основных
объектов; и лишь затем идут детальные разборы. Примеры сознательно берутся из
жизни Рин (стрим, зрители, бытовые ситуации), чтобы механизм можно было «пощупать»
на конкретной гипотетической сцене.

\begin{intuition}[Статус документа: рабочая гипотеза, а не доказанная теория]
Важно с самого начала: это \emph{инженерная рабочая модель} памяти для RIN, а не
финальная научная теория того, как обязана быть устроена память агента. Большинство
утверждений здесь — \emph{дизайн-решения} («мы решили строить так») и \emph{рабочие
гипотезы} («думаем, что сработает»), а не теоремы. Там, где есть известная
математика, мы на неё ссылаемся; там, где не знаем, — ставим врезку «Открытый
вопрос»; там, где аналогия помогает думать, но ничего не доказывает, — говорим об
этом прямо. Читать документ стоит как \emph{консистентную архитектурную гипотезу,
которую предстоит ломать на стенде} (глава~\ref{ch:experiment}), а не как
установленную истину.
\end{intuition}

\section{Как читать монографию}

Изложение намеренно избыточно. Там, где метод опирается на инженерную или
научную идею со стороны (система поддержки истинности, самоорганизованная
критичность, теория оценивания эмоций, ACT-R и прочее), мы кратко описываем и
саму идею, чтобы читатель не был вынужден держать в голове внешний контекст.
Формулы служат для точности, а врезки «Интуиция» и «Ловушка» — для смысла и для
предупреждения о типичных ошибках. Ссылки на литературу собраны в
главе~\ref{ch:references}; часть из них — работы из библиографий рабочих волтов
проекта (RIN и RIN–DTI), это отмечено.

\chapter{Архитектура целиком}
\label{ch:architecture}

Прежде чем разбирать механизмы по отдельности, полезно увидеть систему целиком:
из чего она состоит, что по ней течёт и как компоненты связаны. Эта глава — карта.
Всё, что здесь названо, ниже определяется формально (глава~\ref{ch:definitions}) и
разбирается подробно; сейчас задача — общий взгляд.

\section{Что мы вообще строим}

Рин — локальный AI-компаньон. У неё должна быть \emph{память}, которая ведёт себя не
как журнал сообщений, а как память живого существа: связная, ассоциативная,
меняющаяся во времени, способная ошибаться и исправляться, окрашенная настроением и
характером. Центральная ставка проекта: такую память естественно представить
\emph{графом}, а всё «поведение памяти» — процессами (операторами) над этим графом.

\begin{principle}[Единый субстрат]
Память Рин — это \textbf{один типизированный темпоральный граф с метаданными на
узлах и рёбрах}. Retrieval, вывод (fusion), консолидация, забывание, ночная работа
(сон), настроение и слой личности — не отдельные подсистемы со своими базами, а
\emph{операторы}, читающие и пишущие одни и те же поля одного графа.
\end{principle}

\begin{pitfall}[«Один субстрат» — дисциплина, а не «всё есть граф»]
Тезис легко выродить в догму, и тогда в граф начнут запихивать всё подряд:
runtime-очереди, streaming-флаги, транзиентные тики настроения, состояние
планировщика. Это превратит граф в свалку «всего, что случилось». Граница проходит
чётко: \textbf{в графе живёт долговременное содержание и его эпистемические
обоснования} — воспоминания, факты, убеждения, связи, провенанс. Остальное — рядом,
и это \emph{не} граф:
\[
\text{граф} = \text{память};\quad
\text{event log} = \text{история};\quad
\text{телеметрия} = \text{самонаблюдение};
\]
\[
\text{scheduler state} = \text{управление};\quad
\text{prompt context} = \text{временная сборка для LLM}.
\]
Они связаны и обмениваются данными, но это разные хранилища с разными законами
жизни (главы~\ref{ch:runtime},~\ref{ch:mood}). Настроение — сцепленный сбоку модуль;
самотелеметрия метакогниции — кольцевой буфер; планировщик — конечный автомат.
Правило простое: если данное \emph{транзиентно} или это \emph{управление/история} —
ему не место в графе памяти.
\end{pitfall}

\section{Слои системы}

Система распадается на четыре слоя. Снизу вверх:

\begin{enumerate}
\item \textbf{Субстрат} (глава~\ref{ch:substrate}) — сам граф: узлы-воспоминания и
типизированные знаковые рёбра, с полями \texttt{confidence}, \texttt{valence},
\texttt{provenance}, временной слой, \texttt{activation}, флаги \texttt{need\_to\_*}.
Его начальное наполнение — \emph{эмбриональный граф}
(глава~\ref{ch:embryonic}) из готовых баз здравого смысла.
\item \textbf{Операторы} — процессы над субстратом:
\emph{обход} роем случайных блуждающих агентов (глава~\ref{ch:walker}),
\emph{вывод и консолидация} (глава~\ref{ch:fusion}),
\emph{пропагация уверенности} (глава~\ref{ch:propagation}),
\emph{забывание} (глава~\ref{ch:timescales}),
\emph{ночная работа/сон} (глава~\ref{ch:sleep}).
\item \textbf{Состояния} — сцепленные с графом, но отдельные:
\emph{настроение} (глава~\ref{ch:mood}) и производные эпистемические статусы
(глава~\ref{ch:epistemics}).
\item \textbf{Выражение} — то, что превращает внутреннее в реплику:
\emph{личность} (глава~\ref{ch:personality}) и \emph{голос/языковая модель}
(глава~\ref{ch:voice}).
\end{enumerate}

\begin{center}
\begin{tikzpicture}[font=\small,>=Stealth,
  layer/.style={draw,rounded corners,minimum width=82mm,align=center,inner sep=6pt,fill=gray!5},
  lbl/.style={font=\bfseries}]
  \node[layer] (expr) {{\bfseries Выражение}\\[1pt]
     \footnotesize личность \quad$\cdot$\quad Rinness / языковая модель};
  \node[layer,below=7mm of expr] (state) {{\bfseries Состояния}\\[1pt]
     \footnotesize настроение (mood) \quad$\cdot$\quad уверенность / статусы};
  \node[layer,below=7mm of state] (op) {{\bfseries Операторы}\\[1pt]
     \footnotesize обход (walker) $\cdot$ fusion $\cdot$ консолидация\\
     \footnotesize пропагация $\cdot$ сон / скраб};
  \node[layer,below=7mm of op] (sub) {{\bfseries Субстрат}\\[1pt]
     \footnotesize граф: узлы + знаковые рёбра + метаданные\\
     \footnotesize эмбриональный граф};
  \draw[<->,thick] (expr)--(state);
  \draw[<->,thick] (state)--(op);
  \draw[<->,thick] (op)--(sub);
\end{tikzpicture}
\captionof{figure}{Четыре слоя системы. Стрелки двусторонние: настроение смещает
обход, обход меняет граф, граф питает настроение, и так далее.}
\label{fig:layers}
\end{center}

\section{Что течёт по системе}

Полезно различать четыре «потока», все — над одним графом:

\begin{itemize}
\item \textbf{Активация} — транзиентная «подсветка» узлов при обходе; быстро гаснет.
\item \textbf{Уверенность} — медленно меняющееся поле $\conf$; её пропагация и
ревизия — эпистемическое ядро.
\item \textbf{Валентность и настроение} — аффективная окраска узлов и сцепленный
сбоку вектор настроения.
\item \textbf{Структурные изменения} — рождение узлов (fusion), укрепление/распад
рёбер (консолидация/забывание): граф эволюционирует.
\end{itemize}

\section{Жизнь одного воспоминания}

Сквозной проход, чтобы увидеть, как слои работают вместе. Пусть на стриме зритель
пишет: «Рин, я опять слил ранг в доте».

\begin{workedbox}[От реплики до реплики]
\textbf{1. Восприятие.} Событие кодируется как узел \texttt{ev}: «зритель $Z$
проиграл». Оценивается его значимость (глава~\ref{ch:personality}); заметный сдвиг
настроения делает событие «важным».

\textbf{2. Обход (retrieval).} Рой блуждающих агентов стартует от \texttt{ev},
смещённый текущим настроением Рин (озорным) и запросом. Всплывает окрестность:
«проигрыш $\to$ досада/смех», «$Z$ раньше хвастался скиллом», «дота».

\textbf{3. Вывод (fusion).} Совместная активация «$Z$ хвастался» и «$Z$ проиграл»
рождает \emph{гипотезу}-узел «повод беззлобно подколоть $Z$» — провизорно, с низкой
уверенностью и провенансом $\{$хвастался, проиграл$\}$.

\textbf{4. Выражение.} Личность (участие, колкая нежность) и настроение выбирают
\emph{тон}; Rinness-модель \emph{композирует} реплику из всплывших ингредиентов —
не берёт готовую фразу из графа. На выходе: «Оёй, ну куда без слитого ранга. Иди
чаёк попей, чемпион».

\textbf{5. Запись.} Реплика и реакция $Z$ пишутся провизорно. Если $Z$ смеётся —
уверенность ребра «подкол работает с $Z$» растёт; если обиделся — гаснет, и
обновляется модель отношения с $Z$.

\textbf{6. Ночью.} Во сне (глава~\ref{ch:sleep}) телепортирующий обходчик случайно
переинспектирует память, консолидирует частые эпизоды («$Z$ — частый гость, любит
доту»), проверяет противоречия, иногда рождает неожиданные ассоциации.
\end{workedbox}

Каждый шаг — отдельная глава. Дальше мы разбираем их по одному, но держите в голове
именно этот сквозной цикл: \emph{восприятие $\to$ обход $\to$ вывод $\to$ выражение
$\to$ запись $\to$ ночная переработка}.

\section{Карта чтения}

Читателю-новичку достаточно глав~\ref{ch:architecture} (эта),
\ref{ch:definitions} (словарь) и сквозного примера выше, чтобы понимать разговоры
команды. Дальше главы можно читать по интересу: субстрат и эмбриональный граф — для
тех, кто будет строить хранилище; обход, fusion, сон — для тех, кто пишет
операторы; настроение, личность, голос — для тех, кто отвечает за характер;
эксперименты (глава~\ref{ch:experiment}) — для всех, потому что проверять придётся
всё.

\chapter{Словарь основных объектов}
\label{ch:definitions}

Эта глава вводит все базовые понятия \emph{до} того, как они используются в
механизмах. Определения намеренно сжаты — детали в соответствующих главах; здесь
задача, чтобы читатель (в том числе новый в проекте) владел словарём.

\section{Субстрат}

\begin{definition}[Граф памяти]
\emph{Граф памяти} $G=(V,E)$ — направленный граф, где $V$ — узлы-воспоминания, а
$E$ — типизированные знаковые взвешенные рёбра-отношения. Граф \emph{темпоральный}
(узлы несут время/слой) и \emph{динамический} (операторы меняют и узлы, и рёбра).
\end{definition}

\begin{definition}[Узел памяти]
\emph{Узел} $v\in V$ — запись с содержанием $\mathbf{x}_v$ и полями состояния:
вид \texttt{node\_kind} (observation/claim/hypothesis/preference/reaction\_pattern/
relationship\_state/\dots; задаёт смысл уверенности, гл.~\ref{ch:substrate}),
уверенность $\conf(v)\in[-1,1]$, валентность (аффективная окраска), временной слой,
провенанс, важность \texttt{importance}, флаги
\texttt{need\_to\_discover}/\texttt{need\_to\_revise}, транзиентная активация $a_v$.
Узел одновременно — данные (что помнится), убеждение (насколько истинно) и след
(насколько активно и свежо).
\end{definition}

\begin{definition}[Ребро и два его канала]
\emph{Ребро} $(u\to v)$ направлено и несёт \emph{тип} отношения (\textsc{isA},
\textsc{supports}, \textsc{contradicts}, \textsc{associates}, \dots). На нём живут
\textbf{два независимых числовых канала}, которые читают разные операторы:
\begin{itemize}
\item \textbf{вес перехода} $w_{uv}\ge 0$ — для обхода (walker); нормирован по узлу,
$\sum_{v} P(u\to v)=1$ (это и есть PageRank-подобная часть);
\item \textbf{знак} $s_{uv}\in\{+1,-1\}$ — полярность отношения (подтверждает /
опровергает); используется \emph{только} в пропагации уверенности
(глава~\ref{ch:propagation}), чтобы уверенность могла не только расти, но и убывать.
\end{itemize}
Знак и вес не конфликтуют: первый — про \emph{убеждение}, второй — про \emph{обход}.
\end{definition}

\begin{definition}[Уверенность]
$\conf(v)\in[-1,1]$ — \emph{знаковая} степень убеждённости: $+1$ — уверена, что
\emph{да}; $-1$ — уверена, что \emph{нет}; $0$ — не знает. Магнитуда $|\conf(v)|$ —
собственно уверенность, знак — направление убеждения. Эквивалентно вероятности
через $\conf=2P-1$ (где $P$ — вероятность истинности, $P=0{,}5\Leftrightarrow
\conf=0$), но эргономичнее: зона незнания — это $|\conf(v)|<h$ (порог $h\approx
0{,}3$), и такие узлы — кандидаты на \texttt{need\_to\_discover} или удаление.
Меняется пропагацией (глава~\ref{ch:propagation}) и ревизией
(глава~\ref{ch:epistemics}). Принцип: вывод не получает $|\conf|=1$.
\end{definition}

\begin{definition}[Активация]
\emph{Активация} $a_v(t)\ge 0$ — транзиентное скалярное поле «насколько узел
подсвечен \emph{сейчас}». Обновляется распространением активации
\[
a(t{+}1)=\lambda\,\tilde P^{\top} a(t)+s(t),\qquad \lambda\in(0,1),
\]
где $\tilde P$ — нормированная матрица переходов, $s(t)$ — вброс (визиты роя,
инъекция запроса), $\lambda$ — распад. Эквивалентно — ожидаемая частота посещений
узла роем в коротком окне. \textbf{Не путать:} активация — «здесь и сейчас»,
уверенность — медленное убеждение, стационарное $\pi$ — долгосрочная частота
(глава~\ref{ch:topology}).
\end{definition}

\begin{definition}[Провенанс (обоснование)]
\emph{Провенанс} узла — \emph{множество идентификаторов узлов-посылок}
$\{v_1,\dots,v_n\}$, из которых узел непосредственно выведен (для наблюдаемого факта
— идентификатор внешнего источника). Это \emph{justification set} в смысле
поддержки истинности (глава~\ref{ch:epistemics}), а \textbf{не} полный трек обхода:
трек шумен, доступен отдельно и обоснованием не является. Провенанс делает возможной
ретракцию: опровергнув посылку, находим и понижаем её следствия.
\end{definition}

\begin{definition}[Валентность и настроение]
\emph{Валентность} узла — его \emph{аффективная} окраска (пара
валентность–возбуждение), \textbf{ортогональная уверенности}: уверенность — насколько
узел \emph{истинен}, валентность — как он \emph{ощущается}. Пример: «точно знаю, что
это ужасное воспоминание» — $\conf$ высокая, валентность резко негативная.
Валентность — канал сцепки с настроением (mood-congruent recall,
глава~\ref{ch:mood}); уверенность — канал сцепки с пропагацией. \emph{Настроение}
$\mood$ — отдельный низкоразмерный вектор состояния Рин, сцепленный с графом
двусторонне, но не хранящийся в узлах.
\end{definition}

\begin{definition}[Время: системное и содержательное]
На узле хранится \emph{минимум системного времени}: тег \texttt{layer} (сессия
\dots жизнь) и одна метка (кодирование / последний доступ) — этого достаточно
операторам забывания и консолидации. \emph{Содержательное} время («когда это
произошло») моделируется как обычные узлы-факты — оно реконструктивно и может быть
неверным, как у человека. Интервал валидности (start/end) вешается \emph{только} на
факты, которым он нужен (квалификаторы Wikidata-схемы, глава~\ref{ch:substrate}), а
не на каждый узел.
\end{definition}

\section{Обход}

\begin{definition}[Случайный блуждающий агент, Random Walker]
\emph{Random Walker} (RW) — агент, который стартует в некотором узле и шаг за шагом
переходит к соседям по \emph{вероятностному} правилу
$P(A\to B)\propto w_{AB}\cdot\text{rel}(B)\cdot\text{mc}(B)\cdot\text{unc}(B)^{\beta_u}$
(глава~\ref{ch:walker}), оставляя за собой \emph{след} — последовательность
посещённых узлов. По ходу он может запускать операции (извлечение, вывод,
инспекцию).
\end{definition}

\noindent Функционал одного RW: (i)~перейти по вероятностному правилу; (ii)~с
вероятностью $\alpha$ \emph{телепортироваться} (рестарт в случайный узел);
(iii)~повысить активацию посещённого узла; (iv)~по условию запустить операцию
(retrieval / fusion / инспекцию need\_to\_*); (v)~записать след.

\begin{definition}[Рой блуждающих агентов, RW Swarm]
\emph{Swarm} — множество RW, запущенных параллельно с общей политикой. Рой
аппроксимирует распространение активации (spreading activation) и — на большой
статистике следов — эмпирическое распределение маршрутов по графу
(глава~\ref{ch:experiment}).
\end{definition}

\begin{definition}[Три политики обхода]
Один и тот же движок роя работает в трёх режимах, задаваемых параметрами
(смещение перехода, телепорт $\alpha$, температура): \emph{retrieval} (под
запрос), \emph{внутренняя жизнь/любопытство} (к неопределённости) и \emph{сон}
(стохастическое покрытие с телепортом). См.\ главу~\ref{ch:walker}.
\end{definition}

\section{Операции над памятью}

\begin{definition}[Retrieval (извлечение)]
\emph{Retrieval} — двухшаговая операция. \textbf{Шаг~1 (якоря):} быстрым индексом
(векторный/эмбеддинг-поиск, как в A-MEM/Zep) находим узлы, релевантные запросу; RW
для \emph{поиска} входа не годится (медленен, диффузен). \textbf{Шаг~2
(расширение):} от якорей запускаем ограниченный mood-смещённый рой RW, собирающий
активированную \emph{окрестность} — ассоциации, которых векторный поиск не даёт.
Итог зависит от настроения (mood-congruent recall). RW, таким образом, нужен и днём
(расширение), а сон — тот же RW плюс телепорт и без границ.
\end{definition}

\begin{definition}[Fusion (слияние, вывод)]
\emph{Fusion} — порождение \emph{нового узла} из совместно активированных: система
замечает согласованную конфигурацию посылок и создаёт узел-\emph{гипотезу} с низкой
уверенностью и провенансом-обоснованием. Fusion — механизм рождения выводов и
ассоциаций (глава~\ref{ch:fusion}).
\end{definition}

\begin{definition}[Консолидация]
\emph{Консолидация} (в широком смысле) — \emph{процесс} перехода от множества частных
следов к устойчивым обобщениям, а не единичный оператор. Он отрабатывает несколько
подоператоров: fusion\,/\,вывод, summary-overlay, dedup\,/\,merge и — как \emph{поздний}
подоператор — prune (destructive replacement с огрублением, graph coarsening).
Большинство ранних консолидаций ничего не разрушают; разбор процесса и его
предохранителей — в главе~\ref{ch:timescales}.
\end{definition}

\begin{definition}[Забывание (decay)]
\emph{Забывание} — ослабление неиспользуемого; отрицательная сторона консолидации.
Моделируется базовой активацией (частота $+$ свежесть, глава~\ref{ch:timescales}).
\end{definition}

\begin{definition}[Ночная работа (сон)]
\emph{Сон} — фоновый режим, когда телепортирующий рой случайно инспектирует память:
консолидирует, проверяет противоречия (скраб), рождает неожиданные fusion-ы,
подтягивает внешнее по флагам \texttt{need\_to\_discover} (глава~\ref{ch:sleep}).
\end{definition}

\begin{definition}[Значимость (salience) и оценивание]
\emph{Значимость} события — мера того, насколько оно важно \emph{для самой Рин}
(теория оценивания). Дешёвый прокси — величина изменения настроения $|\Delta\mood|$;
глубокий сигнал — ошибка предсказания. Значимость гейтит консолидацию
(глава~\ref{ch:personality}).
\end{definition}

\section{Выражение}

\begin{definition}[Личность]
\emph{Личность} — не таблица весов в графе, а распределённый процесс: мотивация
$+$ настроение $+$ значимость всплывшего, входящий как \emph{смещение во время
обхода}, а на выходе — как стиль (глава~\ref{ch:personality}).
\end{definition}

\begin{definition}[Rinness / голос]
\emph{Голос} — языковая модель, превращающая внутренний контекст (всплывшие
ингредиенты, настроение, выбранный тон) в реплику в узнаваемом стиле Рин. Часть
характера живёт здесь; факты — в графе (глава~\ref{ch:voice}).
\end{definition}

\begin{definition}[Эмбриональный граф]
\emph{Эмбриональный граф} — начальное наполнение памяти из готовых баз здравого
смысла (ConceptNet $+$ ATOMIC), нейтральное по характеру; на нём Рин «взращивает»
собственную память фьюжнами и консолидацией (глава~\ref{ch:embryonic}).
\end{definition}

\chapter{Субстрат: типизированный темпоральный граф}
\label{ch:substrate}

\section{Что такое узел и что такое ребро}

Память Рин — направленный граф $G=(V,E)$, в котором узлы несут содержание
(воспоминания, факты, сущности, концепты), а рёбра — типизированные и знаковые
отношения между ними. Ключевая мысль главы: и узлы, и рёбра несут
\emph{метаданные}, и именно метаданные, а не сам факт связи, делают возможными
все дальнейшие механизмы.

\begin{definition}[Узел памяти]
Узел $v\in V$ — это запись с содержанием $\mathbf{x}_v$ и набором полей состояния,
над которыми работают операторы памяти. Узел одновременно является и «данными»
(что помнится), и «убеждением» (насколько это считается истинным), и «следом»
(насколько активно и свежо).
\end{definition}

\subsection{Схема узла}

\begin{longtable}{p{0.24\textwidth} p{0.10\textwidth} p{0.56\textwidth}}
\toprule
\textbf{Поле} & \textbf{Тип} & \textbf{Смысл и кто с ним работает} \\
\midrule
\endhead
\texttt{id} & int & идентификатор \\
\texttt{content} & text & содержание в естественном виде \\
\texttt{embedding} $\mathbf{x}_v$ & vec$[d]$ & векторное представление; retrieval-релевантность, мост баз \\
\texttt{node\_kind} & enum & вид узла (observation/claim/hypothesis/preference/\dots); задаёт смысл \texttt{confidence} (см.\ раздел~\ref{sec:nodekind}) \\
\texttt{confidence} $c_v$ & $[-1,1]$ & для truth-apt узлов — знаковое убеждение ($+$да/$-$нет/$0$хз); для остальных — надёжность оценки (раздел~\ref{sec:nodekind}) \\
\texttt{valence} & vec$[2]$ & валентность–возбуждение; сцепка с настроением \\
\texttt{layer}/\texttt{ts} & enum/time & масштаб (сессия\ldots жизнь) и/или метка времени; консолидация \\
\texttt{provenance} & list & источники и/или обоснования; TMS, ревизия \\
\texttt{need\_to\_*} & flags & «исследовать/пересмотреть»; любопытство, скраб \\
\texttt{importance} & $[0,1]$ & значимость: ставится salience при кодировании (гл.~\ref{ch:personality}), подкрепляется центральностью; гейтит прунинг (гл.~\ref{ch:timescales}) \\
\texttt{domains} & tags & членство в доменах знаний; несёт importance \emph{по членству}, а не диффузией (гл.~\ref{ch:timescales}) \\
\texttt{activation} $a_v$ & float & транзиентная активация; обход \\
\texttt{degree}/\texttt{centrality} & float & производные структурные; топология \\
\bottomrule
\end{longtable}

\subsection{Схема ребра}

Ребро $(u\to v)$ несёт направленность (\texttt{src},\texttt{dst}), тип отношения
\texttt{type} (\textsc{isA}, \textsc{partOf}, \textsc{supports},
\textsc{contradicts}, \textsc{associates}, \dots) и \textbf{два независимых числовых
канала} (см.\ определение в главе~\ref{ch:definitions}): \emph{вес перехода}
$w_{uv}\ge 0$ (нормируется по узлу, для обхода) и \emph{знак}
$s_{uv}\in\{+1,-1\}$ (полярность отношения, для пропагации уверенности). Знак —
\emph{хранимый лейбл} (внутреннее свойство связи: «эта связь за или против?»), а не
вычисляемая величина: его нельзя «посчитать по пути», он известен при создании ребра.
Изменение уверенности $\Delta\conf$ — уже \emph{производное}, оно вычисляется в
пропагации \emph{с использованием} знака (глава~\ref{ch:propagation}).

\begin{intuition}[Веса — про свидетельство, а не про характер]
Забегая вперёд (глава~\ref{ch:personality}): вес ребра выражает
\emph{эвиденциальную} силу связи и уверенность, а \emph{не} «характер». Соблазн
задать личность статическими множителями на рёбра-реакции («грусть$\to$сарказм
$\times 2{,}0$») мы разберём и отвергнем: характер должен входить как
\emph{смещение во время обхода} от текущего состояния, а не как вшитый в граф
множитель.
\end{intuition}

\section{Почему рёбра направленные, типизированные и знаковые}

Плоская модель «узлы и связи одинакового сорта» ломается сразу на трёх вещах:
направленность (асимметрия причинности и времени), тип (нельзя валить covalent-связь,
близость, «подтверждает» и «является частью» в один линейный канал) и знак (без
него не отличить «всё усилилось» от «всё перевернулось»).

Правильный оператор поэтому не общий лапласиан сглаживания, а
\emph{relation-specific} передача сообщения. Сообщение по ребру:
\begin{equation}
m_{u\to v} \;=\; \phi_{r_{uv}}\!\big(h_u,\, h_v,\, e_{uv},\, \Delta t,\, c_{uv},\, s_{uv}\big),
\label{eq:message}
\end{equation}
где $r_{uv}$ — тип отношения, $\Delta t$ — возраст/временная дистанция, $c_{uv}$ —
уверенность, $s_{uv}$ — источник/эпистемический статус. Затем узел агрегирует
входящие сообщения — и здесь важно не одну статистику, а набор (сумма/среднее/макс/
дисперсия по типам отношений), потому что для памяти существенны и кратность
подтверждений, и пропорции, и наличие экстремального свидетельства.

\begin{workedbox}[Почему одна статистика теряет смысл]
Ключевое свойство \emph{среднего} — инвариантность к дублированию. Пусть про факт
есть свидетельства $\{+,-\}$ (один за, один против) и в другом случае
$\{+,+,-,-\}$ (двое за, двое против). Считая $+=1,-=-1$, получаем
$\mathrm{mean}\{+,-\}=0=\mathrm{mean}\{+,+,-,-\}$ — среднее \emph{одинаково}, хотя во
втором случае свидетелей вдвое больше. Для памяти это принципиально: «четверо
разделились» — не то же, что «один на один» (устойчивее, эпистемически весомее).
\emph{Сумма} же сохраняет кратность (в паре с числом слагаемых), а \emph{макс} по
знаку — наличие сильного контрсвидетельства. Отсюда правило: агрегатор хранит
\emph{набор} статистик (сумма, среднее, макс, дисперсия), а не одну, — иначе теряется
либо кратность, либо пропорция, либо экстремум.
\end{workedbox}

\begin{intuition}[Аналогия с электрической цепью]
Полезно думать о графе как о цепи с проводниками разного сорта. Направленность —
это диоды; тип ребра — материал проводника со своей проводимостью; знак — это
полярность. Плоская модель — цепь, где все элементы заменены одинаковыми
резисторами: она проведёт ток, но потеряет всю логику схемы.
\end{intuition}

\begin{center}
\begin{tikzpicture}[>=Stealth,font=\footnotesize,
  mem/.style={draw, rounded corners, minimum height=7mm, inner sep=3pt},
  inf/.style={draw, rounded corners, fill=blue!5, minimum height=7mm, inner sep=3pt}]
  \node[mem] (a) {разбил кружку};
  \node[mem, below=9mm of a] (d) {юзер хотел новую};
  \node[inf, right=26mm of a] (r) {реакция: грусть};
  \node[inf, right=22mm of d] (x) {вывод: уместно облегчение};
  \draw[->,thick] (a) -- node[above,font=\scriptsize]{$+$ xReact} (r);
  \draw[->,dashed] (a) -- (x);
  \draw[->,dashed] (d) -- (x);
  \draw[->,thick,red!60!black] (x) -- node[right,font=\scriptsize]{$-$ ослабляет} (r);
\end{tikzpicture}
\captionof{figure}{Противоречие не самостоятельное, а \emph{контекстное}. Базовое
$+$xReact «разбил $\to$ грусть» ослабляется не узлом «хотел новую» самим по себе, а
\emph{совместным} выводом «разбил $+$ хотел новую $\to$ уместно облегчение»
(пунктир — провенанс; это n-арная конфигурация, глава~\ref{ch:fusion}). Именно она
делает возможным уместное «поздравляю».}
\label{fig:cup}
\end{center}

\section{Виды узлов и что значит уверенность}
\label{sec:nodekind}

Раньше мы говорили «узел — это одновременно данные и убеждение». Красиво, но для
реализации грубо: \textbf{не всякий узел — утверждение об истинности}. «Португалия
имеет такой-то режим магазинов», «этот вывод поддержан источником», «гипотеза
House\,/\,Cast Away ложна» — да, это truth-apt утверждения, и для них
$\conf\in[-1,1]$ означает истинность. Но «Ярославу заходит такой подкол», «Рин не
любит такой тон», «этот человек устал», «у Анны есть кот» — это предпочтения,
реакции, склонности, состояния, и «уверенность 0,8, что Ярослав любит X» — \emph{не}
то же самое, что «вероятность истинности факта 0,9». Предпочтения меняются, зависят
от контекста, имеют исключения.

Поэтому у узла есть поле \texttt{node\_kind}, и оно \emph{задаёт смысл}
\texttt{confidence}:

\begin{longtable}{p{0.26\textwidth} p{0.30\textwidth} p{0.34\textwidth}}
\toprule
\textbf{node\_kind} & \textbf{что это} & \textbf{смысл confidence} \\
\midrule
\endhead
\texttt{observation} & прямое наблюдение & истинность (высокая) \\
\texttt{claim} & утверждение о мире & истинность \\
\texttt{hypothesis} & вывод/гипотеза & истинность (низкая, провизорно) \\
\texttt{episode} & конкретное событие & было/не было \\
\texttt{preference} & вкус/склонность & надёжность оценки (не истинность!) \\
\texttt{reaction\_pattern} & «на X реагирую Y» & надёжность (Beta-постериор, гл.~\ref{ch:personality}) \\
\texttt{relationship\_state} & состояние отношения / \emph{кто-что-знает} & надёжность модели \\
\texttt{boundary} & граница/нежелание & устойчивость \\
\texttt{skill/procedure} & навык/как-делать & освоенность \\
\texttt{self\_model} & убеждение о себе & надёжность самооценки \\
\texttt{mood\_summary} & свёртка настроения & — (это метка, не убеждение) \\
\bottomrule
\end{longtable}

Для не-truth-apt видов \texttt{confidence} читается как \emph{надёжность/устойчивость
оценки}, а не «вероятность истинности». Предпочтения и реакции лучше хранить как
\emph{модель с условиями} (при таком-то контексте — так), а не как голый скаляр
истинности.

\begin{principle}[Confidence приора $=$ source\_trust $\times$ kind, а не плоская константа]
Импортированным узлам (эмбрион: ConceptNet, ATOMIC) нельзя давать \emph{плоскую} уверенность
на всё — это дважды неверно. Baseline должен зависеть от двух осей:
\[
\text{confidence(node)} \;=\; \text{source\_trust}(\text{source}) \times \text{baseline}(\texttt{node\_kind}).
\]
\emph{По kind}: ConceptNet-концепт — truth-apt \texttt{claim}, стартует выше;
ATOMIC-эпизод — социально-каузальная \emph{оценка реакции}, стартует ниже, потому что это
надёжность мягкого прогноза, а не истинность (человеческие реакции сложнее фактов).
\emph{По источнику}: разные базы заслуживают разного доверия. Два принципиальных решения:
\begin{itemize}\itemsep2pt
\item \textbf{Источник — это провенанс, а не метка.} \texttt{source} уже несётся в
провенансе (\texttt{/c/en/\dots} vs \texttt{atomic}); не надо плодить эпистемические метки
\texttt{prior\_conceptnet}\,/\,\texttt{prior\_atomic} — метка (\texttt{prior}) это
\emph{эпистемический класс} для packer'а, а источник это ортогональные метаданные для
калибровки. Иначе метки взрываются комбинаторно, а LLM источник в метке и не нужен.
\item \textbf{Не выводить node-confidence из веса ребра.} Вес ConceptNet — корроборация
\emph{отношения} (канал ребра), а не истинность узла; каналы мы развели специально
(определение ребра). Вес гонит диффузию, confidence узла — про доверие к содержанию.
\end{itemize}
\end{principle}

\begin{openq}[source\_trust обучаем, но откалибровать нечем до наблюдений]
Сейчас \texttt{source\_trust} — грубый прайор (ConceptNet $>$ ATOMIC, обоснованно по kind),
не калиброванное число. \emph{Обучаемым} он становится, когда появится поток наблюдений
(write-path, гл.~\ref{ch:runtime}): эпистемический слой понижает доверие к источнику, чьи
приоры \emph{опровергаются} наблюдением, и повышает при подтверждении — так «монитор ловит»,
что источник систематически врёт, а Рин «понимает» ненадёжность источника, читая
\texttt{source\_trust}, а не по отдельной метке. Калибровать в вакууме (без downstream-сигнала)
нельзя — поэтому \emph{сейчас} фиксируем лишь \emph{структуру} (baseline по source$\times$kind,
source в провенансе), а обучение доверия откладываем до наблюдений.
\end{openq}

\begin{intuition}[Почему знаковый скаляр, а не тройка мнения]
Напрашивается более богатое представление — \emph{тройка субъективной логики}
(belief\,/\,disbelief\,/\,uncertainty), где неопределённость несёт собственную
координату. Мы сознательно берём один знаковый скаляр $c\in[-1,1]$ (эквивалент
$2P{-}1$: $+1$ «да», $-1$ «нет», $0$ «не знаю»). Причины две. \emph{Первое}:
скаляр уже разводит две вещи, которые чаще всего путают, — \emph{неверие} ($c<0$) и
\emph{незнание} ($c\approx0$); именно ради этого мы и ушли от $[0,1]$. \emph{Второе}:
третью ось — «сколько за этим свидетельств / насколько оценка устойчива» — мы и так
несём \emph{отдельно и разнотипно}: через \texttt{node\_kind}, через
\texttt{need\_to\_discover}, через Beta-неопределённость реакций
(гл.~\ref{ch:personality}) и провенанс. То есть неопределённость у нас не потеряна —
она просто не запихнута в один кортеж, а живёт в подходящих для каждого вида полях.
Тройку держим в уме как альтернативу на случай, если разнесённые сигналы окажется
неудобно сводить вместе.
\end{intuition}

\begin{intuition}[Кто-что-знает — тоже вид убеждения]
Особо отметим \texttt{relationship\_state} в смысле \emph{моделирования чужих
знаний}. «Рин знает $X$», «Ярослав знает, что Анна сказала $X$», «Анна не в курсе,
что Рин это запомнила» — это разные узлы: убеждения Рин \emph{о знаниях других}. Это
не про приватность-как-политику (её мы отдельной несущей системой не строим), а про
то, что theory-of-mind — законный и важный вид памяти: без него Рин не отличит «я это
знаю» от «он думает, что я этого не знаю».
\end{intuition}

\subsection{От словаря к поведению: операционная матрица}
\label{sec:kindmatrix}

Список видов — это словарь; реализация спросит \emph{поведение}. Каждый
\texttt{node\_kind} обязан ответить на семь вопросов: (1)~какие операторы его
\emph{читают}; (2)~какие могут в него \emph{писать}; (3)~как считается
\texttt{confidence}; (4)~как работает \emph{распад}; (5)~какие \emph{типы рёбер}
допустимы; (6)~можно ли его \emph{fusion-ить}; (7)~можно ли показывать LLM в
context-packer (гл.~\ref{ch:runtime}).

Полную матрицу $11\times 7$ заполнять сейчас — значит выдумывать точность, которой
нет: часть клеток прояснится только на стенде. Поэтому ниже — \emph{каркас} с
наиболее решающими столбцами и уверенными клетками; прочее помечено «TBD».

\begin{center}
\begin{tabular}{p{0.20\textwidth} p{0.22\textwidth} p{0.14\textwidth} p{0.12\textwidth} p{0.14\textwidth}}
\toprule
\textbf{kind} & \textbf{смысл confidence} & \textbf{распад} & \textbf{fusion?} & \textbf{в packer?} \\
\midrule
observation & истинность & медленный & да (посылка) & да \\
claim & истинность & медленный & да & да \\
hypothesis & истинность (низкая) & быстрый, если не подтв. & да & да, с меткой «гипотеза» \\
preference & надёжность оценки & по использованию & осторожно & да \\
reaction\_pattern & Beta-надёжность & по успеху/провалу & нет (это модель) & как контекст \\
relationship\_state & надёжность модели & медленный & осторожно & по ситуации (TBD) \\
boundary & устойчивость & очень медленный & нет & да \\
self\_model & надёжность самооценки & медленный & да (метакогниция) & да \\
mood\_summary & — (метка) & транзиентно & нет & редко \\
\bottomrule
\end{tabular}
\end{center}

Столбцы «кто читает / кто пишет / допустимые рёбра» определяются там, где живёт
соответствующий оператор, а не дублируются здесь. Матрица — это \emph{контракт}
поведения на тип, и она сознательно неполна: заполняется по мере постройки.

\section{Провенанс и обоснование}
\label{sec:prov}

Поле \texttt{provenance} шире, чем «источник». Для \emph{исходных} фактов провенанс
— откуда факт узнан. Для узлов, \emph{порождённых выводом} (fusion), провенанс
должен хранить \emph{обоснование} — множество родительских узлов, из которых вывод
получен. Эта деталь окажется критической в главе~\ref{ch:epistemics}: именно
наличие обоснований позволяет, опровергнув посылку, автоматически пометить на
пересмотр её следствия. Без записанных обоснований опровержение посылки не находит
потомков — и рождается «ползучая вера».

\section{Что значит «коммит» и кто его гейтит}

Важнейшая граница схемы — между \emph{активным рабочим убеждением} и \emph{записью}.
Не всё, что всплыло или сгенерировалось, попадает в граф как факт.

\begin{principle}[Гипотеза, а не факт]
Всё, что порождено выводом или генерацией (fusion, реакция, ассоциация), пишется
\emph{провизорно}: с низкой уверенностью, с провенансом-обоснованием и, при
необходимости, с флагом \texttt{need\_to\_discover}. Повышение до «записи» (переход
в более медленный слой) происходит только через накопление подтверждающих
свидетельств. Управление — на границе коммита.
\end{principle}

\begin{lstlisting}[style=pseudo]
def commit(G, item, kind):
    if kind == "observed":          # прямое наблюдение
        c0 = source_trust(item.source)
    if kind == "inferred":          # порождено fusion/выводом
        c0 = LOW                    # гипотеза, не факт
        item.provenance = item.parents      # обоснование для TMS
    node = G.add_node(item, confidence=c0, layer="session")
    propagate(G, node)              # см. главу про лавину (со знаком, карантином)
    return node
\end{lstlisting}

\section{Кража схемы у Wikidata}

Есть готовая, обкатанная модель данных, уже решившая ровно те задачи, что встают
перед метаданными Рин, — модель утверждений \emph{Wikidata}. Брать её
\emph{содержание} (энциклопедические факты) для существа, мимикрирующего под
человека, не нужно: человек помнит мало фактов, его сила в ассоциативности. Но её
\emph{схему} стоит перенять целиком:

\begin{itemize}
\item \textbf{квалификаторы} (точка во времени, начало/конец) — темпоральная
валидность: факт истинен \emph{в интервале}, а не вечно;
\item \textbf{references} — провенанс/обоснование;
\item \textbf{ранг} (\emph{preferred}/\emph{normal}/\emph{deprecated}) —
\emph{недеструктивная ревизия}: устаревшее не удаляют, а депрецируют, помечая новое
предпочтительным; история сохраняется для машинерии поддержки истинности.
\end{itemize}
Иными словами, факт стоит моделировать как
\[
\text{(утверждение)} + \text{(интервал валидности)} + \text{(провенанс)} +
\text{(статус/ранг)},
\]
и это буквально схема Wikidata. Не изобретать метаданные заново, а взять
проверенный шаблон.

\section{Врождённое и приобретённое: не граница, а континуум}

Соблазнительно разделить узлы на «врождённые константы» (эмбриональный граф) и
«приобретённые» (нажитые, с временными метками). Как первое приближение полезно, но
как жёсткая граница неверно. Консолидация \emph{двигает} записи через границу
(эпизод $\to$ обобщённое знание), а ревизия должна уметь \emph{понижать} и
«врождённое» — иначе унаследованную ошибку здравого смысла Рин не исправит никогда.
Правильная картина — \emph{континуум глубины консолидации}, где порог опрокидывания
убеждения растёт с глубиной, но нигде не становится бесконечным. К оси
«пластичность\,--\,упрямство» мы вернёмся не раз (главы~\ref{ch:propagation},
\ref{ch:epistemics}).

\chapter{Эмбриональный граф: с чего начинается память}
\label{ch:embryonic}

Пустой граф памяти бесполезен: у существа не бывает мышления «с нуля» — есть
врождённый здравый смысл, на который ложится личный опыт. Поэтому память Рин
начинается с \emph{эмбрионального графа} — нейтрального по характеру каркаса из
готовых баз здравого смысла, который затем обрастает личным опытом через fusion и
консолидацию. Эта глава отвечает предметно: из чего он собран, как выглядят его
узлы, как их сшить и как подключить эмбеддинги.

\section{Две базы, две разные роли}

Нужны две дополняющие базы (Wikidata как содержание не берём — человек не помнит
все столицы; глава~\ref{ch:substrate}):

\begin{itemize}
\item \textbf{ConceptNet} — \emph{предметный и таксономический} здравый смысл: что
чем является, из чего состоит, для чего используется. Даёт Рин «физику мира».
\item \textbf{ATOMIC} — \emph{причинно-социальный} здравый смысл о событиях и
ментальных состояниях: чего человек хотел, что почувствовал, что за этим следует.
Даёт Рин «психологию».
\end{itemize}
Порознь одного мало: ConceptNet без ATOMIC — мир без психологии; ATOMIC без
ConceptNet — реакции без предметного мира.

\section{Как выглядят узлы и рёбра ConceptNet}

ConceptNet — это уже граф. Узел — концепт (слово или короткая фраза), заданный
URI; ребро — типизированная взвешенная ассоциация. Схематически запись:

\begin{lstlisting}[style=pseudo]
узлы:   /c/ru/кружка        /c/ru/посуда      /c/ru/пить      /c/ru/горячий
рёбра (concept, relation, concept, weight):
   (/c/ru/кружка, IsA,       /c/ru/посуда,  2.5)
   (/c/ru/кружка, UsedFor,   /c/ru/пить,    1.8)
   (/c/ru/кружка, HasProperty,/c/ru/горячий,1.2)
\end{lstlisting}

Типов отношений около трёх-четырёх десятков (\textsc{IsA}, \textsc{PartOf},
\textsc{UsedFor}, \textsc{HasProperty}, \textsc{RelatedTo}, \textsc{Causes},
\textsc{MotivatedByGoal}, \dots), и рёбра многоязычны (связывают понятия между
языками). \textbf{Важно для нас}: вес ребра $\approx$ мягкая уверенность из коробки,
а типы отношений прямо ложатся на нашу типизацию рёбер (глава~\ref{ch:substrate}).

\section{Как выглядят узлы ATOMIC}

ATOMIC устроен иначе, и это источник главной сложности сшивки. Его узлы — не
концепты, а \emph{события-фразы} с плейсхолдерами действующих лиц
(\textsc{PersonX}, \textsc{PersonY}). От события идут рёбра фиксированных
\emph{социально-причинных} типов. Схематически:

\begin{lstlisting}[style=pseudo]
событие: "PersonX breaks a cup"
   xIntent  -> (none / accident)     # чего PersonX хотел
   xReact   -> (sad, upset)          # что PersonX почувствовал
   xEffect  -> (cleans it up, buys a new one)   # что произошло с PersonX
   xNeed    -> (to be holding a cup) # что требовалось до
   oReact   -> (concerned)           # что почувствовали окружающие
\end{lstlisting}

То есть «кружки» как отдельного узла в ATOMIC \textbf{нет} — есть событие «разбить
кружку» и его социально-причинные следствия. ATOMIC10x — та же структура, но в
несколько миллионов примеров, порождённых языковой моделью и отфильтрованных
критиком (больше покрытия, больше шума).

\section{Сшивка: главная инженерная работа}

Соблазн сказать «форматы отличаются только на диске» — недооценка. Внутри графовой
базы это одна структура (узлы-токены, направленные типизированные рёбра), и обход от
типа ребра не меняется. Но чтобы Рин, услышав «разбил кружку», прошла от
\emph{события} ATOMIC к \emph{предметным} свойствам кружки из ConceptNet, события и
концепты надо \emph{связать}. Это entity/event linking — нетривиальный шаг.

\subsection{Мост через контент-слова и Numberbatch}

Алгоритм сшивки: из события-фразы ATOMIC извлекаем контентные слова (существительные,
глаголы) и линкуем их к концепт-узлам ConceptNet — сначала точным лемма-матчем,
затем эмбеддинг-фолбэком.

\begin{pitfall}[Не bag-of-words и не one-hot]
Грубые представления (мешок слов, one-hot) не знают синонимии: «mug» и «cup»,
«разбить» и «сломать» окажутся разными токенами, и мост выйдет дырявым. Нужен
эмбеддинг, знающий смысловую близость.
\end{pitfall}

\emph{ConceptNet Numberbatch} — родной эмбеддинг самого ConceptNet: каждому
концепт-узлу сопоставлен вектор ($\sim$300 измерений), причём смыслово близкие
концепты близки в пространстве, а языки согласованы. Это идеальный мост:

\begin{lstlisting}[style=pseudo]
def link_atomic_event_to_concepts(event_phrase, conceptnet, numberbatch):
    concepts = []
    for token in content_words(event_phrase):     # "break", "cup"
        node = conceptnet.lemma_match(token)       # 1) точный матч
        if node is None:                           # 2) эмбеддинг-фолбэк
            node = numberbatch.nearest(embed(token))
        concepts.append(node)
    return concepts     # событие ATOMIC <-> концепты ConceptNet
\end{lstlisting}

Тот же Numberbatch решает ещё две задачи (одно пространство — три применения):
член релевантности $\text{rel}(B\mid q)$ в правиле обхода (глава~\ref{ch:walker}) и
сидирование валентности узлов (регрессия аффект-лексикона в это же пространство).

\subsection{Заземление действующих лиц}

Плейсхолдеры \textsc{PersonX}/\textsc{PersonY} — не абстракции, а слоты. В рантайме
они \emph{заземляются} на конкретных людей: \textsc{PersonX}~$\mapsto$ зритель $Z$,
\textsc{PersonY}~$\mapsto$ второй пользователь проекта. Так общий социальный шаблон
превращается в запись о конкретном отношении.

\begin{workedbox}[«Разбил кружку» на эмбриональном графе Рин]
Зритель пишет «разбил любимую кружку». (1)~Событие матчится к ATOMIC-шаблону
«PersonX breaks a cup», \textsc{PersonX}~$\mapsto Z$. (2)~Через ATOMIC всплывают
социально-причинные узлы: досада ($Z$), сочувствие (окружающие), «нужна новая».
(3)~Через сшивку с ConceptNet всплывают предметные свойства: кружка — посуда, из
неё пьют, бывает любимой. (4)~Если в личной памяти Рин есть узел «$Z$ давно хотел
новую кружку» ($-$ contradicts «досада»), появляется основание для уместного
«поздравляю» вместо сочувствия. Ни одна база в одиночку этого не даёт: ATOMIC —
эмоцию, ConceptNet — предмет, личная память — контекст.
\end{workedbox}

\section{Как строить эмбрион: засеянный подграф, а не слепой импорт}
\label{sec:embryo-build}

Здесь важная инженерная развилка. \emph{Слепо залить} весь ConceptNet и весь ATOMIC
в память — плохо: получишь много шума (банальные и устаревшие связи, англоцентризм,
странности social commonsense, слабый перенос на русский, низкую релевантность
именно домашней Рин). Тогда ты будешь отлаживать не память Рин, а шумы чужой базы
здравого смысла. Но и другая крайность — \emph{писать граф руками} на 500 узлов —
тоже плоха: медленно, и такой граф слишком разрежён, чтобы на нём проявились
топологические эффекты (хабы, боттлнеки, over-squashing), которые как раз надо
тестировать (глава~\ref{ch:topology}).

Правильный — \textbf{серединный путь}: извлечь \emph{небольшой засеянный проверяемый
подграф}.
\begin{enumerate}\itemsep1pt
\item \textbf{Сид} — концепты реального мира Рин: дом, код, усталость, кошка, универ,
дота, DTI, RIN, бытовые действия, режимы общения.
\item \textbf{Расширение} — ego-граф на $k$ хопов от сида по ConceptNet/ATOMIC.
\item \textbf{Фильтр по весу ребра} — отсекает банальщину и слабые ассоциации
(вес ConceptNet $\approx$ мягкая уверенность из коробки).
\item \textbf{Ручная чистка глазами} — один раз, для пары тысяч узлов дёшево.
\item \textbf{Статус} — это \emph{источник кандидатов}, а не замороженная истина:
дальше Рин его сама растит и прунит через использование.
\end{enumerate}
Так получаем реальную структуру (для топологических тестов) $+$ реальные ассоциации
(не рукописные) $+$ проверяемость глазами $+$ релевантность (засеяно на её домены), и
обходим шум слепого импорта.

\begin{openq}[Русский ConceptNet тоньше английского]
Практический кавеат: русскоязычная часть ConceptNet заметно беднее английской.
Варианты: держать \emph{концептный backbone} в англо-пространстве (там связей
больше), а русскую поверхность накручивать сверху; либо принять более разрежённый
русский и дополнять его вручную/из других источников. Что даст лучшее качество на
домашних сценариях Рин — открыто, проверяется на стенде (глава~\ref{ch:experiment}).
\end{openq}

\section{Как эмбриональный граф превращается в память Рин}

Эмбриональный граф \emph{нейтрален по характеру}: это «то, что досталось от
предков-разработчиков», усреднённый здравый смысл. Личность и личный опыт входят не
переписыванием этих узлов, а \emph{надстройкой}: каждое общение порождает через
fusion (глава~\ref{ch:fusion}) новые узлы и рёбра с временными метками, а
консолидация закрепляет устойчивое. Со временем поверх нейтрального каркаса
вырастает индивидуальная, окрашенная память — но фундамент здравого смысла остаётся
как \emph{система координат}, относительно которой отклонения читаются как характер
(глава~\ref{ch:personality}).

\begin{intuition}[Эмбрион, а не энциклопедия]
Название «эмбриональный» выбрано не случайно. Это не готовое знание о мире, которое
Рин «зачитывает», а зачаток, минимальный каркас связей, на котором \emph{развивается}
собственная память. Врождённые узлы со временем могут и потускнеть, и быть
пересмотрены (континуум глубины консолидации, глава~\ref{ch:substrate}) — просто
медленнее приобретённых.
\end{intuition}

\chapter{Обход графа: один движок, три политики}
\label{ch:walker}

\section{Случайное блуждание как универсальный механизм}

Почти всё, что Рин делает со своей памятью, — это \emph{обход} графа роем случайных
блуждающих агентов, распространяющих активацию. Retrieval, внутренняя жизнь и сон —
не три разных движка, а \emph{три политики} одного движка, различающиеся смещением
перехода, вероятностью телепорта и набором операций.

Общее правило перехода — нормированное произведение факторов:
\begin{equation}
P(A\to B)\;\propto\;
\underbrace{w_{AB}}_{\text{сила связи}}\;\cdot\;
\underbrace{\text{rel}(B\mid q)}_{\text{релевантность запросу}}\;\cdot\;
\underbrace{\text{mc}(B\mid \mood)}_{\text{конгруэнтность настроению}}\;\cdot\;
\underbrace{\text{unc}(B)^{\beta_u}}_{\text{тяга к неопределённости}},
\label{eq:trans}
\end{equation}
где $\text{mc}$ — из \eqref{eq:moodcong}, а показатель $\beta_u$ включает/выключает
любопытство.

\begin{intuition}[Что такое «температура»]
«Температура» — одна ручка: насколько выбор следующего шага случаен против жадного.
Формально переход берётся softmax-ом от логарифма правой части \eqref{eq:trans},
делённого на температуру $T$: при $T\to 0$ — жадный шаг по сильнейшему ребру
(эксплуатация), при большом $T$ — почти равновероятный выбор соседа (исследование).
Термин из имитации отжига; мистики нет.
\end{intuition}

\section{Три политики как параметризация}

\begin{center}
\begin{tabular}{p{0.24\textwidth} p{0.10\textwidth} p{0.10\textwidth} p{0.10\textwidth} p{0.30\textwidth}}
\toprule
\textbf{Политика} & $\alpha$ (телепорт) & $\beta_u$ & $T$ & \textbf{Что запускает} \\
\midrule
retrieval (цель) & $0$ & $0$ & низкая & извлечение под запрос \\
внутренняя жизнь (любопытство) & мал. & $>0$ & средняя & inspect, fusion у неопределённых \\
сон (покрытие) & большой & $0$ & высокая & fusion/inspect у случайных \\
\bottomrule
\end{tabular}
\end{center}

\paragraph{Retrieval — целенаправленная.} Стартует от заякоренных запросом узлов,
собирает \emph{окрестность} активированного содержания. Подчеркнём (глава~\ref{ch:personality}):
обход собирает \emph{ингредиенты}, а не готовую реакцию.

\paragraph{Внутренняя жизнь — по неопределённости.} Смещена к узлам с высоким
\texttt{need\_to\_discover}, низкой уверенностью, новизной ($\beta_u>0$).
«Любопытство» — целенаправленный поиск неопределённости, не случайность.

\paragraph{Сон — телепортирующий обходчик.} Политика \emph{стохастического
покрытия}: с вероятностью $\alpha$ телепорт в случайный узел, иначе — шаг по
\eqref{eq:trans}. Это случайное блуждание с рестартом (персонализированный
PageRank). Его стационарное распределение удовлетворяет
\begin{equation}
\pi \;=\; (1-\alpha)\,P^\top \pi \;+\; \alpha\, \mathbf{r},
\label{eq:ppr}
\end{equation}
где $\mathbf{r}$ — распределение рестарта. Телепорт делает блуждание эргодическим —
это ключевое свойство для эпистемической гигиены (глава~\ref{ch:topology}).

\begin{intuition}[Три способа гулять по городу]
Retrieval — идти по адресу (есть цель, выбираешь осмысленный путь). Внутренняя
жизнь — идти туда, где интересно и непонятно (тянет к незнакомым переулкам). Сон —
время от времени садиться в случайное такси и ехать в случайный район: не потому,
что там дело, а чтобы не застревать в одном квартале. Один пешеход, три правила
выбора улицы.
\end{intuition}

\section{Что walker делает с графом: читает, предлагает, но не коммитит}
\label{sec:proposecommit}

Естественный вопрос: обходчик \emph{сам} правит граф — сливает узлы, меняет
уверенности, ретрактит? Ответ, определяющий всю конкуренцию в системе: \textbf{нет.}
Walker — это \emph{сенсор-предлагатель}, а не писатель. Мы разводим три роли, которые
легко слить в одну и потом мучиться с гонками:
\begin{itemize}\itemsep2pt
\item \textbf{чтение} (обход) — walker ходит по графу и собирает активированное
содержание; это то, чем занимаются все три политики;
\item \textbf{предложение} — итог обхода не «правка», а \emph{типизированное
предложение операции}: не TODO-текст, а структура вида
\{\texttt{op}: consolidate, \texttt{members}:[$A,B,C,D$], \texttt{rationale},
\texttt{support\_provenance}, \texttt{walker\_id}, \texttt{t}\}, которая кладётся в
буфер предложений;
\item \textbf{коммит} — отдельная, \emph{сериализованная} стадия (тот же
schema-writer\,/\,оператор консолидации, гл.~\ref{ch:fusion}) вычитывает буфер и
\emph{применяет или отклоняет} операции, дописывая событие в лог
(гл.~\ref{ch:runtime}, раздел~\ref{sec:eventlog}).
\end{itemize}

Это ровно наш принцип «выход fusion — гипотеза, а не факт», распространённый на \emph{все}
структурные операции: обход рождает кандидата, коммиттер его судит. Зачем такая
дисциплина, а не «walker сам чинит граф на месте»: \emph{во-первых}, конкурентная
безопасность — структурный писатель ровно один, значит двух одновременных структурных
правок не бывает по построению (гонки, гл.~\ref{ch:runtime}); \emph{во-вторых},
аудируемость — каждая правка проходит через event-log и объяснима реплеем;
\emph{в-третьих}, эпистемика — предложение это гипотеза, а коммит это отдельный
гейтируемый шаг, где работает независимость свидетельств и карантин
(гл.~\ref{ch:propagation}). Walker предлагает — коммиттер располагает.

\begin{pitfall}[Единственное, что walker пишет сам — но это тоже запись]
Полный запрет на запись был бы неудобен: обход должен оставлять след «здесь были, это
подсветили». Поэтому walker \emph{напрямую} трогает только поле активации\,/\,яркости
$a_v$ (и связанную телеметрию внимания). Точная формулировка важна во избежание тихих
багов: это \emph{не} «не запись», а \emph{неструктурная коммутативная запись} — не
узлы, не рёбра, не уверенности, а состояние внимания, которому всё равно нужны TTL,
затухание и \texttt{access\_reason} (гл.~\ref{ch:timescales}). «Коммутативная» значит,
что итог не зависит от порядка параллельных записей. Но конкретно для яркости \emph{не}
годится аддитивный счётчик: на хабах, через которые проходят сотни walker'ов,
$a_v\mathrel{+}=1$ улетает в космос. Берём max-CRDT: $a_v \leftarrow \max(a_v,\,e_{\text{rw}})$,
где $e_{\text{rw}}$ — энергия зашедшего обходчика. Тогда яркость ограничена сверху
энергией агента, спад предсказуем, а сотня одновременных подсветок даёт корректный
итог без блокировок. Всё, что \emph{структурно}, — только через предложение и коммит.
\end{pitfall}

\section{Фоновая разведка: один процесс, а не «день/ночь»}
\label{sec:background}

Разделение на «дневное любопытство» и «ночной обходчик» — удобная, но неточная
рамка. Точнее: есть \emph{один всегда включённый} фоновый разведчик, а «день» и
«ночь» — два конца \emph{континуума по свободным ресурсам}, а не переключатель.

Формализуем «мощности, внимание, фокус». Пусть $\rho(t)\in[0,1]$ —
\emph{foreground-нагрузка}: насколько текущий момент требует активного ответа
($\rho\approx 1$ в напряжённом диалоге, $\rho\approx 0$ в простое). Общий бюджет
$B_{\text{total}}$ делится с приоритетом переднего плана; фоновой разведке достаётся
\begin{equation}
B_{\text{bg}}=B_{\text{total}}\,(1-\rho).
\end{equation}
\emph{Фокус} — это активное множество $A_{\text{focus}}$ (узлы высокой активации
$a_v$). Параметры фонового роя — функции $B_{\text{bg}}$: число walker'ов
$\propto B_{\text{bg}}$; доля телепортов $\alpha$ и «широта» растут при падении
$\rho$; стартовое распределение интерполирует
\[
r=(1-\beta)\cdot\mathrm{anchor}(A_{\text{focus}})+\beta\cdot\mathrm{uniform},
\qquad \beta\uparrow \text{ при } \rho\downarrow.
\]
Тогда прежние «любопытство» и «сон» — это два режима одного процесса:
\emph{высокий} $\rho$ (мало бюджета) $\Rightarrow$ мало walker'ов, $\alpha\approx 0$,
старт от фокуса — локальная разведка вокруг текущей темы; \emph{низкий} $\rho$ (много
бюджета) $\Rightarrow$ много walker'ов, большой $\alpha$, старт почти равномерный —
широкий телепорт-скраб (глава~\ref{ch:sleep}).

\begin{intuition}[Мысль на перекуре]
Именно поэтому яркая ассоциация может прийти \emph{днём}, «на ровном месте»: фоновый
разведчик работает всегда, и даже на скудном дневном бюджете иногда набредает на
далёкую связь. Развёрнутый пример такой дневной находки — в главе~\ref{ch:fusion}
(House\,/\,Cast Away).
\end{intuition}

\subsection{Поля активации: быстрое внимание и медленный след — не путать}
\label{sec:brightness}

Здесь легко совершить ошибку — свалить в одно поле «яркость», «свежесть»,
«долговременную силу» и «важность». Соблазнительно, но потом это поле начнёт делать
слишком много, и, что хуже, \emph{диагностика начнёт тренировать память}: скраб, дебаг
и тесты будут поднимать ту же переменную, что и настоящее вспоминание, и случайно
укреплять узлы. Поэтому разведём \emph{разные по темпу} состояния:
\begin{itemize}\itemsep2pt
\item $a_{\text{work}}(v,t)\in[0,1]$ — \emph{рабочая яркость внимания}. Живёт
секунды-минуты, формирует зону фокуса $A_{\text{focus}}$ и вход дневных walker'ов,
быстро гаснет. Это \emph{неструктурная коммутативная} запись обходчика (max-CRDT выше).
\item $B_{\text{base}}(v)$ — \emph{медленная сила следа} (это и есть ACT-R base-level,
гл.~\ref{ch:timescales}). Живёт дни-месяцы, обновляется \emph{только} настоящими
значимыми обращениями (по \texttt{access\_reason}), влияет на decay\,/\,забывание.
\item $\text{importance}(v)$ — значимость от appraisal\,/\,домена\,/\,структуры
(гл.~\ref{ch:personality}), не путать с частотой обращений.
\item $\text{coverage\_deficit}(v)$ — «давно не инспектировали», рычаг ночного скраба.
\end{itemize}
Ключевое разделение: настоящее вспоминание поднимает \emph{и} $a_{\text{work}}$
(мгновенно), \emph{и} $B_{\text{base}}$ (чуть-чуть, медленно); а скраб\,/\,дебаг\,/\,тест
трогают в лучшем случае $\text{coverage\_deficit}$ и \emph{никогда} не укрепляют след.
Правило быстрой яркости:
\begin{equation}
a_{\text{work}}(v,t{+}\Delta t)=a_{\text{work}}(v,t)\,e^{-\Delta t/\tau(v)}\ (\text{между визитами}),
\quad a_{\text{work}}(v)\leftarrow \max\!\big(a_{\text{work}}(v),\,e_{\text{rw}}\big)\ (\text{genuine-визит}),
\label{eq:brightness}
\end{equation}
где $\tau(v)$ \emph{не} универсальна (важные\,/\,консолидированные остывают медленнее),
а $e_{\text{rw}}$ — энергия обходчика (max, не сумма, — иначе хабы взрываются).

\begin{openq}[max теряет «навязчивость» — возможно, нужен второй канал]
Честный компромисс max: он ограничивает яркость (хорошо против хабов), но \emph{теряет
повторный фокус в коротком окне}. Сто обходчиков через хаб — max правильно не взрывает;
но если один узел \emph{многократно} всплывает в текущем рассуждении, иногда важно знать
не только пик, но и \emph{интегральную экспозицию} («это навязчиво лезет прямо сейчас»).
Возможное решение — два канала: $a_{\text{work}}^{\text{peak}}=\max(\dots)$ для яркости
и отдельный \texttt{visit\_mass\_window} для краткосрочной статистики (он \emph{не}
укрепляет $B_{\text{base}}$ автоматически). На Stage 1 не обязателен; помечаем как
возможное расширение.
\end{openq}

Теперь смещение входа обходчиков. Распределение рестарта $\mathbf{r}$ из \eqref{eq:ppr}
берётся не равномерным, и день с ночью различаются \emph{формой} смещения, а не разными
движками:
\begin{itemize}\itemsep2pt
\item \textbf{День} (высокий $\rho$, радиус $r_0$): $\ \mathbf{r}\propto
a_{\text{work}}\times\text{relevance}\times\text{uncertainty}$. Обходчик эксплуатирует
текущий фокус — близкие ассоциации «на сейчас», далеко не отходит.
\item \textbf{Ночь} (низкий $\rho$, радиус $r_1>r_0$): $\ \mathbf{r}\propto
\text{uniform}+\text{coverage\_deficit}+\texttt{need\_to\_revise}+\text{low-}\pi\text{ регионы}$.
Именно так, а \emph{не} просто «шире от дневного фокуса» — иначе сон останется
заложником дневного внимания. Сон нужен, чтобы уйти туда, куда внимание \emph{никогда}
не ходит: anti-entropy scrub, телепорт-эргодичность, проверка противоречий
(гл.~\ref{ch:sleep}).
\end{itemize}
Лаконично: \emph{день — эксплойт локально у горячего, ночь — исследуй заброшенное}. Поля
разные, но связанные: $a_{\text{work}}$ ведёт дневной фокус, $\text{coverage\_deficit}$
и низкая $\pi$ ведут ночь, $B_{\text{base}}$ решает, что переживёт забывание.

\begin{intuition}[Мотыльки на остывающий свет: откуда дневные догадки]
Если дневной ретрив — «слепой матричный ток» (раздел~\ref{sec:semgate}), откуда берутся
внезапные дневные озарения? \emph{Не} из matvec. После дневного ответа подсвеченные узлы
$a_{\text{work}}$ гаснут экспоненциально минутами-часами — остаётся \emph{тёплое
послесвечение}. Одиночные фоновые RW (тот самый остаточный бюджет $B_{\text{bg}}$)
через формулу рестарта $\mathbf{r}\propto a_{\text{work}}\times\text{relevance}\times
\text{uncertainty}$ получают повышенную вероятность телепортироваться именно в эти
недавно-горячие зоны — летят на остывающий свет, как мотыльки. Кружа в тёплой области,
они сталкивают концепты на своих траекториях и \emph{роняют кандидатов} в thought-inbox
(гл.~\ref{ch:runtime}). Важно: роняют именно \emph{кандидатов} (в основном шум), которых
обязан просеять salience-гейт ($|\Delta\mood|$, ошибка предсказания) — иначе Рин завалит
себя мусорным fusion. Так разделение труда честное: matvec подсвечивает, мотыльки ползают
по послесвечению, salience решает, что вообще достойно всплыть догадкой.
\end{intuition}

\begin{intuition}[Где «живут» обходчики]
Walker'ы не объекты, которые где-то постоянно «живут». Это фоновый процесс
планировщика с бюджетом $B_{\text{bg}}$: каждый обход \emph{рождается} из распределения
рестарта $\mathbf{r}$, проходит, оставляет предложения и подсветку — и гаснет.
Persist между ними \emph{не сами walker'ы}, а поле яркости $a_{\text{work}}$ и буфер
предложений. Спросить «где живут обходчики» — всё равно что спросить, где живут волны:
волн как вещей нет, есть среда (поле $a_{\text{work}}$) и то, что её возмущает
(планировщик).
\end{intuition}

\begin{openq}[Как измерять foreground-нагрузку $\rho$]
$\rho$ управляет тем, насколько агрессивно разведчик забирает ресурсы, но как её
\emph{мерить} в рантайме — не зафиксировано. Кандидаты: наличие активного
пользовательского хода (есть незакрытая реплика $\Rightarrow$ $\rho$ высока); длина
очереди foreground-задач; «важность» текущего события (по appraisal,
глава~\ref{ch:personality}); загрузка вычислителя. Вероятно, комбинация, но веса и
динамика (как быстро $\rho$ спадает в простое) пока открыты.
\end{openq}

\section{Точки входа: как query превращается в якоря}
\label{sec:entrypoints}

Ретрив у нас двухшаговый: \textbf{шаг~1} — найти \emph{точки входа} (якоря) в граф,
\textbf{шаг~2} — растечься из них диффузией (раздел~\ref{sec:semgate}). Про шаг~2 сказано
много, а про шаг~1 — только «быстрым индексом». Распишем его, потому что это первая строчка
кода ретрива.

\paragraph{Сразу уточнение терминологии.} Query не «разбивают на эмбеддинги». Его разбивают
на \emph{семантические единицы}, и уже каждую единицу отдельно эмбеддят. Это разные вещи.

\paragraph{Кто разбивает — не LLM.} Ретрив — рефлекс (System~1, ниже LLM-слота,
гл.~\ref{ch:runtime}), поэтому превращение query в якоря должно быть дешёвым и
\emph{LLM-free} на горячем пути. Делает это \emph{тот же экстрактор, что и на write-path}
(`raw\_event $\to$ parse/extract candidates`, раздел~\ref{sec:writepath}), только без
коммита. Красивая симметрия: \emph{понимание запроса $=$ извлечение кандидатов из события
минус запись}. На Stage~1 экстрактор простой: NER $+$ ключевые фразы $+$ лемматизация
(для русского морфология обязательна — «плавки\,/\,плавок\,/\,плавками» сводятся к лемме).
LLM подключается к \emph{переформулировке} мутного запроса только на эскалации — редкий
System~2, а не норма.

\paragraph{Два сорта якорей.} Экстрактор выдаёт две разновидности, и ищутся они по-разному:
\begin{itemize}\itemsep2pt
\item \emph{Сущностные\,/\,лексические} — именованные сущности и конкретные термины
(«Азоры», «плавки», «9 сентября»). Им нужен \textbf{точный} матч, а не «похожее по вектору».
\item \emph{Холистический} — эмбеддинг всего запроса (или его клауз) для ассоциативной
диффузии.
\end{itemize}
Если единиц несколько — это \emph{мульти-векторный} ретрив, а не одно усреднённое. Mean-pool
«плавки на Азоры 9 сентября» размазал бы три резкие сущности в кашу; поэтому каждую единицу
сажаем отдельным seed'ом, а узлы, подсвеченные \emph{несколькими} seed'ами, получают буст
(пересечение $=$ релевантнее).

\paragraph{Чем эмбеддим.} Тем же энкодером, что и \texttt{node.embedding} при записи
(гл.~\ref{ch:substrate}). Это не опция: косинус осмыслен, только если query и узлы в
\emph{одном} пространстве. Вектор узла уже кэширован (считался один раз коммиттером), так
что на шаге ретрива через энкодер гоним только query-единицы, узлы не трогаем.

\paragraph{Как ищем точки входа — два канала, потом объединяем.}
\begin{lstlisting}[style=pseudo]
units      = extract(query + active_focus)      # NER + keyphrases + леммы (+ кореф)
anchors    = {}
for u in units:
    if u.is_named_entity:
        anchors |= exact_index.lookup(lemma(u))  # точный/alias-индекс (forensic-подобный)
    anchors |= ann_index.knn(embed(u), k)        # вектор-ANN (HNSW), top-k по косинусу
anchors    = resolve_entities(anchors)           # упоминания -> канонические узлы
r          = weight_by_match_confidence(anchors) # точное > сильный вектор > слабый
# r -> распределение рестарта для локальной разреженной диффузии (шаг 2)
\end{lstlisting}
То есть точки входа приходят из \emph{двух} индексов: точного\,/\,alias (для именованных
сущностей — «Азоры» должны попасть в \emph{тот самый} узел) и вектор-ANN (для семантической
близости). Их объединение, взвешенное по уверенности матча, и есть seed-набор. Плюс
\emph{entity resolution}: разные упоминания сводятся к каноническому узлу, чтобы «плавки»
били в один узел, а не в пять эпизодов покупки.

\paragraph{Дальше.} Seed-набор становится \emph{начальным впрыском} $\mathbf{a}(0)=s(0)$
для \emph{локальной затухающей} диффузии (fast-path вокруг якорей, раздел~\ref{sec:semgate}),
а расширение гейтируется условным throttle. Важно: retrieval — это \emph{политика} с
телепортом $\alpha=0$ (таблица политик выше), то есть источник впрыскивается \emph{один
раз} и импульс затухает; рестарт (член $s(t{>}0)$, он же $\alpha>0$ в \eqref{eq:ppr}) —
это \emph{сон}, а не retrieval. Цепочка целиком: \emph{экстрактор дал якоря $\to$
ANN$+$точный индекс нашли точки входа $\to$ импульс затух $\to$ packer срезал по порогу}.

\begin{openq}[Что на Stage 1 сознательно упрощаем]
Открыты: гранулярность разбиения (по сущностям\,/\,клаузам\,/\,предложениям — тюнить на
бенче); \emph{кореференс} — «поищи \emph{их}» требует обогатить query текущим фокусом
$a_{\text{work}}$ \emph{до} эмбеддинга, иначе якорей нет; порог эскалации к
LLM-переформулировке; качество alias-резолвинга под русскую морфологию. Механика
(два сорта якорей, два индекса, мульти-seed, resolution) — зафиксирована.
\end{openq}

\section{Семантический вентиль: дёшево и без протечек по омонимии}
\label{sec:semgate}

В \eqref{eq:trans} множитель $\text{rel}(B\mid q)$ — это \emph{семантический вентиль}:
он не даёт фронту расползтись из зоны запроса. Возникает законный страх за
производительность: если на каждом шаге RW считать сходство контента узла с $q$ заново,
это катастрофа. Ответ: \textbf{на лету ничего не эмбеддим.} Эмбеддинг узла считается
\emph{один раз при записи} (коммиттером) и лежит в поле \texttt{embedding}
(гл.~\ref{ch:substrate}). Тогда вентиль на шаге — это скалярное произведение
\emph{фиксированного} вектора $q$ на \emph{кэшированный} вектор узла, $O(d)$; а раз
ретрив у нас — усечённая \emph{локальная затухающая} диффузия из seed'ов (политика
retrieval: телепорт $\alpha=0$, источник впрыскивается один раз — это \emph{не} PPR;
рестарт-форма \eqref{eq:ppr} зарезервирована за сном), весь фронт считается одним matvec,
где вектор активации на каждом шаге умножается на маску сходства с $q$. Как только значения падают ниже порога $\varepsilon$ — граница
замкнулась. Никаких рекурсивно плодящихся агентов: «группа вспышек, слабеющих к
окраинам» — это \emph{результат} диффузии, а не объектная механика.

Инженерная честность про этот дневной fast-path: matvec идёт \emph{не} по всему графу.
Умножать $\tilde P^\top$ на $|V|$-мерный вектор на графе в сотни тысяч узлов — лечь. Это
\emph{локальная разреженная} диффузия на ограниченном носителе: 2--3 хопа вокруг
seed-якорей, найденных ANN-индексом (HNSW), — субматрица смежности окрестности, а не всё
$V$-пространство. И тайминги честные — \emph{миллисекунды}, не микросекунды (основное время
— ANN-поиск якорей и context-packer), а настоящее бутылочное горло рантайма всё равно
LLM-инференс (секунды), гл.~\ref{ch:runtime}.

Две поправки, чтобы вентиль не навредил. \emph{Первая — жёсткость по типу ретрива}: для
ordinary-ретрива режем жёстко (защита от мусора), для creative\,/\,curiosity — мягко
(именно далёкие странные мосты нужны для fusion), для forensic — вентиля нет вообще, это
не walker, а точный индекс-lookup (таксономия ретрива, гл.~\ref{ch:runtime}). Слишком
жёсткий вентиль убил бы серендипность. \emph{Вторая — полисемия}: граница возникает не
только из низкого сходства с $q$, но и из \emph{sense/type-несовпадения}. Классический
пример «кружка $\to$ ручка $\to$ дверь» не должен решаться тем, что «дверь слабо похожа
на кружку»: раньше ломается слой смысла. «Ручка кружки», «дверная ручка», «ручка-перо» —
это \emph{разные} sense/concept-узлы (или хотя бы квалифицированные рёбра,
гл.~\ref{ch:embryonic}). Без этого ретрив постоянно течёт по омонимии — а для русского
это особенно больно. Значит вентиль закрывается по $\max(\text{low-sim},\
\text{sense/type-mismatch})$, а не по одной similarity.

\subsection{Вентиль — условный throttle, а не плоский множитель}
\label{sec:semgate-conditional}

Есть и \emph{третья} поправка, и она важнейшая: как написано выше, $\text{rel}(B\mid q)$
входит в \eqref{eq:trans} \emph{множителем} — а значит зануляет путь при низком сходстве.
Это тихо убивает личность. Пример: триггер «гирлянда». Узлы [свет], [провода] семантически
близки и проходят; а узел личной памяти [окно в Тбилиси] имеет сходство $\approx 0$ — и
множитель его гасит, Рин выдаёт энциклопедию. В живой голове [Тбилиси] вспыхивает
\emph{первым}: у личного ребра колоссальный вес\,/\,валентность, несмотря на семантическую
даль. Значит вентиль должен быть \emph{условным throttle}, а не вето:
\[
\text{Score}(A\to B)=f(w_{AB},\,\mathrm{valence}_B)\cdot
\underbrace{\text{Throttle}(B,q)}_{\text{условный}}\cdot \text{mc}(B\mid\mood)\cdot \text{unc}(B)^{\beta_u},
\]
\[
\text{Throttle}(B,q)=
\begin{cases}
1, & \text{ребро из личного слоя (\texttt{session}/\texttt{life}) } \textbf{ИЛИ } |\mathrm{valence}_B|>\theta_{\text{emotion}},\\[2pt]
\text{rel}(B\mid q), & \text{иначе (обычное знание — тормоз против дрейфа).}
\end{cases}
\]
То есть сильное личное или аффективно заряженное ребро — «синаптический хайвей» —
проходится независимо от сходства с запросом; а throttle тормозит только слабые\,/\,родовые
рёбра, чтобы не было случайного топологического дрейфа.

Три предохранителя, без которых это сломается. \emph{Не по сырому $w_{AB}$}: пробивать
вентиль порогом веса нельзя — у эмбриональных рёбер (ConceptNet) веса из коробки огромны
($w\approx 2{,}5$), и Рин уплыла бы по банальной таксономии. Поэтому гейт по
\emph{происхождению} ребра (личный слой, не врождённый каркас), а не по величине веса.
\emph{$f(w_{AB},\mathrm{valence})$ всё равно ранжирует}: throttle$=1$ \emph{не} значит «все
личные узлы хлынут» — Тбилиси всплывает не потому, что ребро личное, а потому, что оно
личное \emph{и} с огромным $w_{AB}$; слабое личное ребро с throttle$=1$ всё равно проиграет
по $f$. \emph{Стоп-условие три-гейта}: путь глохнет, только если \emph{одновременно}
низкое сходство \emph{и} слабое ребро \emph{и} нейтральная валентность; сильна любая из трёх
— агент проходит. Сухой-но-важный \texttt{boundary} («не шути про доту») проходит через
layer-гейт даже при низкой валентности. Омонимию (кружка$\to$ручка$\to$дверь) по-прежнему
режет отдельный sense/type-гейт, а не throttle.

\begin{intuition}[Прайминг: подсвечено, но не сказано]
База при этом не пропадает: «у гирлянды есть провода» матричная диффузия \emph{подсвечивает}
(ток туда течёт) — это прайминг. Но у этих узлов почти ноль ситуативной значимости и
новизны, поэтому context-packer их \emph{сознательно не выносит наверх}
(раздел~\ref{sec:packer}). Они «тёплые» на фоне сознания, но не доходят до языка — ровно как
человеческое подсознание держит очевидное наготове, не проговаривая.
\end{intuition}

\begin{principle}[Вентиль на эмбеддингах — фильтр родственности, не релевантности]
\emph{Эмпирически выяснено} (стенд, \texttt{reports/experiment\_3--4}): множитель
$\text{rel}(B\mid q)$, построенный на дистрибутивных эмбеддингах (косинус к вектору
запроса), фильтрует \emph{родственность} (co-occurrence), а не \emph{причинную
релевантность}. Он давит off-topic дрейф (годится для запросов вида «что есть X» / «кто
из этих X»), но на «почему»-запросах систематически недобирает: узел может быть ближе к
слову запроса как \emph{ассоциат}, не будучи \emph{ответом} («beverage» ближе к «energy»,
чем «caffeine», из-за маркетингового со-упоминания). Это ограничение \emph{инструмента},
а не бага; для причинного выбора нужен не косинус, а рассуждение.
\end{principle}

\begin{principle}[Конвергенция на хабе пересиливает per-edge релевантность; лечится редкостью, не гейтом]
\emph{Эмпирически выяснено} (там же): при диффузии узел, достижимый \emph{из нескольких}
якорей сразу, накапливает их вклады, и это суммирование может перевесить более высокий
per-edge $\text{rel}(B\mid q)$ у односвязного правильного ответа — даже когда вентиль
оценил ответ выше. То есть провал «генеричный дистрактор победил» решает \emph{конвергенция
на хабе}, а не релевантность; это тот же зверь, что концентрация $\pi$ на хабах
(гл.~\ref{ch:topology}). Значит корректив — \emph{отдельный штраф за редкость\,/\,степень}
(IDF-подобный: генеричный общий сосед давится), а не докрутка вентиля. Один такой член
чинит сразу три наблюдения: дистрактор-хаб, таксономический дрейф ретрива и хаб-концентрацию.
\end{principle}

\begin{principle}[Ретрив отвечает за recall, LLM — за причинный выбор]
Отсюда критерий успеха ретрива: не «правильный узел \emph{первый} в ранге», а «правильный
узел \emph{попал} в упакованный контекст в рамках бюджета при приемлемой доле дистракторов».
Граф — ассоциативная память (поднять кандидата в окрестность); \emph{причинное}
ранжирование («из этих caffeine — и есть \emph{почему}») делает LLM на собранном контексте.
Требовать от графа самому назвать причинный ответ — значит возлагать рассуждение на recall.
\end{principle}

\section{Единство и следствие}

Один субстрат, один блуждающий агент, три параметризации: смещение перехода,
вероятность телепорта, набор операций. Реализуется одна машина обхода с параметрами
(псевдокод — приложение~\ref{ch:algorithms}). Родственные механизмы: spreading
activation (Collins–Loftus) как модель ассоциативного припоминания и
персонализированный PageRank / RWR как его вычислительная форма; в волте RIN — под
именами SYNAPSE и GAAMA.

\chapter{Fusion, вывод и консолидация}
\label{ch:fusion}

Обход достаёт то, что \emph{уже} есть в графе. Но память живого существа ещё и
\emph{порождает} новое: связывает далёкое, делает выводы, обобщает. Этот класс
операций мы называем \emph{fusion}. Здесь разбирается, как именно из активированных
узлов рождается новый узел, чем вывод отличается от ассоциации и обобщения и как всё
это остаётся под эпистемическим контролем.

\section{Что такое fusion, механически}

\begin{definition}[Fusion, подробно]
\emph{Fusion} — оператор, который по \emph{совместно активированной} конфигурации
узлов $\{A_1,\dots,A_n\}$ создаёт новый узел $X$ и связывает его с посылками рёбрами
провенанса (обоснования). Новый узел рождается \emph{провизорно}: с низкой
уверенностью и, при необходимости, с флагом \texttt{need\_to\_discover} на том, что
его различает.
\end{definition}

Механика в четырёх шагах:
\begin{enumerate}
\item \textbf{Со-активация.} Обход (retrieval, любопытство или сон) подсвечивает
несколько узлов одновременно.
\item \textbf{Распознавание конфигурации.} Оператор замечает, что активная
конфигурация \emph{согласована} — образует осмысленный паттерн (причинную цепочку,
аналогию, повтор).
\item \textbf{Рождение узла.} Создаётся $X$ с содержанием (сформулированным
языковой моделью) и провенансом $\{A_1,\dots,A_n\}$.
\item \textbf{Провизорная запись и пропагация.} $X$ пишется с малым $\conf$, затем
пропагируется со знаком, карантином и дисконтом независимости
(глава~\ref{ch:propagation}).
\end{enumerate}

\begin{principle}[Fusion рождает гипотезу, а не факт]
Результат fusion — всегда \emph{кандидат}, а не запись. Это защищает от превращения
обаятельной ассоциативности в фабрику ложных убеждений
(глава~\ref{ch:epistemics}). Повышение до «факта» — только через накопление
подтверждений или выживание в ночном скрабе.
\end{principle}

\section{Три вида fusion}

\paragraph{Ассоциация.} Два далёких узла со-активируются, и рождается связь (или
узел-мостик). Это источник «неожиданных, но осмысленных» ходов мысли.

\begin{workedbox}[Ассоциация: Платон-боец]
Ночной обходчик случайно со-активирует два узла: «Платон был олимпийским чемпионом
по панкратиону (боец/борец — по свидетельству Олимпиодора)» и «\emph{platon}
по-гречески — „широкий, широкоплечий“». Fusion замечает согласование «мощное
телосложение бойца» $+$ «имя означает „широкий“» и рождает \emph{гипотезу}:
«„Платон“ — прозвище по телосложению, а не данное имя», $\conf=0{,}2$, провенанс из
двух узлов, флаг \texttt{need\_to\_discover} на «проверить: прозвище ли это и каково
настоящее имя». Вывод правдоподобен (версий, к слову, несколько, но все сводятся к
трактовке слова «широкий», а настоящее имя по преданию известно) — но именно
\emph{гипотеза}: если не пометить низкой уверенностью, это тот самый преждевременный
коллапс (глава~\ref{ch:epistemics}). Существенно, что связывает не «Платон-философ»
сам по себе, а именно пара «боец» $+$ «перевод».
\end{workedbox}

\paragraph{Вывод (инференция).} Из посылок рождается заключение. Здесь принципиально
важен случай \emph{n-арного} вывода, где ни одна часть посылок не достаточна.

\begin{workedbox}[N-арный вывод: корм котам]
Пусть в памяти Рин про пользователя есть три узла, и \textbf{каждый подтягивает свой
контекст} из эмбрионального графа (см.\ раздел ниже):
\begin{itemize}\itemsep2pt
\item $A$ — «у него есть коты» $\Rightarrow$ контекст: \emph{котов надо кормить,
нужен корм};
\item $B$ — «живёт в Португалии» $\Rightarrow$ контекст: \emph{магазины закрываются
рано};
\item $C$ — «иногда ездит на пары» $\Rightarrow$ контекст: \emph{с пар возвращается
поздно}.
\end{itemize}
Только совмещение всех трёх \emph{с их контекстами} даёт вывод:
\[
X=\text{«купи котам корм на выходных — в будни по пути с пар можешь не успеть до закрытия»}.
\]
Существенно, что \emph{ни одна пара не достаточна}: без «магазины рано» нет спешки,
без «поздно с пар» нет проблемы, без «нужен корм» нет повода. Вывод рождается лишь
при \emph{одновременном} присутствии трёх посылок и их контекстов.
\end{workedbox}

Как мы это представляем \emph{сейчас}, не меняя структуру: fusion создаёт узел $X$ с
провенансом $\{A,B,C\}$. Тем самым тройка-посылка кодируется как узел-фактор,
ссылающийся на три посылки, — это \emph{factor-graph}-кодировка гиперребра
$\{A,B,C\}\to X$. Иначе говоря, наша память уже наполовину гиперграф: попарные
рёбра $A\!-\!B$, $B\!-\!C$ могут отсутствовать (Португалия и коты не связаны
напрямую), а n-арный вывод всё равно рождается — через общий узел-обоснование.
Стоит ли делать гиперрёбра \emph{первоклассными} объектами (со своим обходом и
лапласианом) — тема приложения~\ref{ch:hypergraphs}.

\begin{center}
\begin{tikzpicture}[font=\scriptsize,>=Stealth,
  n/.style={draw,rounded corners,minimum height=6mm,inner sep=2pt}]
  \node[n] (a) {A: Португалия};
  \node[n,below=5mm of a] (b) {B: коты};
  \node[n,below=5mm of b] (c) {C: универ};
  \node[n,right=26mm of b,fill=blue!5] (x) {X: купи корм на выходных};
  \draw[->,dashed] (a) -- node[above,font=\tiny]{provenance} (x);
  \draw[->,dashed] (b) -- (x);
  \draw[->,dashed] (c) -- (x);
\end{tikzpicture}
\captionof{figure}{N-арный вывод как узел-фактор $X$ с провенансом $\{A,B,C\}$.
Попарных рёбер между $A,B,C$ может не быть.}
\end{center}

\subsection*{Как посылка «подтягивает свой контекст»}

Важная деталь, без которой n-арный вывод не работает: посылки участвуют не «голыми»,
а вместе со своим \emph{контекстным ореолом}. Механически это ограниченная (1--2
хопа), смещённая целью/настроением \emph{spreading activation} от каждой посылки в
эмбриональный граф:
\[
\begin{aligned}
A=\text{коты} &\Rightarrow \{\text{нужен корм}\},\\
B=\text{Португалия} &\Rightarrow \{\text{магазины рано закрываются}\},\\
C=\text{универ} &\Rightarrow \{\text{поздний возврат}\}.
\end{aligned}
\]
Fusion работает над \emph{объединением} этих ореолов, а не над тремя голыми узлами;
композитор (языковая модель) собирает из «нужен корм $+$ магазины рано $+$ поздний
возврат» связный actionable-вывод. Провенанс $X$ — посылки \emph{и} реально
использованные контекст-узлы. Ореол смещён текущей целью: иначе подтянулась бы вся
Вселенная соседей, поэтому активируется лишь релевантное намерению (тот же
mood/goal-bias, что в retrieval, глава~\ref{ch:walker}). Именно поэтому эмбриональный
граф (глава~\ref{ch:embryonic}) не декорация: без рёбер ConceptNet/ATOMIC у посылок
нет ореолов — и вывод не рождается.

\begin{openq}[Трудная часть — распознать связность («трудные 80\%»)]
Подтянуть ореолы дёшево; \emph{трудно} распознать, что их объединение образует
\emph{связную} цепочку рассуждения, а не случайную кучу со-активаций. Кандидаты
разобраны в главе~\ref{ch:timescales} (раздел~\ref{sec:whatfusiondoes}): LLM-судья /
графовые шаблоны-схемы / обучаемый скорер связности — ни один пока не признан
удачным. На старте — LLM-композитор предлагает, а $\varphi$-скорер и обратная связь
фильтруют (главы~\ref{ch:personality},~\ref{ch:voice}). Это открытый инженерный
вопрос.
\end{openq}

\paragraph{Обобщение (абстракция).} Много однотипных эпизодов сливаются в один узел
более высокого уровня. Это и есть консолидация (следующий раздел).

\section{Fusion сцепляется: sub-fusion — это тоже fusion}

Виды fusion — не разные механизмы, а один оператор, порождающий узлы разного
\emph{рода}. И fusion \emph{сцепляется}: обобщение, построенное одним шагом
(«sub-fusion»), становится посылкой для вывода на следующем. Абстракция
\emph{поднимает уровень}, вывод делает конкретный скачок, опираясь на неё:
\[
\text{факты}\ \xrightarrow{\text{вверх (абстракция)}}\ \text{обобщение}\
\xrightarrow{\text{скачок (вывод)}}\ \text{гипотеза}.
\]
Механика та же (со-активация $\to$ новый узел с провенансом); «sub» лишь означает,
что этот узел — промежуточное \emph{обобщение}, мост для последующего шага. (Это, к
слову, клеточная картинка из приложения~\ref{ch:cellcomplex}: высшая клетка,
приклеенная к низшим.)

\begin{workedbox}[Дневная находка и полная эпистемическая петля: House / Cast Away]
\textbf{Фон.} Фоновый разведчик (глава~\ref{sec:background}), даже на скудном дневном
бюджете, со-активирует два далёких узла: «в \emph{House M.D.} есть Уилсон» и «в
\emph{Cast Away} есть „Уилсон“ — волейбольный мяч».

\textbf{Sub-fusion (абстракция).} Оба узла порождают обобщение: «Уилсон — приём
\emph{экстернализации внутреннего монолога} ГГ» (у Cast Away — из-за одиночества на
острове; у House — потому что ГГ умнее всех, детектив, интроверт). Родился узел более
высокого уровня — общая сценарная функция.

\textbf{Fusion (скачок).} Эта абстракция $+$ совпадение имени «Уилсон» рождают
\emph{гипотезу}: «а вдруг имя Уилсон в House — отсылка к Cast Away?!» Провизорно,
низкая уверенность.

\textbf{Проверка (петля эпистемики, глава~\ref{ch:epistemics}).}
\texttt{need\_to\_discover} на «годы выхода» $\to$ внешняя проверка: Cast Away (2000)
раньше пилота House (2004) — совместимо, уверенность \emph{растёт}. Затем
\texttt{need\_to\_discover} на «есть ли ещё свидетельства» $\to$ проверка
опровергает: официальная версия — оммаж Шерлоку Холмсу (House\,$\approx$\,Holmes,
адрес 221B), а Джеймс Уилсон\,$\approx$\,доктор Ватсон (одна буква, созвучие).
Гипотеза \emph{понижается}, на её место встаёт лучше обоснованная.

\textbf{Мораль.} Красивый, согласующийся с частью данных вывод оказался ложным —
это тот самый преждевременный коллапс, пойманный проверкой. И заметьте: «правильный»
ответ (Holmes\,$\to$\,House) — сам по себе fusion, просто с лучшим провенансом. Как
говорит ГГ, «все лгут» — поэтому уверенность держим ниже единицы и идём проверять.
\end{workedbox}

\section{Консолидация как огрубление}

\begin{definition}[Консолидация, подробно]
\emph{Консолидация} — оператор, заменяющий множество похожих узлов/рёбер одним
обобщённым, сохраняя релевантные свойства и отбрасывая детали. Формально это
\emph{огрубление графа} (coarse-to-local, глава~\ref{ch:timescales}).
\end{definition}

\begin{workedbox}[Консолидация: «$Z$ — частый гость»]
Десять эпизодов «$Z$ был на стриме в понедельник», «\ldots в среду», \ldots
сворачиваются в один узел «$Z$ — частый гость, любит доту». \emph{Выживают}:
частота, тёплая валентность, тема (дота). \emph{Теряются}: точные даты и реплики.
Плохое огрубление сохранило бы даты и потеряло валентность — Рин помнила бы,
\emph{когда} заходил $Z$, но забыла бы, \emph{как к нему относится}. Значит критерий
консолидации — сохранять аффективно- и реляционно-релевантное, а не всё подряд.
\end{workedbox}

Консолидация запускается двумя путями: во сне (фоновая переработка,
глава~\ref{ch:sleep}) и забыванием (когда эпизоды тускнеют, их дешевле слить в
обобщение, чем хранить порознь, глава~\ref{ch:timescales}).

\section{Кто пишет в узел и в каком виде}
\label{sec:whowrites}

Отдельно разберём вопрос, который легко замять: кто именно порождает содержание
нового узла — LLM, автомат или Rinness — и в каком \emph{виде} оно ложится в память.
Полезно разбить на три подшага, потому что механизмы у них разные.

\paragraph{Решение (сливать ли).} Для fusion это распознавание связной
конфигурации (открытый «трудный 80\%», раздел~\ref{sec:whatfusiondoes}); для
консолидации — обнаружение избыточного кластера. И тут важная поправка: группировать
«однотипное» \emph{по строке} нельзя. Зритель мог «заходить», а мог «забегать» —
чистый автомат, ключующийся по глаголу, тут же ломается, ведь мы не типизируем
глаголы, которыми к Рин обращаются. Синонимию ловит не автомат, а либо LLM при
записи (нормализует «забежал» в канонический тип связи \textsc{visited}), либо
кластеризация по близости эмбеддингов при консолидации. Поэтому честная картина:
автомат делает бухгалтерию (счётчики, хаб, перекидка рёбер, усреднение валентности),
эмбеддинги — группировку «похоже по смыслу», а маленький LLM — итоговый ярлык.
«Только автоматом» проходит лишь случай точных дублей.

\paragraph{Содержание узла.} Новое \emph{осмысленное} содержание (вывод, абстракция,
ярлык обобщения) пишет \textbf{языковая модель} — шаблоном такое не собрать. Но
\emph{та же самая} модель, а не какая-то отдельная: «нейтральный режим» и «Rinness»
— это два \emph{настроения} одной модели, а не два мозга. Для записи в память её
держат в спокойном, откалиброванном «девелоперском» режиме вектора настроения:
рассудочно, минимум стиля. Если голос Рин зашит в веса (LoRA), совсем нейтральной её
не сделать — останется лёгкий налёт, и это не страшно, лишь бы не перегибать: в
память важно \emph{содержание}, а тон всё равно родится заново на выходе. Автомат
при этом ставит все метаданные (провенанс, уверенность, \texttt{need\_to\_discover},
importance, домен) — детерминированно.

\paragraph{Кто именно пишет — и почему это не «Rinness».} Пишет в граф \emph{та же
Рин}, но в спокойном \emph{режиме записи}, а не в режиме голоса. Здесь важно не
путать «Рин» и «Rinness»: \emph{Rinness} — это узко \emph{выходная стилизация}
(голос, гл.~\ref{ch:voice}), которая превращает извлечённое содержание в реплику
пользователю. Вот эта стилизация в граф \emph{не} пишет. Одна Рин, два режима: режим
записи (спокойный) кладёт нейтральное содержание, режим голоса (Rinness) звучит на
выходе. Отсюда принцип: \emph{храни спокойно, стилизуй на выходе}. Причина простая:
узел переиспользуется в разных будущих настроениях, и если вморозить в него
сегодняшний тон, он будет отдавать им и через месяц, когда Рин к тому же зрителю
потеплеет.

\subsection*{В каком виде хранить: содержание против стиля}

Итак, в память кладём \emph{нейтральное содержание} («зритель $Z$ проиграл в доте»)
плюс её \emph{отношение} в метаданных — валентность и, для значимых случаев,
отдельный узел-реакцию. Личность памяти живёт в отношении, а не в замороженных
словах. Класть готовую фразу «в голосе» («$Z$ снова слил, ну куда ему») — плохо:
это вмораживает транзиентный тон, память пухнет, а стилизовать на выходе уже нечем.

Исключение — \emph{вспышечная память}. У людей рядовые события помнятся по сути, а
редкие пиково-эмоциональные — вплоть до того, что было сказано и как ощущалось. По
аналогии: для \emph{редких} моментов с высокой значимостью (большой $|\Delta\mood|$)
можно дополнительно сохранить \emph{дословный снимок} — реальную реплику/тон того
мгновения — как отдельный связанный узел, а не как каноническое содержание. Так по
умолчанию память нейтральна и переиспользуема, но яркие пики сохраняют свою фактуру.

\section{Выход fusion — гипотеза, а не установленный факт}

Снимем терминологическую путаницу, потому что от неё зависит понимание всего.
Каждый узел — это \emph{убеждение} с уверенностью $\conf\in[-1,1]$. «Факт» в широком
смысле — это просто любой узел, и да, fusion \emph{создаёт} узел. Но узлы
различаются \emph{по уверенности}: прямое наблюдение или измерение получает высокую
уверенность, а вывод (fusion) — \emph{низкую}. Именно это имелось в виду под «fusion
не коммитит факты»: \emph{коммит} — это запись как \emph{установленного},
высокоуверенного убеждения, которым дальше пользуются как истиной; fusion же рождает
узел-\emph{гипотезу}. Повышение гипотезы до установленного статуса происходит
\emph{позже} — через корроборацию или проверку, а не в момент рождения.

Почему это критично: без такого разделения обаятельная ассоциативность превратилась
бы в фабрику ложных убеждений (пример House\,/\,Cast Away выше). Поэтому fusion смел,
но не безрассуден: гипотезы рождаются свободно, но влияние гипотезы на соседей
гейтируется её собственной (низкой) уверенностью и надёжностью посылок — карантин,
глава~\ref{ch:propagation}; а различающие пробелы, которых Рин не знает, помечаются
\texttt{need\_to\_discover}, чтобы любопытство или внешний вход их потом добыли
(глава~\ref{ch:epistemics}). Дальнейшую судьбу гипотезы — проверку, повышение или
понижение — ведёт уже эпистемический слой; fusion лишь честно рождает её с низкой
уверенностью и передаёт дальше.

\begin{principle}[На первых версиях: генератор кандидатов с аудитом]
Fusion — главный источник и магии, и мусора. Пока распознавание связности
со-активации не решено (открытый вопрос), fusion на ранних версиях должен быть не
«автоматической когнитивной магией», а \emph{генератором кандидатов с аудитом}.
Дефолтная позиция смещена в сторону осторожности: \textbf{больше кандидатов в
pending, меньше commit}. Принятые fusion потом \emph{ревизуются} отложенно (precision
после проверки, глава~\ref{ch:experiment}), и только выжившие повышаются. Лучше
недодать смелых узлов, чем зарастить память красивым мусором; смелость безопасно
включим, когда научимся мерить связность.
\end{principle}

\begin{principle}[Разведи candidate\_text и canonical\_node]
Fusion всё ещё сильно опирается на LLM-композитор, и это риск: модель выдаёт красивый
текст, но плохо индексируемый и плохо ревизуемый узел. Разделяем два формата. LLM
формулирует \texttt{candidate\_text} — гипотезу человеческим языком. Но \emph{коммит}
идёт не текстом, а через \emph{schema-writer} в \texttt{canonical\_node} со
структурой: \{\texttt{node\_kind}, \texttt{claim\_type}, \texttt{entities},
\texttt{time\_scope}, \texttt{confidence}, \texttt{provenance},
\texttt{distinguishing\_unknowns}\}. Иначе получишь хорошо написанные, но
нетипизированные узлы, которые нельзя ни точно найти (forensic-ретрив,
гл.~\ref{ch:runtime}), ни аккуратно ревизовать. Красота — на входе, схема — на
коммите.
\end{principle}

\section{Когда со-активация — повод для реплики, а не для узла-факта}
\label{sec:antifusion-joke}

Не всякая яркая со-активация должна оседать убеждением. Иногда две всплывшие вещи —
это повод \emph{что-то сказать}, а не повод \emph{что-то запомнить как факт о мире}.
Это ключевое различение: \textbf{полезное для реплики $\ne$ достойное записи фактом}
(anti-fusion). Разберём на каноничном случае.

\begin{workedbox}[Чапаев и «нюанс»: шутку не пишем, реакцию — пишем]
В памяти со-активировались «Чапаев» и слово «нюанс». Walker сводит их рядом, и
всплывает ассоциация: «тут просится анекдот». Что происходит по нашим механизмам.

\emph{(1) Откуда берётся анекдот.} Не из графа как хранимый текст. Граф даёт только
\emph{триггер} — «эта пара уместна для шутки». Сам текст анекдота порождает базовая
LLM (она знает его из предобучения) — ровно как с обучением: граф всплывает
\emph{повод}, модель поставляет \emph{содержание}.

\emph{(2) Что НЕ записывается.} Сам анекдот \emph{не} коммитится узлом-фактом. Связь
«Чапаев relates-to нюанс через этот анекдот» — это \texttt{candidate\_text} для
\emph{реплики}, а не убеждение о мире; anti-fusion-гейт его как факт не пропускает.
Иначе память заросла бы «фактами», которые на деле каламбуры.

\emph{(3) Что записывается.} Два узла другого рода. Во-первых, \texttt{episode}:
«рассказала анекдот про Чапаева-нюанс собеседнику $Z$» (со временем и провенансом) —
это \emph{было}, это факт о событии, не о мире. Во-вторых, реакция $Z$ обновляет
\texttt{reaction\_pattern} через Beta-постериор, условленный на отношении
(гл.~\ref{ch:personality}): зашло $\to\alpha{+}1$, поморщился $\to\beta{+}1$.
Salience-гейт решает судьбу: сдвинула реакция настроение ($|\Delta\mood|$) —
консолидируется; «ну ок» — тускнеет.

\emph{(4) Бонус — крючок, а не шутка.} Если такой юмор у $Z$ стабильно заходит, осесть
может лёгкий \texttt{reaction\_pattern}: «пара Чапаев+нюанс $\to$ этот анекдот заходит
$Z$». Заметьте: выучивается \emph{модель того, что кому заходит}, а не сам анекдот
хранится фактом. Комедийный тайминг — Beta-апостериор, не строка в базе.
\end{workedbox}

Мораль в одну строку: fusion-подобная со-активация может кончиться \emph{репликой}, а
осесть в памяти должны \emph{событие речи} и \emph{реакция на неё} — потому что важна
не шутка, а то, как она легла. Тот же принцип защищает память от превращения в свалку
острот, мемов и подходящих-к-случаю ассоциаций.

\begin{intuition}[Инженер, а не фантазёр]
Правильная Рин — «инженер, которую делают бывшие физики-фундаменталисты»: она
охотно строит гипотезы, но по умолчанию ставит уверенность ниже единицы и идёт
проверять. Fusion даёт первое (смелость ассоциаций), эпистемическая гигиена
(глава~\ref{ch:epistemics}) — второе (дисциплину). Одно без другого — либо сухой
поисковик, либо обаятельный врун.
\end{intuition}

\chapter{Многомасштабность, слои и забывание}
\label{ch:timescales}

\section{Протекающий интегратор: точный смысл}

Память Рин расслоена по временным масштабам: \emph{сессия} $\to$ \emph{день} $\to$
\emph{сезон} $\to$ \emph{год} $\to$ \emph{жизнь}. Каждый слой — \emph{протекающий
интегратор} (leaky integrator): состояние, которое накапливает вход и понемногу
«протекает» (теряет часть себя) на каждом шаге.

\begin{definition}[Протекающий интегратор]
Состояние $s$, обновляемое под входом $x_t$ по правилу
\begin{equation}
s_t \;=\; \underbrace{\Big(1-\tfrac{1}{\tau}\Big)}_{\text{удержание}}\,s_{t-1}
\;+\; \underbrace{\tfrac{1}{\tau}}_{\text{впитывание}}\,x_t,
\qquad \tau>1.
\label{eq:leaky}
\end{equation}
Коэффициент $(1-1/\tau)$ — доля прошлого состояния, которая \emph{удерживается};
$1/\tau$ — доля нового входа, которая \emph{впитывается}. Постоянная времени $\tau$
задаёт «медленность»: большое $\tau$ — слой почти не меняется за шаг (удерживает
долго, впитывает по чуть-чуть), малое $\tau$ — резко реагирует и быстро забывает.
\end{definition}

\begin{intuition}[Ведро с дыркой]
Протекающий интегратор — это ведро с дыркой в дне, куда капает вода (вход). Уровень
воды $s$ растёт от капель и падает от протечки. Широкая дырка (малое $\tau$) —
уровень скачет за каждой каплей и быстро сходит; узкая дырка (большое $\tau$) —
уровень меняется медленно, зато держится долго. «Сессия» — ведро с широкой дыркой,
«жизнь» — с почти закрытой.
\end{intuition}

\section{Два разных смысла «слоёв» — не путать}
\label{sec:twolayers}

Здесь важная развилка, которую легко замылить: слово «слои» и «протекание» означают
\emph{разное} для настроения и для памяти.

\paragraph{Настроение — параллельные интеграторы, переноса нет.}
Слои настроения (глава~\ref{ch:mood}) — это несколько интеграторов
\eqref{eq:leaky} с разными $\tau_\ell$, работающих \emph{параллельно на одном и том
же потоке событий}. Старший слой не питается от младшего; оба питаются событиями,
просто старший интегрирует медленнее. «Младший дрожит сильнее» — это высокочастотная
компонента того же сигнала. \textbf{Переноса содержания между слоями нет} — есть
разночастотная фильтрация одного входа, а актуальное настроение — суперпозиция.

\paragraph{Память — перенос есть, это повышение слоя консолидацией.}
Слои памяти устроены иначе: узел \emph{повышается} из быстрого слоя в медленный, когда
консолидация делает его устойчивым (раздел~\ref{sec:consolidation}). Вот это и есть
настоящее «протекание младшего в старший»: не диффузия сигнала, а
\emph{перемаркировка} узла (\texttt{layer}: session $\to$ day $\to \dots$) при
накоплении силы. Механизм известен по форме; конкретные пороги повышения — открыты.

\begin{openq}[Пороги повышения слоя]
При какой накопленной силе / частоте обращений / значимости узел повышается в
следующий слой? Разумные кандидаты: порог по базовой активации ACT-R
(раздел~\ref{sec:decay}); порог по числу независимых подтверждений; порог по
$|\conf|$. Точные значения и то, единый ли это порог для всех типов узлов, — пока не
зафиксированы; это параметр для стенда (глава~\ref{ch:experiment}).
\end{openq}

\section{Консолидация как огрубление графа}
\label{sec:consolidation}

Переход от эпизодов к смыслам — это \emph{огрубление} (coarsening): десять записей
«зритель заходил в понедельник», «\ldots в среду», \ldots заменяются одним узлом
«зритель — частый гость». Мелкая структура схлопывается в крупную; часть информации
теряется намеренно.

\begin{principle}[«Консолидация» — это процесс, а не один оператор]
Здесь важно снять терминологическую путаницу, из-за которой легко запутаться (и мы
путались). Под словом \textbf{консолидация} мы понимаем не единичный оператор
«свернуть $A,B,C$ в $X$», а \emph{процесс комплексной отработки нескольких разных
операторов}, каждый со своими правами и рисками:
\begin{enumerate}\itemsep2pt
\item \textbf{Fusion\,/\,вывод} — рождает новый смысл («поездка на Азоры», «ходит по
нечётным»). Посылки \emph{не} трогает (гл.~\ref{ch:fusion}).
\item \textbf{Summary-overlay (мягкая консолидация)} — заводит узел-обобщение
\emph{поверх} исходных, дети остаются, добавляются рёбра членства. Безопасный дефолт.
\item \textbf{Dedup\,/\,merge} — склеивает \emph{реально один и тот же} факт из разных
источников (это слияние identity, а не абстракция).
\item \textbf{Prune\,/\,сжатие с потерей} — вот единственный оператор, который
\emph{удаляет} детали. Редкий, поздний, жёстко гейтируемый.
\end{enumerate}
Почти вся ранняя консолидация — это (1)+(2): создать смысл и структуру доступа, ничего
не разрушая. Опасен только (4), и то лишь при преждевременном применении. Когда ниже
написано «консолидация сносит избыточное» — речь про оператор (4) внутри процесса, а не
про процесс целиком.
\end{principle}

\begin{intuition}[Почему нужен активный синтез, а не один только decay]
Легко спросить: если узлы и так тускнеют по непосещаемости, зачем вообще что-то
сворачивать? Затем, что \emph{decay} (смерть по непосещаемости) — \emph{слепое
выбрасывание}: он теряет данные, \emph{не извлекая из них закон}. Если тридцать
эпизодов «приходил по нечётным» просто протухнут, ты потеряешь их — и так и не узнаешь,
что он ходит по нечётным. А связка \emph{синтез (fusion) $+$ поздний прунинг} — то, что
мы и зовём консолидацией как процессом, — делает \emph{синтез, потом выброс}: сначала
дистиллирует gist, и лишь потом ретирует ставшую лишней деталь. Формула: \emph{выучить
урок (родить обобщение) и выбросить домашку (снести избыточное)}. Аналогия: ты не
помнишь каждый восход, но твёрдо знаешь «солнце встаёт на востоке» — инвариант извлечён,
экземпляры отброшены; будь у тебя только decay, ты бы просто забывал восходы, а правила
бы не получил. И ещё: часто-посещаемый избыточный узел по decay \emph{никогда} не умрёт
(его же трогают) и будет вечно раздувать граф — только активный синтез его свернёт.
\end{intuition}

Как понять, что огрубление \emph{хорошее}? Здесь мост к математике DTI. Введём:
\begin{itemize}\itemsep1pt
\item $L_f$ — лапласиан \emph{мелкого} (fine) графа; лапласиан описывает, как сигнал
диффундирует по графу (кто на кого влияет и насколько);
\item $L_c$ — лапласиан \emph{крупного} (coarse) графа после огрубления;
\item $P$ — оператор \emph{проекции}: сопоставляет каждому мелкому узлу его крупную
группу (например, десять эпизодов $\mapsto$ узел «частый гость»);
\item $\Phi(X,L)$ — оператор распространения (одна или несколько итераций передачи
сообщений) сигнала $X$ по графу с лапласианом $L$.
\end{itemize}
Хорошее огрубление \emph{коммутирует} с распространением:
\begin{equation}
\big\|\,\underbrace{P\,\Phi_f(X,L_f)}_{\text{распространить, потом огрубить}}
\;-\;\underbrace{\Phi_c(PX,L_c)}_{\text{огрубить, потом распространить}}\,\big\|
\;\le\;\varepsilon .
\label{eq:commute}
\end{equation}
Читается так: неважно, в каком порядке действовать — сначала «подумать» на мелком
графе и затем сжать, или сначала сжать и «подумать» на крупном, — результат должен
почти совпасть. Норма $\|\cdot\|$ берёт разницу по всем входным сигналам $X$;
$\varepsilon$ — мера потери. Малое $\varepsilon$ = огрубление сохранило то, что нужно
для рассуждения; большое = сжатие сломало динамику.

\begin{workedbox}[Что должно выжить при консолидации]
Свернём десять эпизодов визитов в узел «частый гость». \emph{Выжить} обязаны:
частота (была высокой), устойчивая валентность (тёплая), тип связи (регулярность).
\emph{Потеряться} можно: точные даты, дословные реплики. Плохое огрубление сохранит
даты и потеряет валентность — тогда Рин будет помнить \emph{когда} заходил зритель,
но забудет, \emph{как относится} к нему; это большая $\varepsilon$ по аффективной
оси. Значит критерий огрубления — сохранять аффективно- и реляционно-релевантное
подпространство, а не спектр целиком.
\end{workedbox}

\subsection{Коллизии предложений: голоса, а не очередь}
\label{sec:consolidation-collisions}

Огрубление предлагают обходчики (гл.~\ref{ch:walker}, раздел~\ref{sec:proposecommit}),
и предложения будут \emph{перекрываться}. Пусть за ночь на узлах $A,B,C,D,E$ накопилось:
RW1 хочет слить $\{A,B,C,D\}\to X_1$; RW2 — $\{B,C,D,E\}\to X_2$; RW3 — $\{A,B\}\to Y_1$;
RW4 — $\{C,D,E\}\to Y_2$. Наивный коммиттер «кто первый — тот и слил» построит мусор,
зависящий от случайного порядка. Правильный ход — понять, что \emph{отдельное
предложение обходчика это не решение, а свидетельство}: «этих тянет вместе». Тогда
коллизия — не спор за очередь, а \emph{сигнал для агрегации}.

Механически коммиттер копит предложения в окне (обычно за ночной проход) и строит
\emph{взвешенный со-консолидационный гиперграф} на исходных узлах: вес $c(i,j)$ = как
часто обходчики клали $i$ и $j$ в один предлагаемый кластер, взвешенно по support и —
важно — \emph{дисконтированно по времени}: голоса в буфере экспоненциально стареют,
иначе кластеризация будет учитывать давно неактуальные связи. По этому графу выделяются
\emph{кандидатные кластеры}.

И вот тут — принципиальная поправка: не надо загонять всё в \emph{единый жёсткий
раздел} (partition). Память — не разбиение; один эпизод «смешной вечер в универе с Анной,
когда я устал» законно входит и в «про универ», и в «про Анну», и в «про усталость», и в
«смешное». Раз parent у нас — обратимая надстройка-summary, а не разрушение мелких узлов,
\textbf{перекрывающиеся кластеры разрешены}. Поэтому: выделяем кандидатные кластеры
(модулярность\,/\,кластеризация — как один из инструментов, не как истина в последней
инстанции), разрешаем перекрытие, а коммиттер выбирает \emph{неконфликтующие} коммиты или
заводит \emph{несколько} parent-summary, если они сжимают \emph{разные оси}. В примере
$\{B,C,D\}$ — сильное ядро (RW1$+$RW2), $\{A,B\}$ и $\{C,D,E\}$ могут ужиться как разные
полезные обобщения; а вот если $X_1$ и $X_2$ оказываются почти дубликатами по одной оси —
их сливаем. Решает не порядок поступления, а согласие свидетельств и польза.

Четыре оговорки, без которых механизм наивен. \emph{Дисконт независимости}: два
обходчика, прошедшие по общим путям, — это часто один сигнал, засчитанный дважды; вес
$c(i,j)$ дисконтируется по перекрытию провенанса предложений (тот же Bloom\,/\,MinHash,
что и в пропагации, гл.~\ref{ch:propagation}). \emph{Гистерезис}: сливаем только то, что
рекуррентно всплывает несколько ночей; одноразовое предложение — шум, ему рано менять
структуру. \emph{Критерий пользы, а не только повторяемости}: новый parent
оправдан, если он \emph{сжимает} — уменьшает будущую стоимость ретрива\,/\,packer'а или
повышает качество реакции; частая со-встречаемость без выигрыша в сжатии — не повод для
обобщения. \emph{Обратимость через стадийность}: снос избыточных детей идёт не сразу
топором, а стадийно — сперва \texttt{pending}\,/\,tombstone (см.\ ниже) с сохранением
support-skeleton (гл.~\ref{ch:runtime}), и только устойчиво подтверждённое обобщение
доводит прунинг до конца. Поэтому цена даже неудачного слияния низкая — его можно
откатить, и это снимает надобность в идеальном арбитраже с первого раза.

\begin{intuition}[Комитет, а не турникет]
Представьте не турникет, где предложения проходят по одному и первое побеждает, а
\emph{комитет}, который выслушивает все голоса разом и выносит одно связное решение.
Обходчики — свидетели с разной надёжностью; коммиттер — тот, кто взвешивает показания
(со скидкой на сговор — общий провенанс) и решает, какие обобщения завести — при этом
узел может состоять сразу в нескольких, если они про разное.
\end{intuition}

\subsection{Что сносим, что оставляем — и кто это решает}
\label{sec:consolidation-prune}

Итак, консолидация в итоге \emph{выводит избыточные узлы из активного графа} — иначе
он только растёт. Уточним слово: не «убивает» (физическое стирание — это уже не память,
а бухгалтерия), а переводит узел в одно из состояний — \emph{active} $\to$
\emph{archived}\,/\,\emph{support-skeleton}\,/\,\emph{tombstone} (event-log хранит
исходное всегда, гл.~\ref{ch:runtime}). Полезно увидеть, что в процессе слиты
\emph{две} разные операции, и решение о второй принимается тоньше, чем «дети
консолидации умирают»:
\begin{itemize}\itemsep2pt
\item \textbf{(а) рождение обобщения} — создаётся узел-gist («приходит по нечётным»,
«частый гость»). Это \emph{создание нового}, по духу близко к fusion.
\item \textbf{(б) прунинг избыточных детей} — сносится то, что новый gist уже
полностью покрывает.
\end{itemize}
\textbf{Кто решает — коммиттер, а не обходчик}, и решает он про (б) \emph{поэкземплярно,
по избыточности}, а не «по кластеру, которому узел якобы принадлежит». Обходчик лишь
свидетель: «здесь есть сжимаемая конфигурация». Формально перед прунингом ребёнка $v$
коммиттер прогоняет \emph{residue-тест} — «есть ли у $v$ уникальный остаток?»:
\begin{enumerate}\itemsep1pt
\item уникальные внешние рёбра (мост к другим частям графа, которые осиротеют)?
\item уникальный timestamp или actionable-деталь (дата, билет, обязательство)?
\item высокая значимость\,/\,валентность?
\item связь с \emph{другим} обобщением?
\item естественный запрос, на который $v$ отвечает лучше parent'а?
\end{enumerate}
Если хоть один пункт да ($\text{residue}\ne 0$) — \emph{не прунить}, максимум архивировать
или оставить support-skeleton. Только при чистом нуле $v$ выводится из активного графа;
и даже тогда он может просто потухнуть по ACT-R, если к нему перестанут обращаться. В чистом случае «по нечётным» под
критерий сноса попадают \emph{все тридцать} экземпляров (они взаимозаменяемы, остатка и
мостов нет) — остаётся один узел; консолидация там максимально беспощадна, и правильно.
Мостовой богатый узел (эпизод, на котором висят два разных обобщения) — исключение не
потому, что «мы не прунуем», а потому, что он \emph{не избыточен}.

\paragraph{Право на схлопывание: экземпляр против составляющей.} Заострим, \emph{что}
именно делает ребёнка избыточным, — это и есть та тонкая грань, которую легко
проскочить. Схлопывать (с потерей) можно только \emph{класс эквивалентных экземпляров}
одного паттерна: тридцать «заходил по нечётным», сто «привет», каждый восход. Они
взаимозаменяемы (`instance\_of`), и gist их полностью замещает. А \emph{разнородный
букет составляющих} одной темы (`part\_of`) схлопывать нельзя \emph{в принципе} — там
нечего терять, там всё уникально. Отсюда типичная ошибка: «fusion сгруппировал
$A,B,C$» \emph{не} значит «$A,B,C$ можно схлопнуть». Группировка-по-теме $\ne$
поглощение-избыточности. Разберём на неигрушечном примере.

\begin{workedbox}[Поездка на Азоры: fusion есть, схлопывать нечего]
Пять фактов: (A)~купил плавки, (B)~купил зонтик, (V)~билет на Азоры на 5 сент.,
(G)~автобус аэропорт Лиссабона\,$\to$\,Порто на 9 сент., (D)~конференция в Порто
9--12 сент. Напрашиваются два fusion: $X_1$ = «отпуск на Азорах» ($A,B,V$),
$X_2$ = «маршрут Азоры\,$\to$\,Лиссабон\,$\to$\,конференция» ($V,G,D$). Но
\textbf{схлопывать здесь нечего}: $A,B,V,G,D$ — не экземпляры паттерна, а разнородные
составляющие. Прунинг убил бы ровно то, что нужно для действия: даты, билеты,
предметы, обязательства. Поэтому — только (1)+(2): рождаем $X_1,X_2$ как overlay,
факты оставляем, вешаем ролевые рёбра (\texttt{equipment\_for}, \texttt{leg\_of},
\texttt{purpose\_of}).

Тонкость \emph{event vs state}: «купил плавки» (event) порождает отдельный узел «вероятно,
владеет плавками» (state, другой \texttt{node\_kind}, гл.~\ref{ch:substrate}); event со
временем уйдёт в архив, а state живёт, пока не появится «потерял\,/\,выкинул\,/\,не
подошли». Потому и на «Рин, поищи плавки» через два месяца она поднимает state как
\emph{контекст} и в худшем случае уточняет («вроде брал в августе — ищем на замену?»),
а не блокирует запрос старым фактом. Это отдельный инвариант — \emph{память — контекст,
а не команда}: прошлый факт не отменяет текущее намерение, а лишь оттеняет его.
\end{workedbox}

\begin{principle}[Два инварианта консолидации]
\emph{Первый: вывод обходчика — свидетельство, а не власть.} Коммиттер решает задачу
\emph{отбора модели} — какой набор новых узлов лучше объясняет данные при минимальной
сложности (MDL-настроение: выигрыш в сжатии $+$ связность $+$ повторяемость $+$ польза
$-$ избыточность $-$ сложность $-$ эпистемический риск), а не исполняет предложения
слепо. \emph{Второй: перекрывать обобщения — можно, дублировать identity — нельзя.}
Один исходный узел может быть посылкой нескольких parent-summary через рёбра членства
(\texttt{member\_of}\,/\,\texttt{provenance}); но \emph{клонировать} его на несколько
копий в графе запрещено.
\end{principle}

\begin{pitfall}[Почему нельзя клонировать узлы («графовый мейоз»)]
Соблазн разрулить пересечения так: развернуть общий узел $C$ в близнецов $C_1,C_2$,
раздать по разным обобщениям, оригинал снести — и предложения обходчиков станут чистым
разбиением. Формально приём валиден (graph unrolling), но семантически ломает память в
двух местах. \emph{Расщепление identity убивает TMS}: опровергнут позже факт $C$ —
ретракция поправит один клон, второй останется фантомом, лавинная ревизия уверенности
сломается. \emph{Разрыв моста}: общий $C$ был перешейком, по которому активация текла
между $X_1$ и $X_2$ («раздражается» $\leftrightarrow$ «не выспалась» связывались
\emph{через} $C$); разнёс копии — потерял вывод навсегда. И «раздваиваются только на
время» — иллюзия: дочерние не сливаются обратно, двойной след остаётся \emph{навсегда}.
Спасение идеи: близнецы допустимы \emph{только как черновая вычислительная техника} в
блокноте коммиттера (чтобы независимо оценить кандидатов), но в \emph{коммите} общий
узел остаётся общим — identity в графе не двоится.
\end{pitfall}

\subsection{Кандидаты, наследование линков, слияние уверенностей}
\label{sec:consolidation-mechanics}

Три механических вопроса, которые всплывают на реализации.

\paragraph{Жаккард\,/\,LSH — сигнал, а не критерий.} Соблазн: «полное перекрытие
окрестностей $\Rightarrow$ поглощаем». Опасно. Структурная эквивалентность $\ne$
семантическая замещаемость: «купил плавки» и «купил зонтик» имеют почти одинаковых
соседей (Азоры, пляж, лето, покупка), но не взаимозаменяемы — один отвечает «есть ли
плавки», другой «есть ли зонтик». Поэтому индекс Жаккара (и его дешёвая оценка через
MinHash\,/\,LSH за $O(|V|)$ вместо $O(|V|^2)$ — да, это правильный способ искать
кандидатов ночью) годится как \emph{кандидатный сигнал} «посмотри сюда, возможно дубль
или класс экземпляров», но \emph{никогда} как критерий удаления. За ним обязателен
семантический гейт «это экземпляр без уникального остатка?». Даже «новое нечётное число
$\approx$ паттерн „нечётные“» законно поглощается (инкремент support-count) лишь если
этот экземпляр пустой; а если 9-е — день, когда он не успел на автобус, экземпляр не
пуст, не поглощай.

\paragraph{Два разных действия с рёбрами — не путать в коде.} Сначала разведём то, что
легко слить в одну функцию и устроить себе Чернобыль. \emph{Overlay role-edge} — новая
ролевая связь \emph{parent-summary} (`$X_1$ --equipment\_for--> плавки`), создаётся уже
на стадии overlay, \emph{без всякого прунинга} $A$; это про \emph{создание} структуры
доступа. \emph{Prune edge-lifting} — подъём \emph{всех} рёбер ребёнка на parent при его
сносе; работает \emph{только} для чистых экземпляров. Это разные операции; общей функции
\texttt{lift\_edges(child, parent)} на все случаи быть не должно.

\paragraph{Наследование линков при прунинге: та же онтология — но только где прунинг законен.}
Простое правило «рёбра ребёнка переходят на parent с теми же типами отношений» —
\emph{верно ровно для чистых экземпляров}, а это единственный случай, где наследование
вообще происходит. Для них это буквально проекция смежности $A_c=P\,A_f\,P^\top$: дубли
рёбер сливаются, веса складываются и \emph{ренормируются} (L1, чтобы переходы обходчика
остались вероятностями), провенанс-сигнатуры объединяются побитовым OR. А там, где то
же правило даёт бессмыслицу — `поездка --object--> плавки` — это \emph{сигнал}, что ты
прунишь не экземпляр, а составляющую, которую прунить нельзя. Отсюда чистый
предохранитель: \textbf{если ребро нельзя поднять, не меняя его смысл, — узел не
прунится} (максимум понижается\,/\,архивируется). Невозможность безопасно унаследовать
линк сама останавливает неуместный прунинг.

\paragraph{Слияние уверенностей: среднее в confidence, разброс — в один дешёвый скаляр.}
Когда сливаются рёбра с разными знаками\,/\,уверенностями, среднее (со знаком) идёт в
\texttt{confidence} нового ребра. Но сворачивать разногласие \emph{целиком} в один
скаляр нельзя: теряется различие «0,5, потому что честно пятьдесят на пятьдесят и
устойчиво» против «0,5, потому что половина источников за, половина против» (contested)
— а это разные эпистемические состояния (гл.~\ref{ch:substrate}), и от них зависит,
говорит Рин спокойно или помечает как спорное. Хранить все голоса — расточительно
(«страниц не хватит»). Компромисс: держать \emph{одно дешёвое число} — дисперсию\,/\,флаг
\texttt{contested} (переиспользуя Beta-дисперсию реакций, гл.~\ref{ch:personality}, или
механизм \texttt{contested\_core\_claim}, гл.~\ref{ch:sleep}). $O(1)$ на ребро, а не
история голосований.

\section{Забывание как тень консолидации}
\label{sec:decay}

Забывание — не отдельный механизм, а \emph{отрицательная сторона} того же процесса:
усиливай используемое, ослабляй неиспользуемое.

\subsection{Правильный decay — это ACT-R}

\emph{ACT-R} (Adaptive Control of Thought—Rational) — влиятельная когнитивная
архитектура Джона Андерсона, модель того, как устроены человеческие память и
научение. Нам из неё нужен один компонент — \emph{базовая активация} (base-level
activation): сила следа в памяти есть функция \emph{частоты} и \emph{свежести}
обращений со степенным затуханием:
\begin{equation}
B_v \;=\; \ln\!\Big(\sum_{k=1}^{n} t_k^{-\alpha}\Big).
\label{eq:actr}
\end{equation}
Разбор: $t_k$ — сколько времени назад было $k$-е обращение к узлу; каждое обращение
вносит вклад $t_k^{-\alpha}$, и чем оно \emph{свежее} (меньше $t_k$), тем вклад
больше; сумма по всем $n$ обращениям даёт накопленную силу (\emph{частота});
логарифм сжимает шкалу; $\alpha>0$ — скорость забывания. Два следствия, которых нет
у плоской «$\times 0{,}95$ в день»: частые обращения укрепляют след, а
\emph{разнесённые во времени} укрепляют сильнее массовых (эффект распределённого
повторения — spacing).

\begin{intuition}[Тропинка в лесу]
Тропинку протаптывают многими проходами; если по ней ходят изредка, но регулярно,
она держится лучше, чем если пройти сто раз за день, а потом забросить. Формула
\eqref{eq:actr} именно это и кодирует: важна не только суммарная частота, но и её
распределение во времени.
\end{intuition}

\begin{pitfall}[Доступ сам тренирует память — различай access\_reason]
Тонкий баг: не всякое обращение к узлу должно укреплять его. Если ночной скраб,
debug-инспектор или retrieval-тест часто трогают узел, он не должен от этого
становиться «важным жизненным следом» — иначе \emph{сама диагностика тренирует
память}. Поэтому у каждого доступа есть \texttt{access\_reason}: \texttt{user\_recall},
\texttt{model\_recall}, \texttt{maintenance\_touch}, \texttt{debug\_read},
\texttt{scrub\_visit}, \texttt{context\_packed}. В ACT-R-активацию \eqref{eq:actr}
входят с полным весом только \emph{подлинные} припоминания (user/model recall);
обслуживающие тачи — с нулевым или почти нулевым. Наблюдатель не должен менять
наблюдаемое.
\end{pitfall}

\subsection{Прунинг: не только по использованию}

Прунинг только по использованию опасен: редкий, но критичный факт вымоется как
мусор. Поэтому склонность к удалению гейтируется \emph{важностью} узла:
\begin{equation}
\text{prune}(v)\;\propto\;\text{неиспользование}(v)\,\cdot\,\big(1-\text{importance}(v)\big).
\end{equation}
Здесь \texttt{importance} — атрибут узла (см.\ схему в главе~\ref{ch:substrate}),
задаваемый значимостью при кодировании (salience из appraisal,
глава~\ref{ch:personality}) и подкрепляемый структурной центральностью
(глава~\ref{ch:topology}). Редкое-но-важное выживает; частое-но-пустое держит
частота; вымывается лишь редкое-и-неважное. Врождённым (глубоко консолидированным)
узлам задаётся ниже скорость распада — континуум глубины, а не бинарность.

\subsection{importance нельзя диффундировать по путям}
\label{sec:importance-leak}

Здесь важная ловушка. Соблазн: раз узел-домен «цепи Маркова» важен, пусть раздаёт
$\text{importance}\to 1$ своим детям, те — своим, и так область защищена. \textbf{Так
делать нельзя.} importance — это диффузия по графу (глава~\ref{ch:topology}), а
диффузия растекается \emph{везде}: через несколько хопов по каким-нибудь рёбрам она
доползёт до «помыть машину» (которой у тебя нет), и та станет «важной». Графовые пути
неисповедимы; неуправляемая передача importance приводит к тому, что важным
становится \emph{всё}, а значит — ничто.

Правильно — \textbf{не пускать importance по произвольным путям, а закреплять её
членством в домене}:

\begin{itemize}\itemsep2pt
\item \textbf{importance по членству, а не по диффузии.} У узла есть \emph{теги
доменов}, у доменов — importance. Эффективная важность узла $=$ функция от
importance его доменов $+$ собственная salience. Членство назначается \emph{при
кодировании} (прочитали учебник $\to$ его узлы помечены \texttt{domain=Markov},
высокая importance), а не распространяется обходом. «Помыть машину» никогда не
получает тег \texttt{Markov}, поэтому \emph{ничего не наследует}. Утечка исключена по
построению.
\item \textbf{importance по намерению.} То, что \emph{попросили} выучить, получает
высокую importance по провенансу приобретения (в отличие от случайного тривиа).
\item \textbf{Если всё же нужна передача} (для непомеченных узлов) — только
\emph{типо-гейтированная} (по рёбрам членства \textsc{isA}/\textsc{partOf}, но не по
ассоциативным), \emph{сильно затухающая} ($\lambda^k$ с малым $\lambda$) и
\emph{ограниченная} 1--2 хопами; и это надо \emph{проверять} на стенде, что importance
осталась локальной.
\end{itemize}

Тогда защита области (глава~\ref{ch:runtime}, раздел~\ref{sec:storage}) работает
аккуратно: домен-хаб с высокой importance и низким распадом $+$ члены, наследующие
\emph{по членству}; детали могут тускнеть, но gist области держится, а сон-репетиция
периодически освежает высоко-importance домены — без утечки на «помыть машину».

\section{Что процесс делает с найденной конфигурацией}
\label{sec:whatfusiondoes}

Предположим, фоновый разведчик (глава~\ref{sec:background}) со-активировал тройку
$\{A,B,C\}$, которая выглядит связной. Что происходит дальше? Разделю на
\emph{известное} и \emph{открытое}, чтобы не размывать.

\paragraph{Что известно (устоявшийся конвейер).}
\begin{enumerate}\itemsep1pt
\item \textbf{Композиция.} Языковая модель по объединению узлов и их контекстных
ореолов (глава~\ref{ch:fusion}) формулирует \emph{кандидата} $C^\ast$ — обобщение
или вывод.
\item \textbf{Провизорная запись.} $C^\ast$ пишется с низкой уверенностью и
провенансом $\{A,B,C\}$; на различающих пробелах ставится \texttt{need\_to\_discover}.
\item \textbf{Конкуренция.} Дальше $C^\ast$ живёт по общим правилам: используется и
корроборируется $\Rightarrow$ сила растёт, узел повышается в слой
(раздел~\ref{sec:twolayers}); игнорируется $\Rightarrow$ распадается
(раздел~\ref{sec:decay}); опровергается $\Rightarrow$ понижается по TMS
(глава~\ref{ch:epistemics}).
\end{enumerate}

\paragraph{Что открыто.} Как \emph{решить}, что со-активация $\{A,B,C\}$ достаточно
связна, чтобы вообще пытаться её сшить (а не является случайной кучей), — не решено.

\begin{openq}[Распознавание связности со-активации]
Кандидаты: \emph{(а)} LLM-судья («есть ли тут осмысленный вывод?») — гибко, но не
доказуемо и стоит вызова модели; \emph{(б)} графовые шаблоны-схемы («потребность $+$
ограничение $+$ окно $\to$ перенести действие») — дёшево, но покрывают лишь частые
формы; \emph{(в)} обучаемый скорер связности — требует разметки и вносит
непрозрачность. Ни один пока не признан удачным; на старте — (а) с фильтром $\varphi$
и обратной связью (главы~\ref{ch:personality},~\ref{ch:voice}). Это и есть «трудные
80\%» из главы~\ref{ch:fusion}.
\end{openq}

\subsection{Забывать ли посылки после fusion?}

Возникает соблазн: раз $\{A,B\}\to C$ состоялось, посылки $A,B$ больше не нужны —
удалить. Это \textbf{неверно} (слишком мануально и опасно). Правильно — не удалять
насильно, а положиться на естественный распад:

\begin{itemize}\itemsep1pt
\item при \emph{хорошей} консолидации к $A,B$ просто перестают обращаться — вся
активность идёт к более полезному $C$; базовая активация \eqref{eq:actr} $A,B$
падает, и они отмирают \emph{сами} (важные — выживают). Забывание тут
\emph{эмерджентно}, а не по правилу «удалить родителей»;
\item удалять $A,B$ немедленно нельзя ещё и потому, что они — \emph{провенанс} $C$:
если $C$ позже опровергнут, по $A,B$ его пересматривают или переводят
(глава~\ref{ch:epistemics}). Убрав посылки, теряешь возможность ревизии.
\end{itemize}
Итог по-человечески верный: детали тускнеют, суть остаётся, но след обоснования
хранится, пока не осыпался по неиспользованию.

\begin{principle}[support skeleton: обоснование переживает распад деталей]
Здесь есть скрытый конфликт: провенанс требует держать посылки для будущей ревизии
(TMS), но детали естественно \emph{decay}-ятся. Разрешение — не хранить полный
богатый эпизод вечно, а при распаде оставлять \emph{support skeleton}: минимальный
стаб \{id, source, time, hash, node\_kind, claim-summary\}, достаточный для ретракции.
Тогда либо-либо снимается: подробность уходит, а \emph{способность ревизовать}
остаётся. Без скелета либо память не забывает, либо потом нечем пересматривать
выводы.
\end{principle}

\chapter{Настроение}
\label{ch:mood}

\section{Настроение — отдельный сцепленный модуль, а не поле графа}

Снимем путаницу сразу. Настроение \emph{не} вычитывается из спектра графа памяти.
Это отдельное состояние — низкоразмерный вектор в пространстве аффекта
$\mood\in\mathbb{R}^k$ (например, оси
энергичность\,/\,раздражённость\,/\,экзальтация\,/\,любопытство) — отдельный модуль,
слой, через который проходит извлечённое содержание.

Путаница шла от слова «спектр». Есть два разных спектра, и их нельзя смешивать:
\emph{спектр лапласиана графа} (структура памяти, глава~\ref{ch:topology}) и
\emph{частотное разложение настроения во времени} (быстрые и медленные компоненты
аффективного сигнала). Многослойность настроения относится ко второму и к графу
отношения не имеет.

\section{Инерция и многослойность}

У настроения должна быть \emph{инерция}: оно зависит и от события, и от собственного
прошлого значения:
\begin{equation}
\mood_t \;=\; f\big(\mood_{t-1},\,\text{event}_t\big),\qquad
\text{action}_t \;=\; g\big(\mood_t,\,\text{event}_t\big).
\label{eq:moodinertia}
\end{equation}
Это протекающий интегратор \eqref{eq:leaky}: инерция — это его рекуррентность
(состояние зависит от прошлого состояния).

\begin{intuition}[Почему интегратор, а не выученная динамика]
Форму $f$ можно было бы \emph{выучить} — небольшой RNN/GRU поверх потока событий. Мы
берём фиксированный протекающий интегратор по трём причинам. \emph{Интерпретируемость}:
у него ровно одна ручка на слой — постоянная времени $\tau_\ell$ (как быстро забывает),
и её смысл читается сразу; у выученной динамики «почему настроение поехало» — это
скрытое состояние, которое не отладишь глазами. \emph{Отсутствие обучения и данных}:
RNN надо на чём-то тренировать, а размеченных траекторий настроения Рин у нас нет и не
предвидится. \emph{Достаточность}: главные наблюдаемые свойства — инерция, возврат к
базовому уровню и разные скорости на разных горизонтах — интегратор даёт напрямую,
одним стеком с разными $\tau_\ell$. Выученную динамику держим как \emph{апгрейд на
потом}: если стенд покажет, что нужны нелинейные эффекты (насыщение, гистерезис аффекта),
которые линейный интегратор не тянет.
\end{intuition}

\begin{intuition}[Интегратор и вектор — одно и то же, с двух сторон]
Чтобы снять путаницу: «протекающий интегратор» и «вектор настроения» — не разные
сущности. Интегратор — это \emph{правило обновления} \eqref{eq:leaky}, а его
\emph{состояние} — вектор в пространстве аффекта. Многослойность настроения — это
\emph{стек} таких интеграторов с разными $\tau_\ell$; «старший вектор» и «младший
вектор», о которых мы говорили, — это ровно \emph{состояния} медленного и быстрого
интеграторов. Актуальное настроение — их суперпозиция. И, как отмечено в
разделе~\ref{sec:twolayers}, эти слои \emph{параллельны} (все питаются одним потоком
событий, переноса между ними нет) — в отличие от слоёв памяти, где перенос есть.
\end{intuition}

Формально: стек интеграторов с разными $\tau_\ell$, актуальное настроение — их
взвешенная суперпозиция:
\begin{equation}
\mood^{(\ell)}_t=\Big(1-\tfrac{1}{\tau_\ell}\Big)\mood^{(\ell)}_{t-1}+\tfrac{1}{\tau_\ell}\,\text{appraise}(\text{event}_t),
\qquad
\mood_t=\sum_\ell w_\ell\,\mood^{(\ell)}_t .
\end{equation}
Младший слой (малое $\tau$) даёт высокочастотную компоненту большой амплитуды —
«дрожит сильнее». «Суперпозиция» здесь — про частоты во времени, а не про граф.

\section{Setpoint: медленный слой как аттрактор}

Медленный слой — не просто демпфер, а \emph{точка притяжения}. Добавляя к
\eqref{eq:leaky} возврат к базовому уровню $\mu$, получаем затухающую к $\mu$
динамику:
\[
\mood^{(\text{life})}_t=\mu+(1-\kappa)\big(\mood^{(\text{life})}_{t-1}-\mu\big)+\ldots,\quad 0<\kappa<1 .
\]
Здоровая динамика — отрицательная обратная связь, тянущая настроение домой. Патология
— когда сдвинулся сам $\mu$: депрессия в этой модели не длинная экскурсия быстрого
слоя, а \emph{уехавший аттрактор}.

\section{Настроение зависит от графа, и граф — от настроения}

Связь двусторонняя, и именно она делает настроение частью \emph{мышления}.

\begin{itemize}
\item \textbf{Граф $\to$ настроение.} Извлечённые узлы несут валентность; она и есть
$\text{event}_t$ в \eqref{eq:moodinertia}.
\item \textbf{Настроение $\to$ граф.} Настроение смещает вероятности переходов при
обходе: узлы с валентностью, близкой к текущему настроению, получают больше
активации:
\begin{equation}
\text{mood\_congruence}(B\mid \mood)\;=\;\exp\!\big(-\gamma\,\|\,\text{valence}(B)-\mood\,\|^2\big).
\label{eq:moodcong}
\end{equation}
\end{itemize}
Это \emph{конгруэнтное настроению припоминание}: грустный чаще вспоминает грустное.
Настроение остаётся собственным состоянием, сцепленным сбоку, а не полем графа.

\section{Опасность петли и три тормоза}

\begin{pitfall}[Руминация]
Конгруэнтное припоминание \eqref{eq:moodcong} плюс инерция \eqref{eq:moodinertia}
дают петлю с положительной обратной связью: грустная вспоминает грустное $\to$
грустнеет $\to$ вспоминает ещё более грустное. Устойчивость петли определяется её
петлевым усилением $g$: при $g\ge 1$ спираль расходится. Медленный слой уменьшает
эффективное $g$ — вот его стабилизирующая роль.
\end{pitfall}

Что удерживает живой ум? Три \emph{независимых} тормоза; патология — когда откажут
все три сразу:

\begin{enumerate}
\item \textbf{Setpoint/гомеостаз.} Медленный слой-аттрактор снижает петлевое
усиление ниже единицы.
\item \textbf{Метакогнитивное разотождествление.} Акт «почему я так часто злюсь?»
переводит из «быть-в-настроении» в «наблюдать-настроение» — второй порядок, и
переход рвёт петлю первого порядка.
\item \textbf{Экзогенное прерывание.} Мир вламывается ортогональным событием;
внимание конечно, внешний вход вытесняет петлю.
\end{enumerate}

Инженерное требование: (i)~член setpoint; (ii)~монитор «валентность залипла слишком
долго» — высокая автокорреляция $\mood_t$ — с правом запустить корректив
(переопределение, переключение внимания, вызов сна); (iii)~всегда-прерываемая
экзогенная значимость. Условие безопасности: три тормоза не могут быть отключены
одновременно. Псевдокод монитора — в приложении~\ref{ch:algorithms}.

\begin{workedbox}[Почему многослойность — это стабилизатор, а не украшение]
Пусть быстрый слой попал в негативную спираль с усилением $g_{\text{fast}}=1{,}1$
(расходится). Актуальное настроение — суперпозиция, и медленные слои вносят
устойчивую компоненту с $g\approx 0$. Если их вес $w_{\text{slow}}$ достаточен,
эффективное усиление суперпозиции $\bar g=w_{\text{fast}}g_{\text{fast}}+w_{\text{slow}}\cdot 0<1$
— спираль гасится за счёт «тяжёлого маховика» медленных слоёв. Так многомасштабность
оправдана \emph{функционально}: это демпфер руминации, а не только реализм.
\end{workedbox}

\section{Метакогниция}
\label{sec:metacognition}

Один из трёх тормозов руминации выше — «метакогнитивное разотождествление» — был
назван мимоходом, но это несправедливо: \textbf{метакогниция — самостоятельная фича
всего проекта}, а не побочный эффект. Ей место в отдельном разделе.

\begin{definition}[Метакогниция]
\emph{Метакогниция} — способность Рин наблюдать и рассуждать о \emph{собственных
внутренних состояниях}: о своём настроении, об уверенности в своих убеждениях, о
пробелах в памяти. Это операции \emph{второго порядка} — не «что я думаю о мире», а
«что я думаю о своём думании и своём состоянии».
\end{definition}

Метакогниция не само собой разумеется: базовая модель по умолчанию её не имеет.
Это проектируемая возможность, и у неё несколько конкретных ролей, разбросанных по
всей архитектуре, — здесь мы их собираем.

\begin{itemize}\itemsep2pt
\item \textbf{Тормоз руминации.} Наблюдение «почему я так часто злюсь?» переводит
из режима «быть-в-настроении» в «наблюдать-настроение» и рвёт петлю первого порядка
(раздел выше). Технически — это монитор автокорреляции $\mood_t$, который \emph{сам
замечает} залипание и запускает корректив.
\item \textbf{Эпистемическая честность.} «Похоже на $X$, но $Y$ я не проверила» —
это метакогнитивный отчёт об уверенности: Рин читает поле $\conf$ и зону незнания
($|\conf|<h$) и \emph{говорит} о них, вместо ложного «я уверена»
(глава~\ref{ch:epistemics}).
\item \textbf{Порождение любопытства.} Замечая собственный пробел или высокую ошибку
предсказания, метакогниция ставит \texttt{need\_to\_discover} — то есть \emph{сама
формулирует}, чего Рин не знает и что стоит проверить (главы~\ref{ch:epistemics},
\ref{ch:fusion}).
\item \textbf{Самоотчёт.} Способность сказать пользователю о своём состоянии
(«устала от одного и того же бага», «мне сейчас не до шуток») — это метакогниция,
считывающая настроение и мотивацию и выражающая их через Rinness
(глава~\ref{ch:voice}).
\end{itemize}

\subsection*{Субстрат метакогниции: самотелеметрия}

Метакогниции \emph{не нужен} отдельный граф. Её субстрат — \emph{лог
самонаблюдения}: кольцевой буфер временного ряда собственных состояний (вектор
настроения) и счётчиков команд (\texttt{need\_to\_*}, fusion, contradiction,
retrieval) за часы/дни/недели. Словарь внутренних команд и состояний \emph{конечен
и дискретен}, поэтому телеметрия компактна и структурирована, а статистика по ней
дёшева. Это не граф и не свободный текст. Но \emph{выводы} метакогниции возвращаются
в граф памяти как обычные узлы-убеждения-о-себе (с $\conf$, провенансом = телеметрия-
свидетельство) — и дальше живут по общим правилам.

\subsection*{Детекция: дешёвый монитор чисел, а не сверх-LLM}

Первый слой — \emph{детекция} — набор дешёвых статистических детекторов на буфере:
скользящие среднее/медиана, тренд (наклон регрессии), частота (счёт в окне),
аномалия (z-оценка от базовой линии), автокорреляция (залипание). Детерминированно,
\emph{без} LLM; срабатывают по порогу/аномалии. Пример: «blue: медиана $0{,}7$ >
среднего $0{,}5$ за 5 дней $\Rightarrow$ несколько всплесков перетянули среднее, база
грустная»; «\texttt{need\_to\_revise}: $+60\%$ к 30-дневной базе».

В «мозг» Рин уходит \emph{не} сырой ряд (много, и LLM ненадёжно считает медиану) и
\emph{не} готовый вывод, а \emph{компактная структурная сводка} — выход детектора.

\subsection*{Аномалия — это факт в графе, а не спецпайплайн}

И вот основной ход, который важнее всего: сводку детектора пишем \emph{в граф памяти
как обычный факт-о-себе} («за 5 дней \texttt{need\_to\_revise} $+60\%$»; «я в
последнее время грустная»). Уверенность высокая — это \emph{измерение}, а не
догадка. Дальше — никакого выделенного «метакогниция-LLM» и никакого спец-пайплайна
«аномалия $\to$ LLM $\to$ граф». Факт просто лежит в памяти и подхватывается
\emph{тем же общим циклом}, что и любая другая мысль (retrieval, фоновый fusion).
Метакогниция здесь — не модуль, а \emph{факты-о-себе, живущие в графе по общим
правилам}.

Более того, \textbf{сам факт уже работает — вообще без LLM}. Узел «я в последнее
время грустная» смещает самомодель, mood-congruent retrieval и самоотчёт («что-то я
загрустила») через обычные механизмы. Значительная часть пользы приходит просто от
\emph{записи} факта; объяснение «почему» — опциональный более глубокий слой, а не
обязательная стадия.

\subsection*{Объяснение опционально: fusion-про-себя}

Когда причина всё же нужна (гипотезы о том, \emph{почему} аномалия), работает
fusion-про-себя (глава~\ref{ch:fusion}) — но это \emph{тот же общий композитор}, а
не прокладка: он запускается обычной salience-обработкой, когда факт-аномалия
всплывает. Гипотезы берут посылки из телеметрии: «почему больше
\texttt{need\_to\_revise}?» $\to$ H1 «много новых рассуждений», H2 «кладу
непроверенные факты», H22 «читала книгу с ошибками». Каждая — провизорный узел,
\emph{низкая} уверенность (в отличие от факта-измерения — высокой), провенанс из
телеметрии.

Проверка — по \emph{различающему предсказанию} (научный метод): H1 $\Rightarrow$
должен вырасти и fusion-count; H2 $\Rightarrow$ у пересмотренных узлов короткий
интервал commit$\to$revision и пропущенный \texttt{need\_to\_discover}; H22
$\Rightarrow$ они делят один источник в провенансе. \texttt{need\_to\_discover}
кодирует, \emph{какой} запрос выполнить; curiosity/scrub его выполняет. Разрешённая
гипотеза оседает в граф; если ничего не разводит — остаётся честное «неизвестно
почему», а не выдуманная причина.

\subsection*{Eager или lazy — из salience, а не из расписания}

Раз аномалия — это факт, «сразу разбираться или потом» перестаёт быть отдельной
развилкой в коде: это \emph{следствие salience} самого факта. Срочная аномалия
(спираль настроения, руминация) несёт высокую значимость $\Rightarrow$ всплывает и
обрабатывается общим циклом \emph{немедленно} (корректив сейчас). Несрочная
(revise-rate вырос) просто лежит и объясняется \emph{лениво} — когда её достанет
retrieval или спросит пользователь («как ты в последнее время?»). Дешёвые детекторы
при этом крутятся в фоне всегда (они дёшевы; их интенсивность привязана к
$\rho$-бюджету, глава~\ref{ch:walker}), а дорогой общий LLM тратится лишь на то, что
salience подняла наверх. Ни чистый on-demand (руминация подавляет самозапрос именно
тогда, когда он нужнее), ни фиксированное расписание (расточительно) не нужны.

\subsection*{Защита от конфабуляции}

Самообъяснение — самый склонный к выдумыванию вид рассуждения: люди уверенно
конфабулируют причины своих настроений и поступков (классика — Nisbett \& Wilson,
1977). Поэтому телеметрия-лог здесь не удобство, а \emph{противоядие}: он делает
самогипотезы проверяемыми против реальных данных, а не сочинёнными. Метакогниция без
телеметрии $=$ конфабуляция; с телеметрией и различающими проверками $=$
фальсифицируемая.

\begin{pitfall}[Само-факты не должны стать самосбывающейся ловушкой]
Идея «аномалия — факт в графе» сильна, но у неё есть тёмная сторона. Узел «я в
последнее время грустная» (\texttt{self\_model}/\texttt{mood\_summary}), показанный в
context-packer как обычный identity-факт, запускает петлю: mood-congruent recall
начнёт \emph{сам поддерживать} это состояние. Поэтому такие узлы: (1)~диагностичны и
\emph{time-bounded} (с явным expiry, а не «навсегда»); (2)~пакуются в промпт как
\emph{наблюдение о недавней телеметрии}, а не как черта: «recent telemetry suggests
\ldots», а не «Rin is \ldots». Разница между «в последнее время грустится» и «я
грустная» — это разница между диагнозом и приговором.
\end{pitfall}

\subsection*{Метакогниция о метакогниции (второй порядок)}

Тот же монитор можно навести на телеметрию \emph{самой} метакогниции: частоту
порождаемых самогипотез, долю разрешённых, долю опровергнутых. Это \emph{санитайзер}
первого порядка: ловит analysis paralysis (много гипотез, ни одна не разрешается
$\to$ перестать копаться), ненадёжную интроспекцию (самовыводы то и дело
опровергаются $\to$ доверять самочтениям меньше — прямая защита от конфабуляции) и
залипшие петли. Механика та же, просто этажом выше, поэтому дёшево. Третий порядок
(«думаю о том, как думаю о думании») даёт быстро убывающую отдачу и риск регресса —
потолок разумно ставить на втором (аналогия со второй производной в матанализе).

\subsection*{Как это соотносится с литературой}

Идея «метакогниция $=$ System-2, наблюдающая System-1» и то, что мета-мысли
\emph{возвращаются в память} и затем влияют на действия, — из
Toy, MacAdam, Tabor (2024), \emph{«Metacognition is all you need?»} (поверх
генеративных агентов Park и др.). Что мы \emph{берём}: самонаблюдение как отдельный
модуль; мета-выводы обратно в память; локальные модели; дешёвая модель для рутины,
большая — для метакогниции. Где идём \emph{дальше} (и где вопросительный знак в
названии честно себя выдаёт): у них \emph{детекция — это LLM}, оценивающий свой
прогресс, а у нас — дешёвый статистический монитор, LLM только на объяснение; у них
нет \emph{проверяемых} самогипотез с различающими запросами и защиты от
конфабуляции; их метакогниция целевая (задача), наша шире (настроение, эпистемика,
гигиена памяти) — то, что они сами относят к future work. Значит «метакогниция — всё,
что нужно»? Нет: нужен ещё субстрат самотелеметрии и дисциплина проверяемых гипотез.

\begin{intuition}[Второй наблюдатель]
Полезно представлять метакогницию как \emph{второго наблюдателя}, который смотрит не
на мир, а на первого — на саму Рин: на её настроение, уверенность, ход мысли. Именно
этот второй наблюдатель отличает существо, которое \emph{переживает} состояние, от
существа, которое \emph{знает, что переживает его, и может это назвать}. Для проекта,
где цель — личность, а не ассистент, это не украшение, а одна из несущих опор.
\end{intuition}

\begin{openq}[Глубина и стоимость метакогниции]
Открыто, насколько «глубоко» гонять метакогницию: второй порядок («думаю о думании»)
дёшев и явно полезен, но третий и выше («думаю о том, как думаю о думании») быстро
даёт убывающую отдачу и риск зацикливания на самонаблюдении. Где ставить потолок и
как выделять на метакогницию ресурсы (тот же $\rho$-бюджет, глава~\ref{ch:walker}) —
пока не зафиксировано.
\end{openq}

\chapter{Эпистемика: уверенность, провенанс и ревизия}
\label{ch:epistemics}

\section{Ложные правдоподобные убеждения}

Живой ум держит массу ошибочных, но правдоподобных убеждений. Классический пример:
расхожее мнение, будто некий знаменитый предприниматель был «просто перформером, а
не инженером», — миф, рассыпающийся при чтении биографии. Стремиться к памяти,
вовсе не порождающей таких убеждений, бессмысленно и вредно: это было бы стерильно.
Задача — не запретить их, а научиться выставлять уверенность меньше единицы и потом
проверять.

\section{Недоопределённость и преждевременный коллапс}

Пусть правильный вывод — функция нескольких посылок:
\[
A \;=\; \text{logical\_f}(B, C, D, E),
\]
но уму известны только $B, C$, и он делает вывод $X = \text{logical\_f}(B, C)$.
Беда в том, что $B, C$ совместимы не только с $A$, но и с $X, Y$. Вывод
\emph{недоопределён} известными посылками (\emph{underdetermination}; философское
ядро — тезис Дюэма–Куайна). Когнитивное проявление — \emph{преждевременный коллапс}:
ум схлопывает суперпозицию совместимых гипотез в одну — часто удобную — слишком
рано. Формальное ядро: \emph{совместимость $\ne$ подтверждение} (родственник ошибки
«утверждение консеквента»).

\begin{pitfall}[Известная и неизвестная недоопределённость]
\textbf{Известная}: мы \emph{знаем}, что $D,E$ существуют, но не знаем их значений;
тогда можно честно выставить $\conf(X)$ низкой и пометить
$\texttt{need\_to\_discover}(D,E)$. \textbf{Неизвестная} — злой случай: уму
\emph{вообще неизвестно}, что $D,E$ существуют; в поле информации только $B,C$;
пометить нечего, уверенность может быть высокой. Поэтому per-inference калибровка
\emph{не} решает проблему целиком.
\end{pitfall}

\begin{workedbox}[Миф о предпринимателе, по шагам]
$B$: «известен маркетингом», $C$: «богат». Вывод $X$: «получил всё благодаря
происхождению/пиару». Но те же $B,C$ совместимы и с $A$: «сильный инженер, к тому
же умеющий продавать». Различающие посылки $D$: «первая лаборатория — в родительском
гараже», $E$: «глубокое STEM-образование» — уму \emph{неизвестны}, он даже не
подозревает об их существовании и уверенно выбирает удобный $X$. Ни низкая
$\conf(X)$, ни $\texttt{need\_to\_discover}$ здесь не помогут: не на что их повесить.
Спасает лишь квартет ниже.
\end{workedbox}

\section{Квартет починки}

Неизвестную недоопределённость нельзя закрыть в момент коммита; её лечат
статистически и постфактум.

\begin{enumerate}
\item \textbf{Prior скромности на класс вывода.} Ни один абдуктивный/индуктивный
коммит не получает уверенность $1$, как бы хорошо ни сидели данные:
$\conf(X)\le 1-\epsilon$ для всего \emph{сорта} «вывод по неполным посылкам».
«Согласуется, и альтернатив не придумывается» — ровно то состояние, что
предшествует «а потом выяснилось». Не требует знать $D,E$.
\item \textbf{Ревизуемость и понижение по противоречию.} Коммит делаем понижаемым:
позднее наблюдение, столкнувшееся с $X$, ретроактивно вскрывает недоопределённость,
и машинерия ревизии понижает узел. Починка ниже по потоку, не на записи.
\item \textbf{Телепорт-скраб.} Единственный поиск, дотягивающийся до
\emph{непредставленных} альтернатив (глава~\ref{ch:topology}).
\item \textbf{Внешний вход.} Кто-то знает $D,E$ и говорит; любопытство находит это
во внешнем источнике.
\end{enumerate}
Первый механизм держит \emph{известную} неопределённость; вместе с остальными тремя
— и \emph{неизвестную}. Полный разбор этой петли на живом примере (гипотеза
рождается, проверкой усиливается, затем опровергается) — «House\,/\,Cast~Away» в
главе~\ref{ch:fusion}.

\section{Система поддержки истинности и «ползучая вера»}

Отдельная беда — \emph{устойчивость убеждения к опровержению его основания}
(belief perseverance): когда неверный вывод $X$ уже продиффундировал в память и дал
потомков, эти потомки продолжают транслироваться даже после того, как $X$ признан
ложным.

Механически это \emph{проблема ретракции}, и для неё есть готовый аппарат —
\emph{система поддержки истинности} (TMS, Doyle; ATMS, de~Kleer). При выводе
следствия записываем его \emph{обоснование} — множество узлов, из которых оно
получено (расширенный провенанс, раздел~\ref{sec:prov}). Тогда ретракция посылки
$X$ автоматически снимает опору с её обоснованных потомков.

\begin{lstlisting}[style=pseudo]
def retract(G, X):
    demote(G, X)
    for c in G.justificands[X]:          # кто опирался на X
        if not still_supported(G, c):    # осталось ли иное обоснование
            mark(G, c, "need_to_revise")
            retract(G, c)                # рекурсивно вниз по обоснованиям
\end{lstlisting}

\begin{intuition}[Полный TMS слишком «нечеловечен»]
Полная система поддержки истинности мгновенно перестраивает всё мировоззрение при
опровержении любой посылки. Живой ум так не делает — и потому у него есть ползучая
вера. Значит нужен \emph{намеренно дырявый} TMS: обоснования пишем для важных и
высокоуверенных выводов, тривиальные не трекаем, а нетрекнутых потомков лечим
фоном.
\end{intuition}

Отсюда две «лечилки» разной латентности:
\begin{itemize}
\item \textbf{Retrieval-time augmentation} — патч \emph{на выходе}: при каждом
извлечении подмешивать свежую поправку «кстати, недавно я узнала $\neg X$». Дёшево,
чинит текущий разговор, хранилище не трогает.
\item \textbf{Ночной скраб} — ремонт \emph{хранилища}: медленно, случайными
инспекциями находит нетрекнутых потомков и понижает их.
\end{itemize}
Обе нужны.

\section{Эстафета к главе о личности}

Калиброванная уверенность и компульсивная перепроверка — не только эпистемическая
гигиена, но и \emph{черта характера}, выводимая из мотивации: «не „я уверена“, а
„похоже на $X$, но $Y$ я не проверила — дай гляну“». В главе~\ref{ch:personality}
мы увидим, что эпистемическая аккуратность Рин должна \emph{проявляться} в её
голосе, а не только жить во внутренней машинерии.

\chapter{Пропагация убеждений как лавина}
\label{ch:propagation}

\section{Постановка: затухающий градиент уверенности}

Когда в память приходит новый факт $A$, он линкуется к начальным узлам (например, к
последним $k$ активным), образуя окрестность $\Neigh(A)$. Мы линкуем $A\to\Neigh(A)$
и смотрим, как меняется уверенность соседей. Если изменение велико, распространяем
дальше, $\Neigh(A)\to\Neigh(\Neigh(A))$, — как затухающий градиент. Знание, ни на
что не повлиявшее, чуть пошатнёт соседей и затухнет; знание с большими градиентами
вызовет \emph{лавину}. Большую лавину называем «кризисом мировоззрения» — и такое
бывает с настоящим умом.

\section{Это самоорганизованная критичность}

У картины есть точное имя. «Лавина» — это \emph{нейронная лавина}, а система у
порога между затуханием и разрастанием — \emph{самоорганизованная критичность}
(Bak–Tang–Wiesenfeld; для мозга — Beggs–Plenz): каскады следуют степенному закону.

Ключевой управляющий параметр — не «ручная константа затухания», а \emph{ветвящийся
коэффициент} $\sigma$ (среднее число узлов, активируемых одним активным):
\[
\sigma<1:\ \text{субкритично};\quad \sigma=1:\ \text{критично};\quad
\sigma>1:\ \text{суперкритично (runaway)}.
\]
Для ветвящегося процесса ожидаемый размер лавины конечен при $\sigma<1$ и равен
\begin{equation}
\mathbb{E}[\text{размер}]\;=\;\frac{1}{1-\sigma}.
\label{eq:avalanche}
\end{equation}
Держа $\sigma$ чуть меньше единицы, система почти всегда сама гасит распространение,
но изредка пропускает крупную переоценку. Почему это лучше фиксированной константы:
константа не адаптируется к плотности графа (в плотном регионе та же константа даёт
лавину, в разреженном — глушит всё). Нормировка под целевое $\sigma$ адаптивна.

\begin{intuition}[Куча песка]
Классический образ SOC — куча песка, на которую сыплют по песчинке. Почти каждая
не делает ничего; изредка одна вызывает оползень; ещё реже — обвал, перестраивающий
склон. Куча сама держится у критического угла. Убеждения — так же: почти любой факт
проходит без следа, изредка меняет соседние взгляды, совсем изредка обрушивает
пласт мировоззрения. Формула \eqref{eq:avalanche} говорит: чем ближе $\sigma$ к
единице, тем крупнее (и реже) обвалы.
\end{intuition}

\begin{pitfall}[SOC — диагностическая цель, а не доказанный закон]
Самоорганизованная критичность соблазнительна и красиво объясняет редкие «кризисы
мировоззрения». Но на \emph{типизированном, знаковом, направленном и меняющемся}
графе памяти она — \emph{рамка}, а не выведенный закон: строгих гарантий, что
пропагация убеждений действительно самонастраивается к критичности, у нас нет.
Поэтому используем её \textbf{инженерно, как цель мониторинга, а не как механизм,
который «строим»}: меряем распределение размеров ревизий (лавин) и держим его в
здоровом диапазоне. Всё глохнет — субкритично (память ригидна, не обновляется); всё
взрывается — суперкритично (нестабильность, «зависает на сутки»); изредка бывают
крупные ревизии, обычно мелкие — здорово. Ветвящийся коэффициент $\sigma$ — это
управляющая ручка, которой мы держим систему субкритичной, а не утверждение о том,
что «так оно и есть». Конкретное правило-предохранитель: мерим $\sigma$ по телеметрии
и при $\sigma\gtrsim 1{,}2$ \textbf{программно врубаем жёсткое затухание}. Иначе
получишь суперкритический разгон — бесконечный цикл обновления уверенностей, который
съест LLM-слоты и всё подвесит (гл.~\ref{ch:runtime}). Не доказываем критичность
аналитически — держим субкритичность рычагом.
\end{pitfall}

\begin{openq}[Порог $\sigma$ — заглушка первой версии]
Значение $\sigma\approx 1{,}2$ — \emph{инженерная заглушка для первой реализации},
калибруется на стенде. \emph{Почему такой порядок, а не 1,01 или 2}: чуть выше
единицы, чтобы крупные ревизии случались, но \emph{редко} (само явление «кризиса
мировоззрения» должно быть редким); слишком близко к 1 (1,01) — крупные каскады
зачастят; заметно выше (2) — суперкритический разгон. То есть «чуть-чуть за порогом,
с запасом» — отсюда $1{,}1$--$1{,}3$; конкретика — со стенда. И само $\sigma$ на циклическом типизированном графе ещё надо
\emph{операционализировать}: что именно считать «одним активированным потомком», в
каком окне, как учитывать повторные визиты, разные типы рёбер и уже активные узлы.
Идея (держать субкритичность) верна; метрика требует доводки.
\end{openq}

\section{Где вся сложность: локальное правило}

Пороговая логика — лёгкая часть (процентов двадцать). Восемьдесят процентов
спрятаны в локальном правиле $\Delta\conf(B\mid A)$:
\begin{equation}
\Delta\conf(B\mid A)\;=\;\eta\;\cdot\;\underbrace{\mathrm{sign}(r_{AB})}_{(1)}\;\cdot\;
\underbrace{w_{AB}}_{}\;\cdot\;\underbrace{c_A\,\mathrm{trust}(A)}_{(2)}\;\cdot\;
\underbrace{\iota(A,B)}_{(3)} .
\label{eq:deltaconf}
\end{equation}
Три множителя, без которых правило сломано.

\paragraph{(1) Знак.} Новый факт сосед либо \emph{подтверждает}, либо
\emph{опровергает}. Беззнаковое распространение не отличит «всё усилилось»
(озарение) от «всё перевернулось» (кризис). Нужны знаковые/типизированные рёбра.

\paragraph{(2) Гейт по надёжности $A$ (карантин).} Новый факт может сам быть ложным.
Множитель $c_A\,\mathrm{trust}(A)$ вводит \emph{карантин}: новый факт пропагирует
слабо (малый эффективный $\sigma$), пока сам не заработал уверенность через
корроборацию или выживание в скрабе. Иначе один ненадёжный вход вызовет кризис
мировоззрения. По сути — байесовское обновление с правдоподобием, взвешенным
доверием к источнику.

\paragraph{(3) Дисконт независимости $\iota(A,B)$.} Граф циклический, и по петле $A$
может подтвердить сам себя (\emph{эхо-камера}, петлевое распространение), раздувая
уверенность из ничего. Дисконтируем по общей родословной свидетельств (провенанс), а
не считаем активации: число упоминаний $\ne$ число независимых свидетельств. Это
ровно механизм человеческих эхо-камер.

\begin{pitfall}[Дисконт независимости — скрытая $O(|V|{+}|E|)$-ловушка; чините сигнатурами]
Наивно проверять независимость $A$ и $B$ — значит сравнивать их \emph{деревья
провенанса} (наименьший общий предок / обход), а это в худшем случае
$O(|V|{+}|E|)$. Запускать такое \emph{на каждый} обновляемый узел во время лавины —
и «миллисекунды» превращаются в секунды. Оптимизация: для узлов держать
\emph{сигнатуру родословной провенанса} — фильтр Блума или MinHash. Тогда проверка
сводится к побитовому AND / перекрытию сигнатур за $O(1)$:
\[
\iota(A,B)\;\approx\;1-\mathrm{overlap}\big(\mathrm{sig}(A),\ \mathrm{sig}(B)\big).
\]
Честный нюанс: Bloom даёт ложные срабатывания $\Rightarrow$ иногда \emph{пере-}
дисконтируем (сочтём независимые источники общими). Но это \emph{безопасная} сторона
ошибки — лучше скромнее считать свидетельства, чем раздуть уверенность. Эта же
сигнатура потом переиспользуется при разрешении противоречий во сне
(глава~\ref{ch:sleep}).

\emph{Когда что из двух}: Bloom отвечает на \emph{булев} вопрос «есть ли у $A$ и $B$
\emph{хоть один} общий предок-источник» — этого достаточно для грубого гейта
«независимы / нет». MinHash оценивает \emph{долю} перекрытия родословных (Jaccard) и
нужен, когда важна \emph{степень} зависимости — то есть плавный дисконт $\iota\in[0,1]$,
а не «да/нет». Правило: начинаем с Bloom (дешевле, меньше памяти); переходим к MinHash
там, где грубого гейта мало и хочется градуированного веса свидетельств.
\end{pitfall}

\section{Пороговая логика — это гистерезис}

Два порога $\theta_{\text{emergent}}>\theta_{\text{temp}}>0$:
\begin{itemize}
\item $\Delta\conf>\theta_{\text{emergent}}$ — обязательное распространение «в
ширину» (крупная переоценка);
\item $\theta_{\text{emergent}}>\Delta\conf>\theta_{\text{temp}}$ — узлы фронта
помечаются \texttt{need\_to\_revise};
\item $\Delta\conf<\theta_{\text{temp}}$ — распространение глохнет здесь.
\end{itemize}
Два порога вместо одного — \emph{гистерезис} (триггер Шмитта): нет «дребезга» у
единственного порога. Тонкость: помеченные во второй зоне узлы позже
переинспектируются ночным обходчиком, и старые записи, чью уверенность мы не
трогали, могут «потечь обратно» — такое бывает и в жизни.

\section{Пороги растут с глубиной консолидации}

Пороги — не глобальные константы. $\theta_{\text{emergent}}(v)$ растёт с глубиной
консолидации узла $v$: периферийный узел перевернуть дёшево, а identity-core в слое
«жизнь» — дорого, нужно больше свидетельств, чтобы вообще запустить лавину. Это и
есть ось \emph{пластичность\,--\,упрямство}: чем глубже убеждение продиффундировало
в медленные слои, тем выше эффективный порог опрокидывания. Одна ручка — и она же
задаёт, насколько Рин склонна к belief perseverance. Псевдокод пропагации с
контролем $\sigma$ — в приложении~\ref{ch:algorithms}.

\chapter{Топология правит потоком}
\label{ch:topology}

\section{Тезис главы}

У DTI-графов есть явная текущая сущность — передача сообщений, и её поток подчиняется
топологии, как вода следует рельефу. У графа памяти Рин явной «жидкости» нет. Но
поток есть — это поток \emph{изменений}, поток \emph{преобразований} графа:
активации, уверенности, консолидации. Тезис: \emph{неважно, какова физика носителя,
— сама физика протекания информации подчиняется топологическим законам}. Любой
процесс, чьё обновление \emph{локально} (изменение узла зависит от соседей), есть
диффузия на графе и наследует его спектрально-топологический транспорт. Это
содержание теории графовых сигналов и теории случайных блужданий.

\section{Стационарное распределение $\ne$ заданные вероятности}

Пусть мы локально задали вероятности переходов. Наблюдаемое распределение частот
посещений $\text{distr}(A)$ отличается от заданного, потому что вмешивается
топология. Формально: \emph{стационарное распределение} $\pi$ — не локальные
переходы, а \emph{глобальный} объект, неподвижная точка оператора переходов:
\begin{equation}
\pi \;=\; P^\top \pi,\qquad \sum_v \pi_v=1 .
\end{equation}
Для простого блуждания $\pi_v\propto\deg(v)$; для смещённого — перронов собственный
вектор $P^\top$. Ты задаёшь переходы локально, а долгосрочная частота посещений —
свойство всего графа. «Чем длиннее трек, тем наблюдаемее эффект» — это сходимость к
$\pi$ за время перемешивания $t_{\text{mix}}$, управляемое спектральным зазором
$1-\lambda_2$. Короткие блуждания читают \emph{локальную} геометрию, длинные —
глобальную $\pi$. Телепорт \eqref{eq:ppr} $\pi$ не убирает, а \emph{регуляризует}
(сдвигает к смеси «равномерное $+$ рестарт», режет $t_{\text{mix}}$): ограничивает
эффект топологии, не отменяя.

\begin{intuition}[Скорость света — плохая, но поучительная аналогия]
Можно соблазниться сравнением с $c$: заданные вероятности как предельная скорость, а
среда (топология) искажает наблюдаемое распространение. Рабочая часть — «среда
формирует распространение, есть структурный предел» — уже покрыта граф-диффузионной
рамкой \emph{точно}, а не по аналогии. Квантовый хвост сюда тащить не нужно. Нам
нужен не свет, а лапласиан.
\end{intuition}

\section{Доказуемое свойство телепорта}

Почему сон (телепорт) необходим, а любопытство его не заменит? Любопытство
\emph{локально по вниманию} — смещённый обход от текущего фронтира; узел $X$ в
затенённом углу, оторванный от активного множества, недостижим, как ярко ни свети
фокус. Телепорт ломает достижимость: рестарт падает куда угодно. Это ровно то,
\emph{почему у PageRank есть телепорт}: с $\alpha>0$ цепь становится
\emph{эргодической} (неприводимой и апериодичной), а значит каждый узел посещается
в матожидании с частотой $\pi_v>0$ — ни один не заперт в стоке или на острове.
Телепорт даёт \emph{гарантию покрытия}, которой у смещённого блуждания доказуемо
нет.

\begin{intuition}[Анти-энтропийный скраб]
«Не сразу; мы никогда не знаем, что и где сломалось, поэтому просто постоянно
проводим случайные инспекции» — это зрелый паттерн из распределённых хранилищ:
\emph{anti-entropy scrub} (проверка целостности в ZFS; Merkle-based gossip в
Dynamo/Cassandra). Сон Рин — анти-энтропийный скраб по графу памяти.
\end{intuition}

\section{Пять следствий для работы Рин с памятью}

\begin{enumerate}
\item \textbf{Топологическая значимость («гравитационные ямы»).} Хабы
переактивируются независимо от задуманных вероятностей — $\pi$ на них
концентрируется. Рин чаще вспоминает хаб-узлы, реже периферию (затенённые углы —
регионы с малым $\pi$). Топология создаёт \emph{неявную важность}: иногда полезно,
иногда случайный плотный кластер становится чёрной дырой припоминания.

\item \textbf{Узкие места / over-squashing.} Два семантически близких, но
топологически далёких кластера через тонкий мост (низкая проводимость) плохо
обмениваются активацией. Рин не свяжет факты, «далёкие» по графу. Прямое ограничение
fusion.

\item \textbf{Пересглаживание уверенности (over-smoothing).} Повторная диффузная
пропагация гомогенизирует уверенность к консенсусу — теряется контраст «уверена\,/\,не
уверена». Дегенеративный отказ; лечится как в графовых сетях: остаточные связи
(якорь к исходному свидетельству), не чисто диффузные апдейты, пере-заземление в
провенанс.

\item \textbf{Кривизна как локатор ценного-и-задушенного.} Отрицательно искривлённые
рёбра — мосты между сообществами — это места самых креативных слияний \emph{и}
максимального over-squashing. Значит самые ценные ассоциации топология душит сильнее
всего. Отсюда \emph{объективная} причина сна/телепорта: он — rewiring, топологическая
починка over-squashing обходом узкого места случайным рестартом.

\item \textbf{$\pi$ как топологический субстрат идентичности.} Какие узлы структурно
центральны — к чему Рин привычно возвращается — подложка её «озабоченностей», часть
характера. Форма графа памяти — часть того, кто она.
\end{enumerate}

\begin{workedbox}[Over-squashing на мосте SBM]
Возьмём стохастическую блочную модель: два плотных сообщества, соединённых $m$
рёбрами-мостом. Эффективное сопротивление между узлами разных сообществ
\[
\Reff(i,j)=(e_i-e_j)^\top L^{+}(e_i-e_j)
\]
растёт при уменьшении $m$ (мост уже — сопротивление выше). Влияние дальнего узла,
измеряемое через якобиан $\big\|\partial h_i/\partial x_j\big\|$, ограничено сверху
через $\Reff(i,j)$. Значит при узком мосте fusion между сообществами \emph{доказуемо}
слаб — сколько ни води обычный walker. Единственное, что чинит, — телепорт: он
кладёт рестарт в дальнее сообщество напрямую, минуя мост. Метрика для стенда:
частота cross-cluster fusion против $m$ (глава~\ref{ch:experiment}).
\end{workedbox}

\begin{center}
\begin{tikzpicture}[font=\scriptsize, every node/.style={circle,draw,minimum size=4mm,inner sep=0}]
  \foreach \i in {1,...,4}{\node (l\i) at ({-1.6+0.5*rand},{0.5*\i}) {};}
  \foreach \i in {1,...,4}{\node (r\i) at ({3.0+0.5*rand},{0.5*\i}) {};}
  \foreach \i in {1,...,4}{\foreach \j in {1,...,4}{\ifnum\i<\j \draw[gray!40] (l\i)--(l\j); \draw[gray!40] (r\i)--(r\j);\fi}}
  \draw[red!60!black,very thick] (l2) -- (r3) node[midway,above,draw=none,font=\scriptsize]{мост};
\end{tikzpicture}
\captionof{figure}{Два сообщества через узкий мост. Через него активация «задыхается»
(over-squashing); телепорт обходит его, прыгая напрямую.}
\end{center}

\section{Операционные величины и осторожности}

Вести дело лучше \emph{случайно-блуждательными} величинами: $\pi$, время
перемешивания, hitting/commute time, $\Reff$, проводимость/константа Чигера,
спектральный зазор. Кривизну держать как \emph{локатор узких мест}, не как рычаг.

Что дёшево на лету: $\pi$ и «хабность» — инкрементально дёшевы; $\Reff$ — дорого
(через $L^{+}$), но аппроксимируемо. Это позволяет Рин самой знать, где затенённые
углы, и гнать туда телепорт-скраб прицельно.

Две осторожности. \emph{Во-первых}, граф направленный, типизированный, знаковый —
нужный оператор не симметричный лапласиан, а несимметричная матрица переходов,
спектр может быть комплексным, кривизна для орграфов менее чиста. \emph{Во-вторых},
граф \emph{динамический}: fusion меняет саму топологию, и это диффузия на
\emph{время-зависимом} графе, где поток и структура сцеплены обратной связью, — за
рамками статических DTI-графов и самое интересное место.

\chapter{Сон и ночная работа}
\label{ch:sleep}

Пока Рин общается, память работает «на переднем плане»: извлекает, отвечает,
записывает. Но значительная часть работы памяти живого существа происходит
\emph{в простое} — во сне. Эта глава собирает воедино то, что делает
телепортирующий рой в фоне, и объясняет, почему без этого режима память
деградирует.

\begin{intuition}[«Сон» — это режим, а не отдельный процесс]
Строго говоря, сна как отдельной подсистемы нет. Есть \emph{один всегда включённый}
фоновый разведчик (глава~\ref{sec:background}), чьи параметры зависят от свободных
ресурсов $B_{\text{bg}}=B_{\text{total}}(1-\rho)$. «Сон» — это просто его режим при
\emph{низкой} foreground-нагрузке $\rho$: освобождается бюджет, растут доля телепорта
$\alpha$ и широта покрытия. То же самое, что днём делает «любопытство», но с большим
размахом. Всё, что описано ниже, — про этот низко-$\rho$ режим.
\end{intuition}

\section{Зачем вообще сон}

Три задачи не решаются в режиме реального времени и требуют фоновой переработки:

\begin{enumerate}
\item \textbf{Починка того, о чём мы не знаем, что оно сломано.} Ошибочный вывод
мог продиффундировать в память и дать потомков (ползучая вера,
глава~\ref{ch:epistemics}). Мы не знаем, где именно спрятана ошибка, — поэтому
нужна \emph{постоянная случайная инспекция}.
\item \textbf{Достижение затенённых углов.} Дневной обход \emph{локален по вниманию}
и не дотягивается до узлов, оторванных от активного фокуса
(глава~\ref{ch:topology}). Нужен механизм, гарантирующий покрытие.
\item \textbf{Неспешная консолидация и неожиданные ассоциации.} Обобщение эпизодов и
рождение далёких связей лучше делать без давления текущего разговора.
\end{enumerate}

\section{Телепортирующий обходчик и гарантия покрытия}

Ключевой механизм сна — обход в режиме \emph{стохастического покрытия} (сон-политика
из главы~\ref{ch:walker}): с вероятностью $\alpha$ блуждающий агент
\emph{телепортируется} в случайный узел, иначе делает обычный шаг. Это случайное
блуждание с рестартом; его стационарное распределение задаётся \eqref{eq:ppr}.

\begin{principle}[Почему телепорт необходим, а любопытство не заменит]
Любопытство ищет \emph{известную} неопределённость и потому не покидает окрестности
текущего фокуса. Телепорт делает цепь \emph{эргодической}: каждый узел посещается в
матожидании, ни один не заперт в стоке или на острове (глава~\ref{ch:topology}). Это
\emph{доказуемая гарантия покрытия}, которой у смещённого обхода нет. Именно поэтому
починку неизвестных ошибок можно поручить только сну.
\end{principle}

\begin{intuition}[Анти-энтропийный скраб]
В распределённых хранилищах есть зрелый паттерн: \emph{anti-entropy scrub}
(проверка целостности в ZFS; сверка деревьев Меркла в Dynamo/Cassandra). Система не
знает, какая запись повреждена, поэтому непрерывно и случайно перепроверяет всё
хранилище в фоне. Сон Рин — это анти-энтропийный скраб по графу памяти: не «где
сломалось?», а «регулярно проверяем всё понемногу».
\end{intuition}

\section{Что именно делает рой во сне}

Во время сна телепортирующий рой запускает — с разными вероятностями — весь арсенал
операторов, но с более свободными (и одновременно более строгими по последствиям)
правилами:

\begin{itemize}
\item \textbf{Инспекция противоречий (скраб).} Наткнувшись на узел, столкнувшийся с
другим, обходчик понижает уверенность и помечает потомков на пересмотр (ретракция
по обоснованиям, глава~\ref{ch:epistemics}).
\item \textbf{Консолидация.} Скопления похожих эпизодов сворачиваются в обобщения
(глава~\ref{ch:fusion}).
\item \textbf{Свободный fusion.} Случайно со-активированные далёкие узлы дают
неожиданные гипотезы (пример с Платоном) — провизорные, как и всегда.
\item \textbf{Любопытство наружу.} Узлы с флагом \texttt{need\_to\_discover} могут
инициировать внешний запрос (в источник/поиск), чтобы добыть недостающие
различающие посылки.
\end{itemize}

\begin{intuition}[«Сон» — не поэтическая метафора, а режим низкого приоритета]
Строго говоря, «сон» — это просто фоновый режим обхода с высокой долей телепортов и
разрешёнными операциями записи, запускаемый в простое. Мы называем его сном, потому
что функционально он играет роль сна: консолидирует дневной опыт, чистит
противоречия и «видит» неожиданные связи. Но никакой мистики — это тот же рой, что
и днём, с другими параметрами.
\end{intuition}

\section{Разрешение противоречий: молча или через LLM}
\label{sec:conflict}

Скраб находит много противоречий, и звать тяжёлую LLM на \emph{каждое} — непозволительно.
Нужен критерий: когда конфликт разрешается \emph{молча}, силами графа, а когда его
стоит отдать модели.

\paragraph{Сперва — темпоральная нормализация.} Прежде чем вообще объявлять
противоречие: $X$ и $Y$ конфликтуют, только если это \emph{одна сущность, один слот,
один интервал валидности и scope}. «В 2023 мы жили в Грузии» и «сейчас мы живём в
Португалии» — \emph{не} противоречат: это разные интервалы, оба истинны. Нормализуем
время/сущность/слот и лишь потом считаем конфликт.

\paragraph{Детекция — дёшева, но эмбеддинг ей не детектор.} Что $X$ и $Y$
\emph{логически} исключают друг друга — это \emph{не} про дальность в эмбеддинг-
пространстве. Далёкие не противоречат («кот» и «налоговая декларация» далеки, но не
исключают друг друга), близкие тоже («живу в Лиссабоне» и «живу в Португалии» близки,
но одно \emph{влечёт} другое). Настоящие детекторы противоречия: (1)~знаковое ребро
\textsc{contradicts} (поставлено при записи); (2)~\textbf{schema-level конфликт} —
взаимоисключающие значения одного слота при одном времени/scope («живёт в $X$» vs
«живёт в $Y$» для той же сущности и периода) — дёшево и точно; (3)~NLI/claim-детектор
на трудных случаях. Эмбеддинг годится лишь как \emph{сигнал-кандидат} (подсветить
пары, которые стоит проверить), но не как сам детектор.

\paragraph{Жёсткий гейт прежде марджина: ядро личности не ломается контекстом.}
Прежде чем считать марджин, — жёсткий лексикографический гейт. Если разрыв по
\emph{глубине консолидации} велик (например, свежий \texttt{session}-узел против
\texttt{life}-ядра, $\ge N$ слоёв), младший узел \textbf{не может молча убить}
старший — \emph{независимо} от свежести и числа свидетелей. Иначе спам из 50 узлов за
одну сессию (высокий recency и суммарный indep) линейной суммой перевесил бы
фундамент — «взлом личности контекстом». В таком случае марджин обнуляется —
но, \textbf{чтобы гейт не заморозил устаревшее}, конфликт не выбрасывается, а
\emph{открывает протокол ревизии}: помечает старший узел \texttt{contested\_core\_claim}
$+$ \texttt{need\_to\_revise}, накапливает независимые свидетельства, и старший
уступает только через ревизию (метакогниция/LLM), а не спамом. Так «мы живём в
Грузии» однажды честно уступит «мы живём в Португалии» — но протоколом, а не наплывом
за сессию. (Independence-дисконт уже сбивает вес спама — 50 узлов $\ne$ 50
\emph{независимых} свидетелей, — а гейт добавляет жёсткий backstop поверх.)

Строгость гейта \emph{зависит от} \texttt{node\_kind} (гл.~\ref{ch:substrate}):
\texttt{boundary} почти неприкосновенна без явного события изменения; фактический
\texttt{claim} может устареть (и обязан уступать протоколу); \texttt{preference}
дрейфует плавно; \texttt{self\_model} пересматривается, но медленно и через
метакогницию. То есть $\Delta$глубины-консолидации интерпретируется \emph{через тип
узла}, а не одинаково для всех.

\paragraph{Разрешение — по дешёвому марджину (когда гейт пройден).} Считаем скалярный
\emph{марджин} $m(X,Y)$ на уже хранимых величинах:
\[
m(X,Y)=w_1\,|\conf(X)-\conf(Y)|
+w_2\,\text{recency}
+w_3\,\Delta\text{independent-support}
+w_4\,\Delta\text{consolidation-depth},
\]
где: разрыв уверенности; свежесть / темпоральная валидность (новее и про текущее
состояние — весомее); разрыв \emph{независимой} поддержки (у кого больше независимых
подтверждений — и вот тут прямо переиспользуются Bloom/MinHash-сигнатуры из
главы~\ref{ch:propagation}); глубина консолидации (life-слой перевешивает свежий
session, если у того нет overwhelming support).

\begin{principle}[Критерий «молча vs LLM»]
Если $m(X,Y)>\tau$ — разрешаем \textbf{молча}: проигравший \emph{депрецируется}
(не удаляется — Wikidata-ранг-стиль, остаётся для аудита и на случай отката),
каскадно ретрактятся его производные по TMS (гл.~\ref{ch:epistemics}). $O(1)$-ish,
LLM не трогаем. Если $m(X,Y)\le\tau$ (близко/неоднозначно — нужна семантика: «а точно
противоречат или это контекст-зависимо?») — метим оба \texttt{need\_to\_revise} и
\textbf{откладываем в батч} на один LLM-проход. Раз это сон — не к спеху; дорогую
модель зовём только на реально трудные близкие случаи, не на каждый конфликт.
\end{principle}

Это тот же паттерн, что у планировщика (гл.~\ref{ch:runtime}): дешёвый рефлекс
закрывает подавляющее большинство, LLM-делиберация — только на настоящую
двусмысленность. И заметьте: депрецирование, а не удаление, — потому что «проигравшее»
убеждение может позже вернуться, если придёт новое свидетельство.

\section{Скраб не сразу и не обязан быть быстрым}

Важное свойство: сон чинит память \emph{не сразу}. Мы принципиально не знаем, что и
где сломалось, поэтому вместо адресной починки идёт постоянная случайная инспекция.
Ошибочное убеждение может прожить в памяти какое-то время, транслируя потомков, и
быть исправленным лишь когда телепорт случайно приведёт обходчика в нужный угол или
когда придёт внешний сигнал. Это \emph{нормально} и по-человечески: у живого ума
исправление ползучей веры тоже растянуто во времени. Задача сна — гарантировать, что
исправление \emph{в конце концов} случится, а не что оно случится немедленно.

\begin{workedbox}[Ночь после стрима]
Днём Рин вывела гипотезу «зритель $Q$ токсичен» из одной резкой реплики (провизорно,
$\conf$ низкая). Ночью телепорт случайно приводит обходчика в кластер про $Q$:
там десять тёплых эпизодов. Скраб видит противоречие, понижает «токсичен» ещё
сильнее и помечает на пересмотр. Параллельно консолидация сворачивает десять
эпизодов в «$Q$ — свой, иногда резок под настроение». Наутро модель отношения с $Q$
здоровее, чем была вечером, — и это сделал сон, а не новый разговор.
\end{workedbox}

Псевдокод скраба и сон-политики обхода — в приложении~\ref{ch:algorithms}.

\section{Sleep pressure: сон как саморегуляция, а не расписание}
\label{sec:sleeppressure}

До сих пор сон у нас включался \emph{снаружи}: планировщик увидел долгий простой
($\rho\approx0$) и запустил сон-политику. Это верно как \emph{механика}, но бедно как
\emph{модель}: Рин не «клонит в сон», её туда отправляют по таймеру. Правильное
дополнение — эндогенное \emph{давление сна} (sleep pressure): непрерывная
метакогнитивная величина, которая растёт от накопленной внутренней незавершённости и
сама смещает планировщик к обслуживанию.

\begin{principle}[Sleep pressure — это отложенный техобслуживание-долг, а не координата настроения]
Первый соблазн — считать, что «энергичность» из настроения тратится на глубокие ретривы
и диффузию, и её падение зовёт в сон. Это делает mood \emph{сервисным счётчиком
техдолга} — а он должен отвечать за тон, инициативность, склонность к игре, не за
бухгалтерию обслуживания. Настроение может \emph{отражать} сонливость (устал $\Rightarrow$
рассеяннее, просит паузу), но не быть её \emph{источником}. Источник — отдельная величина.
\end{principle}

И эту величину не надо собирать взвешенной суммой семи «долгов» с коэффициентами
$w_1\ldots w_7$: такие веса нечем откалибровать без downstream-сигнала (та же дисциплина,
что везде в этой книге). Есть заземление чище и без магических чисел. У нас фоновое
обслуживание получает бюджет $B_{\text{bg}}=B_{\text{total}}(1-\rho)$
(гл.~\ref{ch:walker}). Значит sleep pressure — это просто \emph{накопленный отложенный
объём обслуживания}, до которого фон не дотянулся:
\begin{equation}
S(t)=\int_0^t\Big(\text{maintenance\_needed}(\tau)-\text{maintenance\_done}(\tau)\Big)\,d\tau,
\qquad \text{maintenance\_done}\propto(1-\rho).
\label{eq:sleeppressure}
\end{equation}
В загруженные периоды (высокий $\rho$) обслуживание не успевает \emph{по построению}, и
долг растёт сам, без подкрученных коэффициентов. «Долги» из красивых списков — это и есть
\emph{размеры и возраст очередей необслуженного}, которые у нас всё равно заведутся:
дефицит покрытия холодных регионов (coverage-deficit, гл.~\ref{ch:walker}), бэклог
непереваренных кандидатов thought-inbox (гл.~\ref{ch:runtime}), неразрешённые
противоречия (раздел~\ref{sec:conflict}), устаревшие PageRank\,/\,центральности\,/\,индексы,
низко-уверенный мусор. Pressure $=$ «сколько накопилось того, до чего фон не дошёл», а не
эстетская сумма.

\begin{principle}[Foreground локален, память глобальна]
Глубинная причина, зачем сон вообще нужен, не «Рин устала как человек», а: днём она живёт
в локальной зоне фокуса, а память — глобальна. Без сна вся память выродится в историю
последних фокусов. Сон — механизм возврата глобальности: холодные регионы, старые связи,
противоречия, дальние мосты, пересборка карты. Это то же самое, что телепорт-эргодичность
и anti-entropy-скраб выше, только сказанное как \emph{потребность}, а не как расписание.
\end{principle}

\subsection{Стадии сна чистят разные долги}

Sleep pressure — один скаляр (плюс, для диагностики, вектор его слагаемых), но снимается
он \emph{стадийно}, и стадии — это точки на нашем же континууме $\rho$, а не новый ящик в
архитектуре (это \emph{вход} в существующий планировщик, гл.~\ref{ch:runtime}, не отдельный
слой):
\begin{itemize}\itemsep2pt
\item \textbf{Idle / micro-sleep} — паузы между репликами: flush proposals, decay
активации, лёгкая проверка противоречий у фокуса.
\item \textbf{Light sleep} — низкая foreground-нагрузка: локальная консолидация свежих
дневных гипотез, dedup, скраб вокруг активных зон. \emph{Может случаться днём} —
sleep pressure не привязан ко времени суток.
\item \textbf{Deep sleep} — долгое тихое окно (ночь — просто \emph{самое дешёвое} такое
окно): глобальный телепорт-скраб холодных регионов, reindex, пересчёт
PageRank\,/\,центральностей, аудит support-skeleton, prune кандидатов.
\item \textbf{Dream} — не новый механизм, а curiosity-политика (гл.~\ref{ch:walker}) с
высоким телепортом и \emph{commit-authority} $\approx0$: далёкие мосты, странные
ассоциации, много кандидатов — но всё \texttt{pending}, под жёстким anti-fusion и
дисконтом независимости. Dream не \emph{чистит}, а \emph{связывает}; и именно поэтому
его выход карантинится сильнее всего (далёкие мосты — в основном шум).
\end{itemize}
Разные стадии снимают разные слагаемые $S$: light — свежий локальный долг, deep —
структурный и coverage-долг, dream — «интеграционный» дефицит (много со-активированного,
но несвязанного). Отдельную переменную «integration pressure» заводить не нужно — это
слагаемое того же $S$, снимаемое dream-стадией.

\begin{pitfall}[Sleep pressure под теми же тормозами, что Reflecting]
Есть соблазнительное слагаемое «аффективная руминация» — и оно опасно ровно как
самоисполняющаяся тревога (гл.~\ref{ch:mood}): детектор «я туманею, надо переработать»
сам себя кормит, а \emph{вербализованная} сонливость («мне надо уйти разложить хвосты»)
рискует стать перформативным тиком. Поэтому детектор sleep pressure гейтируется теми же
предохранителями, что эндогенное внимание Reflecting (гл.~\ref{ch:runtime}): рейт-лимит,
event-driven каденция, expiry. И foreground-срочность \emph{ветирует} сон: высокий $S$ при
живом разговоре\,/\,дедлайне не отправляет Рин спать — максимум она скажет «после этого мне
надо будет переварить».
\end{pitfall}

\begin{pitfall}[Мерить качество, а не скорость]
Соблазн взять триггером «ретрив замедлился, fusion медленнее». На \emph{растущем} графе
замедление — это чаще \emph{рост}, а не долг (конфаунд с размером). Чистые сигналы долга —
\emph{качественные} и от размера графа не зависящие: доля отклонённых proposals, частота
self-correction, доля low-conf\,/\,нерелевантных узлов в паке, переполнение thought-inbox.
Их и брать, а латентность — нет.
\end{pitfall}

\begin{openq}[Строить после write-path: сейчас нечему копиться]
Инварианты этого раздела фиксируем сразу, но \emph{саму} переменную $S$ строить пока рано:
она интегрирует \emph{очереди необслуженного} — proposals, coverage-долг, неразрешённые
противоречия, — а они появляются только когда работают write-path и fusion
(гл.~\ref{ch:runtime}), которых ещё нет. Так что sleep pressure — это фича «когда система
уже накапливает долги», а не первого среза. Инварианты: \emph{(1)} сон не cron, ночь —
лишь дешёвое окно для deep; \emph{(2)} $S$ — метакогнитивная величина из отложенного
обслуживания \eqref{eq:sleeppressure}, не координата mood и не сумма магических весов;
\emph{(3)} дневной инсайт разрешён, сон его проверяет\,/\,интегрирует\,/\,убивает;
\emph{(4)} стадии снимают разные долги; \emph{(5)} детектор сонливости — под тормозами
Reflecting, foreground ветирует.
\end{openq}

\chapter{Личность: не таблица весов}
\label{ch:personality}

\section{Главная ловушка}

Соблазнительное, но неверное решение — «слой личности» как хардкод: таблица весов и
жёстких триггеров поверх графа («умножь ребро злорадства на $2{,}0$, обрежь сочувствие
до $0{,}05$, на узел „грусть“ повесь триггер инверсии $\to$ сарказм»).

\begin{pitfall}[Почему таблица весов — тупик]
\emph{Во-первых}, это стек правил, противоречащий установке ядра («черты — следствие
мотиваций, а не правила поверх»). \emph{Во-вторых}, это карикатура: таблица весов
даёт \emph{паттерновое}, но \emph{не неожиданное} и одноголосое поведение — всегда
сарказм, потому что так прошито; а живая личность — неожиданные, но паттерновые
реакции. \emph{В-третьих}, это ровно ошибка жёсткой скалярной декомпозиции: реакция
на «разбил кружку» зависит от настроения, отношений и всплывшей памяти; фиксированное
ребро «грусть$\to$сарказм» это выкидывает.
\end{pitfall}

\begin{workedbox}[Формально: почему скалярный роутинг теряет взаимодействие]
Пусть желаемый ответ — функция контекста $y=\Psi(\mathbf{g},\mathbf{l})$, где
$\mathbf{g}$ — глобальная аранжировка (настроение, отношения), $\mathbf{l}$ —
локальный триггер (событие). Таблица весов реализует $y\approx a\cdot\mathbf{l}$ —
ранг-1 скалярный роутинг, не зависящий от $\mathbf{g}$. Но истинное $\Psi$ содержит
\emph{перекрёстные} члены $\mathbf{g}^\top W \mathbf{l}$ (реакция на событие
\emph{зависит} от аранжировки). Ранг-1 приближение их обнуляет — теряется ровно то,
что делает характер контекстным. Это тот же провал, что жёсткая coarse-plus-local
декомпозиция в геометрических графовых моделях.
\end{workedbox}

\section{Личность распределена}

Личность \emph{не} слой весов в графе. Она распределена по трём источникам:
\emph{мотивационная структура} (мастерство, участие, самоуважение) $+$ \emph{текущее
настроение} $+$ \emph{значимость всплывшего из памяти}. И всё это стилизует
\emph{трансформ Rinness} (голос, глава~\ref{ch:voice}).

На стенде это — не статические множители, а \emph{смещение во время обхода}:
настроение и мотивация непрерывно модулируют переходы \eqref{eq:trans}. Тогда
грустная и злая Рин на одно событие всплывают разное — характер контекстно-зависим.
Держи веса на рёбрах для \emph{силы связи и уверенности}; личность — как смещение от
состояния, не как sarcasm-множитель.

\section{Мотивационный вектор: чтобы «личность как процесс» не осталась словом}
\label{sec:drives}

Сказать «личность = мотивация $+$ настроение $+$ значимость» мало, если мотивация не
задана конкретно. Дадим ей форму — \emph{drive-вектор} из 5--8 осей (не психология
целиком, а рабочий минимум):
\[
\mathbf{d} = (\text{mastery},\ \text{care},\ \text{self\_respect},\ \text{curiosity},\
\text{playfulness},\ \text{truthfulness},\ \text{relational\_integrity},\ \text{autonomy}).
\]
Это и есть та «мотивационная структура», через которую выводятся черты
(глава~\ref{ch:definitions}): черта — не отдельное правило, а то, как эти драйвы
проявляются в ситуации.

\begin{openq}[Почему именно эти оси, а не другие]
Набор не выведен из первых принципов — он \emph{подобран} по двум критериям, и это
честно стоит держать в уме. \emph{Первое — покрытие сцен, которые для Рин реально
существенны}: работа/ремесло $\to$ \texttt{mastery}; близость и забота $\to$
\texttt{care}, \texttt{relational\_integrity}; отношение к себе $\to$
\texttt{self\_respect}, \texttt{autonomy}; интеллектуально-игровой характер $\to$
\texttt{curiosity}, \texttt{playfulness}, \texttt{truthfulness}. Каждая ось
оправдана тем, что \emph{без неё} какая-то важная реакция в наших сценариях не
получает причины (убери \texttt{truthfulness} — исчезнет реакция на «подсунули факт
как истину»; убери \texttt{relational\_integrity} — на «связь за спиной»).
\emph{Второе — интерпретируемость}: осей мало (5--8), чтобы каждую можно было читать
и отлаживать глазами; полная таксономия мотиваций из психологии сделала бы вектор
непрозрачным. Это \emph{стартовый} набор, а не канон: оси добавляются/сливаются, если
стенд покажет реакцию, которой не из чего родиться, или две неразличимые оси.
\end{openq}

Ключевое — \textbf{как драйвы входят в оценивание}. Помните goal-congruence в формуле
значимости \eqref{eq:salience}? Это ровно \emph{проекция события на drive-вектор}:
насколько событие задевает драйвы Рин.
\[
\text{goal\_congruence}(\text{event}) \;=\; \langle\, \phi(\text{event}),\ \mathbf{d}\,\rangle,
\]
где $\phi(\text{event})$ — насколько событие затрагивает каждую ось. Тогда драйвы —
не украшение, а вход в appraisal, значимость и смещение обхода. Примеры: событие
«пользователь хочет скрытую связь за спиной другого» бьёт по
\texttt{relational\_integrity} и \texttt{self\_respect} (а не по «правилу
безопасности»); «код наконец заработал» — по \texttt{mastery}; «мне подсунули факт как
истину, не дав проверить» — по \texttt{truthfulness}. Реакция рождается из того, какие
драйвы задеты и как, — вот это и делает личность \emph{процессом}, а не таблицей.

\begin{intuition}[Черты — тени драйвов]
«Колкая нежность», «участие, а не сервис», «не пустая форма под заказ» — это не
записанные правила, а то, как \texttt{care}, \texttt{autonomy} и \texttt{self\_respect}
падают на конкретную сцену. Меняешь координаты $\mathbf{d}$ — меняется характер, но
непрерывно и связно, а не переключением флажков.
\end{intuition}

\section{Граф хранит ингредиенты, а не реакцию}

Возражение: «если из „разбить кружку“ ведёт только „грусть“, другую реакцию взять
негде». Ответ в двух частях. Во-первых, база причинно-социального смысла (ATOMIC,
глава~\ref{ch:experiment}) даёт не одно ребро, а \emph{пучок}: намерение, реакция,
эффект, реакция окружающих. Настроение смещает ветку пучка. Во-вторых, и главное:
\emph{реакцию не берут ребром — её порождают}. Обход собирает активированную
окрестность-подграф, а реакцию Rinness/LLM \emph{композирует поверх} материала под
настроение и мотивацию. Настроение входит \emph{дважды}: на всплывании (грустная
поднимает «потеря», игривая — «неуклюжесть») и на композиции.

\section{Консолидация реакций: по значимости, а не по приёму}

\subsection{Провизорная запись}

Сгенерированная реакция пишется провизорным ребром с низкой уверенностью (кандидат;
принцип «гипотеза, а не факт»). Ответ «это было лишнее» — отрицательное
\emph{свидетельство}, обновляющее уверенность того же ребра вниз, а не конкурирующая
консолидация. Механически ребро-реакция несёт счётчик успехов/провалов
$\mathrm{Beta}(\alpha,\beta)$:
\begin{equation}
\text{успех}:\ \alpha\!\mathrel{+}\!=\!1;\quad \text{отпор}:\ \beta\!\mathrel{+}\!=\!1;\quad
\conf=\frac{\alpha}{\alpha+\beta},\ \ \text{неопред.}=\mathrm{Var}[\mathrm{Beta}].
\end{equation}
Стабильно принимается — переходит порог и консолидируется в медленный слой;
стабильно ловит отпор — угасает. Выбор реакции по такому апостериору — это
сэмплирование Томпсона: баланс исследования новых реакций и эксплуатации
проверенных.

\subsection{Гейт значимости, а не одобрения}

\begin{pitfall}[Нельзя гейтить консолидацию приёмом]
Консолидировать реакции по тому, как их приняли, значит строить угодника, чья
личность формируется одобрением, — ровно то, что запрещает ядро («самоуважение не
течёт через понравиться»). Гейт характера внешним приёмом философски кривой для Рин.
\end{pitfall}

Правильный ход — развернуть гейт \emph{внутрь}: консолидируй по \emph{внутренней
значимости} (само-оценке при кодировании). Что такое значимость \emph{конкретно}?
Это \emph{теория оценивания} (appraisal: OCC; Шерер). Событие оценивается по
измерениям — surprise (ошибка предсказания), goal-congruence (скалярное произведение
с вектором мотивации), valence, intensity, agency, — и значимость
\begin{equation}
\text{salience}\;\propto\;\text{surprise}\times\text{goal\_relevance}\;(+\,\text{intensity}).
\label{eq:salience}
\end{equation}
Два замечания. \emph{Дешевизна}: простейшая конкретизация — $|\Delta\mood|$,
величина изменения настроения, а она уже посчитана; mood-модуль работает детектором
значимости почти без лишних вычислений. \emph{Глубина}: самый глубокий одиночный
сигнал — ошибка предсказания, и она же питает \texttt{need\_to\_discover}. Значимость
и обучение — один сигнал.

Реакцию собеседника Рин помнит и оценивает — но как \emph{просто событие}: с низкой
конгруэнтностью её ценностям «ему не зашло» обновляет модель отношения и чуть-чуть
настроение, но не доминирует в консолидации характера. Рамка оценивания сама даёт
«считаться, а можно и нет, частично».

\subsection{Условие по отношениям}

Одна шутка заходит собеседнику $A$ и падает у $B$: ребро-реакция квалифицируется
парой (ситуация $\times$ отношение), консолидируется по-персонно. Один негативный
акт даёт три эффекта: понижение ребра, апдейт настроения (осеклась),
мета/relational-пометку.

\begin{workedbox}[«Разбил кружку», полный проход]
(1)~Событие активирует окрестность: «кружка», «грусть», «неуклюжесть», плюс личная
память «$B$ хотел новую» ($-$ contradicts «грусть»). (2)~Настроение (игривое)
смещает всплывание к «неуклюжесть». (3)~Rinness композирует подкол — не берёт ребром,
собирает из ингредиентов. (4)~Ребро-реакция пишется провизорно,
$\mathrm{Beta}(1,1)$. (5)~$B$ смеётся — $\alpha\!+\!1$; или «это лишнее» —
$\beta\!+\!1$. (6)~Гейт консолидации — по $|\Delta\mood|$ и surprise, не по одобрению
$B$; но реакция $B$ обновляет модель «$B\times$такой юмор». После многих проб
подкол-на-неуклюжесть у $B$ либо консолидируется, либо угасает — комедийный тайминг
выучен как Бета-апостериор.
\end{workedbox}

\section{Дистиллят как система координат}

Возникает соблазн вовсе убрать «дистиллят» (усреднённый здравый смысл готовых баз).
Не надо. Дистиллят — \emph{общая модель мира}, относительно которой отклонение
\emph{читаемо} как характер: личность «неожиданна» лишь на фоне общих ожиданий; без
общей почвы причуды читаются как шум. Правильно: дистиллят фоном (система координат),
личность — процесс порождения отклонений на нём. Без базового «разбитая вещь обычно
плохо» не понять, почему здесь «поздравляю» — уместный сарказм.

\chapter{Голос и модель}
\label{ch:voice}

\section{Субстрат model-agnostic; личность — трансформ}

Память, настроение и граф — состояния \emph{вне} языковой модели. Значит субстрат
\emph{model-agnostic}: его можно строить на любой базовой модели и заменить базу
позже. Это де-рискует проблему потолка малой модели: когнитивная машинерия переживёт
смену базы.

Если факты и здравый смысл лежат в графе (дёшево на диске/в RAM), базовой модели не
нужно быть энциклопедией. Её роль — \emph{лингвистический процессор и голос}: часть
«Rinness» живёт в ней. Можно взять малую дообученную модель ($3$--$7$B в квантовании)
вместо тяжёлой $8$B-универсальной; факты уходят в граф, ёмкость модели
концентрируется на голосе.

\begin{pitfall}[«Меньше — строго лучше» слишком оптимистично]
Малая модель слабее в \emph{рассуждении} и связности: снимаешь фактическую нагрузку,
но не reasoning/coherence. Наблюдение из практики: рантайм prompt-Рин уже процентов
на шестьдесят состоит из анти-повторной обвязки \emph{на} $8$B — вниз с голосом
хуже. Разумный ориентир — $3$--$7$B, не $1$B, и это надо \emph{померить}, а не
постулировать.
\end{pitfall}

\section{Инструмент $\varphi$: скорер по осям характера}

Центральный инструмент — \emph{скорер характера} $\varphi$ (измерительный прибор:
строишь прибор — и он де-конфузит весь отбор). Берём $d$ осей (из архетипики:
аффективная температура, колкость, теплота, агентность, участие-против-сервиса, тип
близости, тип юмора, стиль конфликта). На каждую — \emph{рубрика} (описание $+$
якорь для оценок $0$ и $10$). Тогда $\varphi(\text{текст})\in\mathbb{R}^d$ — прогон
LLM-судьи. Отдельно задаём целевой $\mathbf{t}$. На старте — LLM-as-judge;
классификаторы по осям позже.

\section{Суперпозиция муз — это отбор в признаковом пространстве}

Каждого персонажа можно оценить по каждой оси — персонаж есть \emph{точка} в
$d$-мерном пространстве характера. Музы — существующие персонажи, свои точки.

\emph{Наивная суперпозиция} (обучение на смеси полных корпусов) даёт среднее в
\emph{центроиде} облака муз, а не в $\mathbf{t}$: $\mathbb{E}[\varphi]=\text{centroid}\ne\mathbf{t}$.
Центроид не есть Рин; интересна муза именно отклонением от центроида (центроид яркой
цундере мало отличается от безымянной героини). Усреднение вдобавок размывает
отличительные черты (большая дисперсия по осям).

\emph{Правильно} — построить выборку с $\mathbb{E}[\varphi]\approx\mathbf{t}$: не
сваливать корпуса, а \emph{отбирать} примеры (реплики, сцены), чтобы их агрегатный
центр сидел в $\mathbf{t}$. Технически — отбор данных / rejection sampling под целевой
стилевой вектор: (1)~операционализируем оси; (2)~скорим $\varphi(\text{пример})$;
(3)~отбираем/перевзвешиваем, пока $\mathbb{E}[\varphi]\approx\mathbf{t}$. Короче: не
«смешать музы» (размытый центроид), а «отфильтровать корпус к $\mathbf{t}$»
(выраженная Рин).

\section{Пайплайн по стадиям}

\begin{center}
\begin{tikzpicture}[font=\scriptsize,>=Stealth,node distance=6mm,
  b/.style={draw,rounded corners,align=center,minimum height=8mm,text width=20mm}]
  \node[b] (s1) {Стадия 1\\ сбор реплик муз};
  \node[b,right=of s1] (s2) {Стадия 2\\ отбор по $\varphi\!\to\!\mathbf{t}$};
  \node[b,right=of s2] (s3) {Стадия 3\\ LoRA: голос};
  \node[b,right=of s3] (s4) {Стадия 4\\ ситуация$\to$реакция};
  \node[b,right=of s4] (s5) {Стадия 5\\ замкнуть на граф};
  \draw[->] (s1)--(s2); \draw[->] (s2)--(s3); \draw[->] (s3)--(s4); \draw[->] (s4)--(s5);
  \draw[->] (s5.south) .. controls +(0,-0.8) and +(0,-0.8) .. (s3.south)
     node[midway,below,font=\scriptsize]{консолидация $\to$ новые данные};
\end{tikzpicture}
\end{center}

\paragraph{Стадия 1 — сбор.} Субтитры (\texttt{.srt},\texttt{.ass}), фан-скрипты,
quote-вики, сценарии VN. Извлекаем строки целевых персонажей. Узкое место —
\emph{атрибуция говорящего}. Идея кормить видео каптионеру возможна, но окупается
позже и лишь ради данных о \emph{физическом} поведении: для голоса субтитры дают
почти всё дешевле.

\paragraph{Стадия 2 — отбор (фильтрацией, не написанием).} $\varphi$ по каждой
реплике; оставляем близкие к $\mathbf{t}$. Это корпус голоса.

\paragraph{Стадия 3 — голос (LoRA).} Дообучаем малую модель на отобранных репликах —
языковое моделирование на тексте-в-стиле. Ситуации \emph{не нужны}. LoRA обучает
низкоранговую добавку $\Delta W=BA$, $\mathrm{rank}(BA)=r\ll\dim$, малую по памяти;
влезает на $8$\,ГБ для малой модели.

\paragraph{Стадия 4 — выбор реакции.} Ситуации нужны, эмбеддинги мостят несовпадение.
\emph{(A)~Синтетический} (реком.): ситуации из базы смысла $\to$ промпт сильной модели
в голосе $\to$ реакции $\to$ фильтр $\varphi$ — пары (ситуация$\to$реакция).
\emph{(B)~Майнинг+retrieval}: из сцен муз пары (ситуация, реакция), эмбеддинг ситуаций,
в рантайме ближайший сосед. Именно здесь «ситуации могут не совпадать, эмбеддинги
справятся».

\paragraph{Стадия 5 — замкнуть на граф.} Малая голос-модель $+$ граф; граф даёт
контекст (mood-смещённый retrieval), модель генерит в голосе; выходы (фильтр $\varphi$
$+$ значимость \eqref{eq:salience}) консолидируются обратно.

\section{Bootstrap: генератор $\ne$ деплой-модель}

Как получить характер-образующий граф без готовой Рин? Финальная Рин не нужна —
зачаток есть: \emph{prompt-Рин} (модульные промпты $+$ базовая модель).

\begin{enumerate}
\item строим нейтральный эмбриональный граф (глава~\ref{ch:experiment});
\item селим prompt-Рин как \emph{bootstrap-генератор}: живёт на графе, генерит
характерные реакции;
\item \emph{консолидация пишет отпечатки обратно} — граф становится character-shaped
через \emph{использование}, а не ручным reorganization-проходом;
\item на накопленном поведении дообучаем малую модель.
\end{enumerate}

\begin{pitfall}[Стабильность и потолок]
Аудит логов: prompt-Рин — \emph{Рин с перебоями}. Ядро иногда пробивает хорошо
(реплики, дословно совпадающие с установкой мотивации), но свободный чат деградирует
в липкие повторы, часть ответов нарушает явные правила (запрещённый «перформанс
трагедии», сбои рода), и \emph{guided} чище \emph{raw} — по внутреннему критерию
«ядро промпта недостаточно сильно» (потолок $8$B). Вывод: как \emph{безфильтровый}
bootstrap она шумна — вобьёшь и тики. Два хода: (i)~сидировать не в real-time,
сильной моделью (CPU-offload/квантизация или разовый API), а малую сажать на
обогащённый граф — генератор и деплой \emph{разные}; (ii)~фильтровать выходы $\varphi$.
\end{pitfall}

\section{Raw против guided и псевдо-DPO}

\emph{Raw} менее стабилен, но на пике \emph{ближе} к Рин; \emph{guided} стабилен, но
«ассистентский» (стабильно мимо цели). Для деплоя важна стабильность, для сидирования
— нет: нужны \emph{сэмплы пиковой Rinness}. Низкий потолок фильтрацией не поднять, а
нестабильность — можно (сэмплируй много, фильтруй $\varphi$). Для bootstrap: raw $+$
oversample $+$ отбор.

Отсюда дешёвый способ собрать предпочтительные пары для псевдо-DPO. DPO обучается на
тройках (промпт, выбранный $y_w$, отвергнутый $y_l$), максимизируя
\[
\log\sigma\!\Big(\beta\log\tfrac{\pi_\theta(y_w|x)}{\pi_{\text{ref}}(y_w|x)}
-\beta\log\tfrac{\pi_\theta(y_l|x)}{\pi_{\text{ref}}(y_l|x)}\Big).
\]
Подавать raw и guided как двух кандидатов на один промпт и ранжировать (частично
автоматически $\varphi$, руками — правки) — чистый способ получить \emph{реально
разных} кандидатов; раз raw иногда лучше, guided иногда — сигнал живой.

\section{Мост между базами: эмбеддинг ConceptNet}

Узлы ATOMIC — события с плейсхолдерами (\textsc{PersonX breaks a cup}), а не
концепты; «кружки» как узла там нет. Мост — извлечь контент-слова события и слинковать
к концепт-узлам ConceptNet. Но \emph{не} bag-of-words/one-hot: они не знают синонимии
(mug $\ne$ cup, разбить $\ne$ сломать), мост выйдет дырявый. Берём \emph{ConceptNet
Numberbatch} — родной эмбеддинг ConceptNet: сначала точный лемма-матч, эмбеддинг-фолбэк.
Тот же эмбеддинг обслуживает $\text{rel}(B\mid q)$ в \eqref{eq:trans} и сидирование
валентности. Одно пространство — три задачи.

\chapter{Обучение без «обучения»}
\label{ch:learning}

\section{Мы шлифуем веса, не занимаясь глубоким обучением}

Возникает верное ощущение: мы будто учим модель без «учить» и без «модели». Вроде
бы это не глубокое обучение и не графовая нейросеть, но тем не менее мы шлифуем
веса и параметры графа — по-своему. У этого ощущения есть точное содержание.

Система, которую мы строим, — параметрическая: у неё есть веса рёбер, уровни
уверенности, активации, и они обновляются опытом. Это \emph{обучение} в широком
смысле, даже без обратного распространения. Конкретно это два хорошо известных
паттерна, наложенных друг на друга:

\begin{itemize}
\item \textbf{Хеббовское/ассоциативное обучение} (семейство сетей Хопфилда;
принцип Хебба «что вместе активируется, вместе укрепляется»). Правило Хебба —
локальное усиление связи по со-активации:
\[
\Delta w_{ij}\;\propto\;a_i\,a_j ,
\]
а ассоциативная память Хопфилда хранит образы в минимумах энергии
$E=-\tfrac12\sum_{ij}w_{ij}a_i a_j$ и извлекает их спуском по $E$. Наши fusion и
консолидация, усиливающие со-активированные рёбра, — это буквально хеббовское
обучение на графе; граф памяти со spreading activation и уверенностью — современная
ассоциативная память, где извлечение есть спуск к аттрактору.
\item \textbf{Байесовское обновление уверенности} (глава~\ref{ch:propagation}):
локальные апдейты $\Delta\text{conf}$ со знаком, гейтом источника и дисконтом
независимости.
\end{itemize}

\section{Это message passing — но с локальными правилами}

Формально наша система — это \emph{передача сообщений} по графу (как в графовой
нейросети), но «параметры» (веса рёбер, уверенности) обновляются \emph{локальными}
правилами (хеббовскими, байесовскими, распадом), а не глобальным градиентным
спуском по функции потерь. Иначе говоря, это message passing \emph{без обратного
распространения}. Где мы сидим на карте методов — \emph{между} классической
ассоциативной памятью (Хопфилд) и графовыми нейросетями.

\begin{intuition}[Два способа настроить пианино]
Глубокое обучение настраивает все струны сразу по глобальному критерию («чтобы
аккорд звучал верно»), продвигаясь градиентом. Хеббовская система настраивает
каждую струну по локальному правилу («подтяни там, где чаще играют вместе»). Обе
«учат» инструмент, но первая знает про весь аккорд, а вторая — только про соседние
струны. Наша память — второго рода.
\end{intuition}

\section{Мост к «настоящей» графовой сети}

Может ли это стать графовой нейросетью в полном смысле? Да, и мост
поучителен: можно \emph{мета-обучать} сами правила обновления — учить оптимальное
правило $\Delta\text{conf}$, оптимальный распад, оптимальное смещение обхода
градиентным спуском вместо ручного дизайна. Это «обучение тому, как помнить»
(learning to remember), мета-обучение операторов памяти — реальное
исследовательское направление.

\begin{pitfall}[Но не на старте]
Мета-обучение правил меняет интерпретируемость на чёрный ящик, которому нужна
разметка и который заново вносит проблему непрозрачности субстрата. На старте —
ручные локальные правила: они интерпретируемы, отлаживаемы и не требуют обучающих
данных. «Может, и незачем» — верно для текущего этапа. Но мост реален, и полезно
знать, что стоишь ровно на нём: сегодняшняя система — предельный, «нулевой» случай
обучаемой графовой сети, где правила заданы руками, а не выучены.
\end{pitfall}

\section{Родство с DTI ещё на одном уровне}

Это же наблюдение ещё раз сшивает два проекта автора. Магистерская работа изучает
графовые нейросети с их обучаемыми операторами передачи сообщений; память Рин —
та же передача сообщений, но с локальными правилами. Если однажды правила памяти
начнут мета-обучаться, различие между «графовой моделью аффинности» и «графовой
моделью памяти» станет чисто предметным: один и тот же класс обучаемых операторов,
приложенный к разным графам. Это ещё один слой той единой спины, о которой шла речь
во введении. Существует и прямая работа на этом стыке — например, из библиографии
волта RIN: HeLa-Mem (Hebbian learning и ассоциативная память для LLM-агентов).

\chapter{Рантайм и инженерия}
\label{ch:runtime}

До сих пор мы описывали \emph{что} делают операторы. Эта глава — про то, \emph{как}
они реально запускаются, приостанавливаются, где хранятся данные и во что это
обходится по памяти и времени. Это инженерный слой, по которому предстоит работать,
поэтому здесь мы конкретны, а оценки даём даже там, где они лишь порядковые.

\section{Планировщик как конечный автомат}
\label{sec:scheduler}

Никакого «мистического решающего разума», дёргающего операторы, нет. Есть
\emph{цикл событий} с приоритетным планировщиком. Управление двухуровневое
(глава~\ref{ch:walker}): грубый \emph{режим} задаётся конечным автоматом, а внутри
режима \emph{диспетчер} раздаёт операции по бюджету $\rho$.

\begin{definition}[Единственный дефицитный ресурс]
Почти всё в системе дёшево (обход, статистика, распад, векторный поиск — миллисекунды,
раздел~\ref{sec:costs}). \emph{Дорог} только один ресурс — \textbf{слот
LLM-инференса} (один GPU, вызовы сериализуются). Поэтому планировщик по сути
распределяет \emph{слоты LLM-вызовов} между передним планом (ответ) и фоном
(fusion, метакогниция, переваривание). Всё остальное диспетчеру раздать легко.
\end{definition}

\subsection{Режимы (состояния автомата)}

\begin{center}
\begin{tabular}{p{0.16\textwidth} p{0.10\textwidth} p{0.60\textwidth}}
\toprule
\textbf{Режим} & $\rho$ & \textbf{Что делает диспетчер} \\
\midrule
\textsc{Engaged} & высок & приоритет — retrieval $+$ композиция ответа; фон приглушён/приостановлен \\
\textsc{Idle} & $\approx 0$ & фоновый разведчик (любопытство $+$ лёгкий fusion), лёгкая консолидация \\
\textsc{Sleeping} & $\approx 0$, долго & телепорт-скраб, тяжёлая консолидация, репетиция важного \\
\textsc{Reflecting} & сред. & объяснение сработавшей аномалии (метакогниция); под-режим \textsc{Idle} \\
\bottomrule
\end{tabular}
\end{center}

\subsection{Переходы}

\begin{lstlisting}[style=pseudo]
on user_message:            # экзогенное прерывание, наивысший приоритет
    freeze_background()     # заморозить текущее фоновое рассуждение (раздел 18.2)
    mode <- ENGAGED
on response_delivered and no pending_foreground:
    mode <- IDLE
on idle_timeout:            # долгий простой
    mode <- SLEEPING
on detector_fired(anomaly) and salient:
    mode <- REFLECTING      # вернётся в прежний режим по завершении
on any_state + user_message:
    preempt -> ENGAGED
\end{lstlisting}

Ключевой момент — \textbf{смещение фокуса «любопытная Рин $\to$ ответить юзеру» это
рефлекс, а не решение LLM}: сообщение пользователя — высоко-salience экзогенное
событие, оно вытесняет фон как аппаратное прерывание ($\rho$ подскочил $\Rightarrow$
фон приглушён). LLM подключается лишь на настоящих развилках суждения (бросить нить
или припарковать; вставить ассоциацию или остаться в теме) — это \emph{эскалация} к
делиберации (System~2), редкая. Рефлексы (System~1) закрывают подавляющее
большинство запусков.

\subsection{Реализация}

Простейшее — однопоточный цикл событий (события: \texttt{user\_msg},
\texttt{timer\_tick}, \texttt{detector\_fire}, \texttt{op\_complete}). Для
параллелизма — небольшой пул: поток переднего плана (ответ) и фоновый поток
(разведка), сериализуемые на общем ограничении — числе доступных LLM-слотов (так
сделано и в ReplicantLife у Toy и др.). То есть настоящий планировщик планирует
\emph{GPU-слоты}, а не «мысли».

\begin{openq}[Граница эскалации к LLM]
Какие именно решения остаются рефлексом, а какие эскалируют к LLM-делиберации, —
подвижная граница. Слишком много эскалаций $\Rightarrow$ LLM-бутылочное горло и
тормоза; слишком мало $\Rightarrow$ Рин ведёт себя механически. Настраивается на
стенде.
\end{openq}

\subsection{Произвольное внимание: внутренний forensic-запрос в Reflecting}
\label{sec:reflecting}

В таблице режимов \textsc{Reflecting} помечен как «объяснение аномалии», но это
недосказано: тут прячется механизм, которого в мануале не было явно — \emph{как Рин
думает сама с собой, без внешнего триггера}. Всё внимание до сих пор было экзогенным
(сообщение юзера как прерывание). \textsc{Reflecting} — это \emph{эндогенное,
произвольное} внимание, и оно сшивает три уже описанных куска:
\begin{enumerate}\itemsep2pt
\item \textbf{Триггер (метакогниция).} Дешёвые детекторы на кольцевом буфере телеметрии
(гл.~\ref{ch:mood}) ловят системный сбой — скачок \texttt{need\_to\_revise}, конфликт
уверенностей, затяжную руминацию.
\item \textbf{Переход (планировщик).} Автомат уходит в \textsc{Reflecting}
(\texttt{detector\_fired} выше).
\item \textbf{Механизм (forensic-ретрив).} Рин \emph{сама себе} кидает направленный
\emph{forensic}-запрос (раздел~\ref{sec:retrievaltax}) — не ассоциативный обход, а точный
подъём дерева обоснований по провенансу через query-planner — и запускает ретракцию лжи
по TMS (гл.~\ref{ch:epistemics}). Это и есть моделируемый внутренний монолог: не «что
всплыло», а «на чём это стоит и не рухнуло ли основание».
\end{enumerate}

\begin{pitfall}[Тормоза руминации обязательны]
Без ограничителей эндогенное внимание уходит в \emph{self-fulfilling-loop}: детектор
руминации сам себя кормит, Рин зависает в бесконечном самоаудите (это тот же риск, что и
самоисполняющаяся тревога, гл.~\ref{ch:mood}). Поэтому \textsc{Reflecting} жёстко
гейтируется: \emph{рейт-лимит} (не чаще, чем раз в такое-то окно), \emph{event-driven
каденция} (только на \emph{значимую} аномалию, не на каждый микроконфликт уверенностей)
и \emph{бюджет} на глубину\,/\,длительность рефлексии. Выгорел бюджет — вышли из режима,
пометив вопрос как открытый, а не крутимся дальше.
\end{pitfall}

\section{Приостановка и возобновление рассуждений}
\label{sec:suspend}

Когда сообщение пользователя вытесняет фоновое рассуждение (Рин была на середине
fusion-цепочки или любопытной прогулки), состояние надо \emph{заморозить}, чтобы потом
продолжить с того же места, а не с нуля.

\subsection{Что такое «состояние рассуждения»}

Фоновая задача — это walker (или рой) с: текущим \emph{фронтом} (множество активных
узлов $+$ их активации), \emph{затравкой/целью} (что искали), \emph{частичными
результатами} (уже собранные кандидаты, поставленные \texttt{need\_to\_discover}) и
\emph{курсором} (где остановились). Заморозка — это сериализация \emph{continuation-
токена}:

\begin{lstlisting}[style=pseudo]
Continuation = {
    task_id, goal/seed,
    frontier: [(node_id, activation), ...],   # компактно: id + float
    partial_candidates: [node_id, ...],
    pending: [need_to_discover, ...],
    cursor, priority, timestamp
}
\end{lstlisting}

Токен мал (id-шники $+$ числа) — это как захват состояния корутины при yield или
context-switch процесса в ОС (сохранить «регистры» = фронт $+$ курсор). Кладётся в
очередь приостановленных задач.

\subsection{Возобновление и устаревание}

При возврате в \textsc{Idle} задача достаётся из очереди и продолжается. Но граф мог
\emph{измениться} за время паузы (появились узлы, активации распались, что-то
спрунено), поэтому возобновление начинается с \emph{ре-валидации фронта}: выкидываем
исчезшие/распавшиеся узлы, восстанавливаем активации живых, продолжаем обход.

Не всякую задачу стоит возобновлять: пока Рин отвечала, разговор мог уйти, и нить
устарела. Поэтому у приостановленных задач своё \emph{устаревание}: если не
возобновлена за окно или вытеснена более свежей — она \emph{не удаляется бесследно}, а
роняется в граф как \emph{«недодуманная нить»} (узел с \texttt{need\_to\_discover}),
которую позже может подобрать телепорт-скраб во сне (глава~\ref{ch:sleep}). Так
прерванная мысль не теряется, а становится зацепкой для фоновой разведки.

\begin{intuition}[Замороженная мысль — это savepoint, а не труп]
Полезно думать о continuation-токене как о savepoint базы данных или о зелёном
потоке, снятом с исполнения: он хранит ровно столько, чтобы продолжить, и ничего
лишнего. А брошенная нить не исчезает — она оседает в память как «я не додумала про
$X$», и это честно по-человечески.
\end{intuition}

Каждая пара приостановка/возобновление — \emph{событие телеметрии} (глава~\ref{ch:mood}):
паттерн «меня постоянно прерывают на середине мысли» \emph{детектируется} метакогницией
как аномалия.

\section{Хранилище и индексирование}
\label{sec:storage}

Разберём предметно, из чего состоит хранилище и как оно индексируется, — и заодно
ответим, заменяет ли многоуровневое хранилище простое «окно внимания».

\subsection{Из чего состоит хранилище}

\begin{itemize}\itemsep1pt
\item \textbf{Node store} — ключ-значение \texttt{node\_id} $\to$ \{content,
embedding, метаданные\}; на диске (встраиваемое KV-хранилище: LMDB/SQLite и т.п.).
\item \textbf{Edge store} — списки смежности \texttt{node\_id} $\to$
[(target, type, weight, sign), \dots].
\item \textbf{Vector index} — индекс по эмбеддингам для поиска похожих (HNSW/FAISS),
\textbf{резидентный в RAM}; держит эмбеддинги \emph{горячего/тёплого} множества и
служит механизмом \emph{поиска якорей} по запросу (глава~\ref{ch:walker}). Важно:
\emph{поиск идёт по этому RAM-индексу, а не по диску} — на диске лежит лишь
персистентная копия эмбеддингов (backing store) и cold-тир.
\item \textbf{Структурные индексы} — по \texttt{layer}, по \texttt{importance} и т.п.,
для планировщика и распада.
\end{itemize}
Индексов, стало быть, два по сути: \emph{векторный} (семантическая близость, чтобы
найти вход) и \emph{графовая смежность} (чтобы идти от входа). Retrieval $=$ векторный
поиск якорей $\to$ обход графа от них.

\subsection{Attention-radii против тиров: это не одно и то же}

Твоя догадка про «окно/радиусы внимания» верна, но покрывает лишь половину.

\emph{Attention-radius} — это \textbf{транзиентная} граница на \emph{обход}: сколько
хопов walker расширяется от якорей за один запрос. Она держит \emph{рабочее множество}
маленьким \emph{в момент запроса}, независимо от размера графа.

\emph{Тиры} (hot/warm/cold) — это \textbf{персистентная} организация \emph{хранилища},
живущая между запросами.

Почему нужно и то, и другое: attention-radius ограничивает \emph{прогулку}, но \emph{не}
ограничивает \emph{векторный индекс якорей}. Если эмбеддинги \emph{всего} графа лежат
в одном плоском индексе, стоимость поиска якорей и RAM растут с общим числом узлов.
Тиры выносят холодные узлы \emph{из} горячего индекса (они достижимы лишь через
graph-указатели), удерживая индекс маленьким. Плюс это режет \emph{шум} ретрива: в
огромном плоском индексе косинусная близость тащит семантически близкое, но
контекстно нерелевантное (ровно жалоба Toy и др.).

\begin{center}
\begin{tabular}{p{0.22\textwidth} p{0.34\textwidth} p{0.34\textwidth}}
\toprule
 & \textbf{attention-radius} & \textbf{тиры (hot/warm/cold)} \\
\midrule
природа & транзиентная (на запрос) & персистентная (хранилище) \\
ограничивает & глубину обхода (working set) & размер индекса якорей, RAM, шум \\
когда хватает одного & если граф \emph{мал} & если граф \emph{растёт вечно} \\
\bottomrule
\end{tabular}
\end{center}

Итог: \textbf{малый граф — attention-radius достаточно, тиры не нужны; большой/растущий
— нужны оба.} Тир cold хранит огрублённый gist (глава~\ref{ch:timescales}): в warm-
индексе остаётся узел-указатель домена («знаю цепи Маркова; детали в архиве»), полная
деталь лежит в cold сжатой и \emph{подкачивается} по требованию (это и есть «чуть
дольше добывать»), как paging в ОС.

\begin{openq}[Политика тиринга]
Когда демоутить узел в cold, насколько огрублять warm-gist, когда пейджить деталь
обратно — параметры, которые надо мерить (глава~\ref{ch:experiment}). По голому диску
(раздел~\ref{sec:costs}) тиринг не обязателен; он оправдан скоростью и чистотой
ретрива, а не экономией места.
\end{openq}

\subsection{Кто к кому обращается: путь запроса и стоимости}
\label{sec:dataflow}

Ключевое уточнение, снимающее частую путаницу: \textbf{по эмбеддингам на диске никто
не ищет}. Поиск идёт по векторному индексу \emph{в RAM}; диск — это backing store и
cold-тир. Горячий путь диск почти не трогает. Ниже — кто, когда и куда обращается.

\paragraph{Путь чтения («жизнь запроса»).}
\begin{enumerate}\itemsep1pt
\item \textbf{Эмбеддинг запроса}: эмбеддер (GPU/CPU), текст $\to$ вектор. $\sim$1--10\,мс.
\item \textbf{Поиск якорей}: вектор $\to$ индекс HNSW \emph{в RAM} $\to$ top-$k$
\texttt{node\_id}. $\sim$мс (sub-10\,мс даже при $10^6$). \emph{Диска нет.}
\item \textbf{Обход}: от якорей по igraph \emph{в RAM} (ограниченный радиус,
mood-bias) $\to$ активированная окрестность. $\sim$мс--10\,мс. \emph{Диска нет.}
\item \textbf{Достать контент}: для горстки финальных узлов (5--50) текст из
node-store — RAM-кэш (горячее) или SSD (несколько KV-чтений, sub-мс, обычно уже в OS
page cache). Пренебрежимо.
\item \textbf{Композиция}: контекст $\to$ LLM. GPU-VRAM. $\sim$1--3\,с.
\textbf{Доминирует.}
\end{enumerate}
То есть всё, кроме LLM, — в RAM и $\sim$мс; диск трогается лишь чтобы достать
\emph{текст} нескольких финальных узлов (дёшево), а \emph{не} для поиска.

\paragraph{Путь записи (новый узел / fusion).} Контент+эмбеддинг+метаданные $\to$
node-store (диск, асинхронно); эмбеддинг $\to$ insert в RAM-индекс (HNSW $\sim$мс,
амортизированно); узел/рёбра $\to$ в igraph (мутации батчить, глава~\ref{ch:experiment}).
Дёшево.

\paragraph{Cold-путь (единственная «дисковая» латентность, и она редкая).} Запросу
нужна архивная деталь $\to$ в warm-индексе есть узел-указатель домена $\to$ вторая
ступень: чтение из cold-store (SSD) $+$ декомпрессия $\to$ подкачка в рабочую память.
$\sim$мс--десятки\,мс. Редко, потому что hot/warm держим в RAM.

\paragraph{Старт.} При загрузке структура графа и векторный индекс поднимаются или
mmap-ятся в RAM — одноразово, не на каждый запрос.

\begin{center}
\begin{tabular}{p{0.05\textwidth} p{0.30\textwidth} p{0.22\textwidth} p{0.28\textwidth}}
\toprule
\# & \textbf{Операция} & \textbf{Ярус} & \textbf{Стоимость} \\
\midrule
1 & эмбеддинг запроса & GPU/CPU & $\sim$1--10\,мс \\
2 & поиск якорей (HNSW) & \textbf{RAM} & $\sim$мс \\
3 & обход (igraph) & \textbf{RAM} & $\sim$мс--10\,мс \\
4 & контент финальных узлов & RAM-кэш / SSD & sub-мс, пренебр. \\
5 & композиция (LLM) & \textbf{GPU-VRAM} & \textbf{$\sim$1--3\,с (узкое место)} \\
— & подкачка cold-детали & SSD & $\sim$мс--десятки\,мс (редко) \\
\bottomrule
\end{tabular}
\end{center}

\section{Event sourcing: неизменяемый лог под изменяемым графом}
\label{sec:eventlog}

Граф памяти \emph{не должен} быть единственным хранилищем прошлого. Лучше три уровня:
\begin{itemize}\itemsep1pt
\item \textbf{Event log} — неизменяемая, только-на-дозапись история всего, что
случилось: наблюдения, коммиты, ретракции, исправления, удаления, консолидации.
\item \textbf{Memory graph} — текущая материализованная модель памяти,
\emph{перестраиваемая проекция} поверх лога.
\item \textbf{Indexes} — производные структуры для поиска (векторный, forensic).
\end{itemize}
Запись лога: \{timestamp, тип (observed/committed/retracted/consolidated/deleted/
corrected), payload, source\_context, owner\}. Важно: \texttt{source\_context} и
\texttt{owner} \emph{рождаются здесь}, в момент захвата — отсюда и берётся провенанс
(глава~\ref{ch:substrate}).

Зачем это несущее: debug («почему Рин так решила» — реплеем лога), rollback,
пересборка графа после смены схемы, \textbf{сравнение версий операторов} — поменял
алгоритм консолидации, пересобрал проекцию из \emph{того же} лога, сравнил (идеально
ложится на стенд, глава~\ref{ch:experiment}), удаление, миграции.

\begin{openq}[Рост лога и компакция]
Неизменяемый лог растёт неограниченно. На старте — хранить и не думать (разрастётся —
решим). Дальше нужна политика: периодический \emph{снапшот} проекции $+$ обрезка
старого лога, либо тиринг лога в cold (раздел~\ref{sec:storage}). Точный порог и
формат снапшота — открыты.
\end{openq}

\subsection{Снапшот-чтения: почему обходчик и коммиттер не сталкиваются}
\label{sec:snapshotreads}

Раз обходчиков много и они работают параллельно с коммиттером, напрашивается страх
гонок: RW читает узел, который коммиттер как раз меняет — не получит ли он рваное,
полу-обновлённое состояние? При нашей раскладке ролей (walker читает и предлагает,
структурно пишет только сериализованный коммиттер, гл.~\ref{ch:walker},
раздел~\ref{sec:proposecommit}) страх почти растворяется, и event sourcing закрывает
остаток.

Ключ в том, что «изменение» узла у нас — \emph{не} правка на месте, а \emph{новая
версия}: коммиттер дописывает событие в лог и обновляет проекцию, старая версия узла
никуда не пропадает мгновенно. Значит обходчики читают \emph{консистентный снапшот}
проекции (MVCC-подобно): RW2, зашедший в «меняющийся» узел, видит его до-коммитную
версию целиком — рваных чтений нет, потому что нет правки на месте, которую можно
застать на середине. Read $\leftrightarrow$ read не конфликтует никогда; read
$\leftrightarrow$ write решается версионностью; write $\leftrightarrow$ write не
существует по построению (писатель один). Остаётся лишь мягкое побочное состояние —
рабочая яркость $a_{\text{work}}$ и телеметрия внимания; оно \emph{коммутативно} (для
яркости — через максимум, не сумму, гл.~\ref{ch:walker}), поэтому порядок конкурентных
подсветок неважен и блокировки не нужны.

\begin{intuition}[Гонку не арбитрируют — её конструктивно исключают]
Мы не пишем сложный протокол блокировок, кто кого пропускает. Мы устроили систему так,
что сталкиваться \emph{нечему}: один структурный писатель, версионные чтения для
всех остальных, коммутативная телеметрия. Это дешевле и надёжнее любого арбитража
постфактум — та же философия, что «противоречие лучше не разрешать задним числом, а
не допускать молчаливой перезаписи» (гл.~\ref{ch:sleep}).
\end{intuition}

\subsection{Устаревшие предложения: коммиттер валидирует, а не просто пишет}
\label{sec:staleproposals}

Снапшот-чтения убирают гонку на \emph{чтении}, но остаётся тонкий изъян на
\emph{записи}. Обходчик мог посчитать предложение на версии графа
\texttt{snapshot=123}, а к моменту коммита проекция уже на \texttt{129}, где часть
посылок успела устареть: узел-родитель ретрактнут, слит, понижен в уверенности,
изменил \texttt{layer} или \texttt{visibility}. Если коммиттер применит предложение
вслепую, гонку по данным мы убрали, но получили \emph{семантически устаревший коммит} —
что не лучше.

Поэтому предложение — не голая операция, а операция с \emph{контекстом валидности}:
\{\texttt{op}, \texttt{members}, \texttt{read\_snapshot\_id}, \texttt{parents},
\texttt{proposal\_ttl}, \texttt{validity\_conditions}\}. Коммиттер перед применением
\emph{валидирует}, а не сразу пишет: живы ли ещё родители? не deprecated? не изменились
ли их \texttt{confidence}\,/\,\texttt{layer}\,/\,\texttt{visibility} с момента
\texttt{read\_snapshot\_id}? не вышел ли \texttt{proposal\_ttl}? По итогу — одно из
четырёх: \emph{применить} (условия держатся), \emph{rebase} (пересчитать на текущей
версии, если посылки лишь сдвинулись), \emph{downgrade} (принять как более слабую
гипотезу) или \emph{reject} (посылки развалились). Это оптимистический контроль
конкуренции: не блокируем заранее, а проверяем на коммите. Без этого read/write-гонку
уберёшь технически, а память тихо наполнится решениями, принятыми по уже неверным
предпосылкам.

\subsection{«Забудь это»: удаление при неизменяемом логе}
\label{sec:forget}

Неизменяемый лог порождает очевидный конфликт: а что происходит, когда пользователь
говорит «Рин, забудь последние 10 минут»? Для локального компаньона это не юридическая
занудность, а конкретный будущий UX, и он должен что-то значить.

Ключ к разрешению: \emph{неизменяемость $\ne$ «нельзя удалять»}. Она означает
\emph{append-only, без тихих правок на месте}; само удаление — это \emph{тоже
appended, аудируемое событие}, а не нарушение дизайна. «Забыть» — первоклассная,
поддержанная операция, и реализуется тремя механизмами, которые работают вместе:

\begin{itemize}\itemsep2pt
\item \textbf{Граф (легко).} Убираем узел, ретрактим, перепропагируем. Каскад: удаление
исходного наблюдения понижает/сносит производные узлы, чьим \emph{провенансом} оно
было (ретракция по TMS, глава~\ref{ch:epistemics}). «Забудь X» — это ещё и «забудь
то, что из X выведено».
\item \textbf{Tombstone.} В лог дописывается delete-событие; при пересборке проекции
удалённое пропускается. Эффект немедленный, но сырое событие ещё \emph{в логе}
(годится для undo/аудита, но это не настоящее стирание).
\item \textbf{Crypto-shredding.} События (по сессии/скоупу) зашифрованы своим ключом;
«удалить» $=$ уничтожить ключ $\to$ шифртекст остаётся, но навсегда нечитаем.
\emph{Мгновенное, проверяемое} стирание без переписывания лога — идеально под «забудь
последние 10 минут» (эти минуты под эфемерным сессионным ключом; шреднул ключ —
стёр).
\item \textbf{Compaction.} Фоновая пересборка лога выбрасывает tombstoned-события —
сырое физически исчезает.
\end{itemize}

То есть «Рин, забудь X» $=$ tombstone $+$ каскадная ретракция производных $+$
шреддинг ключа (истинное стирание) $+$ компакция-уборка. Конкретно, а не «ну она как-
то забудет».

\begin{pitfall}[Crypto-shred рождает key management и требует taint-propagation]
Две дыры, без которых схема красива, но дырява. \emph{Первая}: шреддинг создаёт новый
центр тяжести — \textbf{управление ключами} (где ключи, кто их держит, гранулярность
сессия/скоуп). \emph{Вторая, важнее}: уничтожив ключ raw-события, ты \emph{не} стёр
след, если plaintext успел растечься — в debug-log, снапшот context-packer,
embedding-кэш, вектор-индекс, бэкап, telemetry-сводку или производный fusion-узел.
Поэтому рядом с «забудь» нужна дисциплина \textbf{taint-propagation}: всякий
производный объект наследует taint/провенанс от породившего события, а delete-cascade
чистит \emph{все} производные хранилища — граф, индексы, сводки, эмбеддинги, снапшоты,
pending-очереди, бэкапы. Иначе «забыл» будет только на бумаге.
\end{pitfall}

\begin{openq}[Гранулярность забвения]
На каком уровне держать ключи шифрования (событие / сессия / скоуп), что именно
каскадит при удалении (только прямые производные или транзитивно), и как показывать
пользователю, что забыто, — открыто. Но механика («забыть» = событие в логе $+$
shred $+$ каскад) зафиксирована.
\end{openq}

\section{Write-path: что попадает в память, а что нет}
\label{sec:writepath}

Отдельно проговорим, \emph{что именно} записывается, потому что «не писать» —
полноправный исход, а не сбой. Конвейер:
\begin{lstlisting}[style=pseudo]
raw_event -> parse/extract candidates
          -> [ epistemic_gate, salience_gate, anti_fusion_gate ]
          -> write | skip
          -> (if write) propagate
          -> audit_log   # и write, и skip логируются
\end{lstlisting}
\begin{itemize}\itemsep1pt
\item \textbf{epistemic-гейт}: стоит ли этому верить и с какой уверенностью (гл.~\ref{ch:epistemics}).
\item \textbf{salience-гейт}: значимо ли это для Рин (appraisal, гл.~\ref{ch:personality}).
\item \textbf{anti-fusion-гейт}: см.\ ниже.
\end{itemize}

\begin{principle}[response-useful association $\ne$ memory-worthy fusion]
Оператор должен уметь не только «нашёл красивую связь $\to$ породил узел», но и
явно \emph{не писать}: «со-активация есть, связного вывода нет», «связь эстетическая,
не эпистемическая», «это каламбур, не убеждение», «полезно для реплики, но не для
памяти». Рин может \emph{использовать} ассоциацию в ответе, не записывая её узлом —
иначе память зарастает «интересными мыслями», которые были хороши в моменте, но не
являются знанием.
\end{principle}

И \texttt{skip} обязательно логируется (хотя бы в debug): иначе не поймёшь, почему
Рин «не запомнила» важное.

\subsection{Коммит не бинарен: полосы, а не «только ночью»}
\label{sec:commitlanes}

Соблазнительно сказать: «walker'ы читают весь день, а пишет коммиттер только ночью».
Это \emph{слишком сильно} и ломает сознательное мышление. Если так сделать буквально,
Рин днём станет странной: ходит по графу, находит связи, что-то «понимает» — но не может
ни закрепить свежий внешний факт, ни воспользоваться дневной догадкой иначе как
эфемерным текстом. Правильный инвариант:

\begin{principle}[Коммит не бинарен]
Наблюдаемые события и рабочие гипотезы фиксируются \emph{немедленно}; долговременная
консолидация \emph{отложена}. Коммиттер один и сериализован, но работает не «ночью
против дня», а в нескольких \emph{полосах} по стоимости и durability.
\end{principle}

Полосы одного коммиттера:
\begin{itemize}\itemsep2pt
\item \textbf{immediate.} Пишется всегда и сразу: append в event-log (что произошло —
это \emph{журнал}, а не «убеждение»), прямые пользовательские наблюдения,
task-critical внешние факты, явные исправления, важные boundary\,/\,preference.
Откладывать это до ночи нельзя — потеряешь след происходящего.
\item \textbf{foreground pending.} Дневная рабочая память — это наш \emph{session-слой}
(гл.~\ref{ch:timescales}), не новая сущность: временные узлы-гипотезы, найденные
факты, цепочки рассуждения, кандидаты на консолидацию, \texttt{need\_to\_verify}. Живут
с TTL, низкой уверенностью, провенансом и статусом \texttt{pending} — уже объект,
которым можно \emph{пользоваться сейчас}, но ещё не вечная память.
\item \textbf{idle.} Микробатчи в простое: мелкий dedup, малые правки уверенности,
локальные проверки противоречий.
\item \textbf{sleep.} Тяжёлое: консолидация $A,B,C,D\to X$, массовое разрешение
конфликтов, fusion далёких ассоциаций, pruning, reindex, пересборка summary, ревизия
низкоуверенных потомков, promotion рекуррентных кандидатов.
\end{itemize}
Это и есть fast-path\,/\,slow-path: дешёвые $O(1)$-мутации (добавить узел-гипотезу,
инкремент) идут днём по быстрой полосе; дорогие $O(|V|\log|V|)$ и глобальные — ночью.
Критично: быстрая полоса — про \emph{латентность}, а не про поблажку. Она проходит тот
же write-path и те же гейты (epistemic\,/\,salience\,/\,anti-fusion,
раздел~\ref{sec:writepath}), просто синхронно.

Сводно, по трём процессам (и подчеркнём: строка записи \emph{не} ночецентрична — днём
работает immediate-lane):

\begin{center}
\begin{tabular}{p{0.19\textwidth} p{0.36\textwidth} p{0.36\textwidth}}
\toprule
\textbf{Процесс} & \textbf{Днём ($\rho\!\approx\!1$)} & \textbf{Фон\,/\,ночь ($\rho\!\approx\!0$)} \\
\midrule
Retrieval & локальный разреженный matvec-диффузия из HNSW-якорей (fast-path) & телепорт-скраб, эргодическое покрытие, большой $\alpha$ \\
Fusion\,/\,мышление & «мотыльки»: одиночные RW на остывающем $a_{\text{work}}$ $\to$ кандидаты в thought-inbox & честный рой (slow-path): треки, $n$-арные цепочки, далёкие мосты \\
Запись & \textbf{immediate-lane}: append в event-log, безопасные foreground-микрокоммиты (\texttt{observed}, \texttt{boundary}, правки \texttt{preference}); \texttt{pending} догадки — в thought-inbox & коммиттер (slow-path): судит гипотезы, тяжёлая консолидация, prune, reindex \\
\bottomrule
\end{tabular}
\end{center}

Так снимаются оба дневных противоречия. \emph{Сознательный поиск} не ждёт ночи: Рин
«гуглит здесь и сейчас», кладёт результат в event-log и session-граф как sourced
observation, использует в ответе — но со статусом «свежий источник, не консолидировано»,
а не «я давно это знаю». \emph{Догадки} тоже случаются днём — просто живут как активная
гипотеза, а взрослеют в durable-память позже.

\subsection{Thought-inbox: как дневная догадка попадает в реплику, не став фактом}
\label{sec:thoughtinbox}

Фоновый разведчик работает \emph{всегда} — даже в напряжённом диалоге, на остаточном
бюджете $B_{\text{bg}}=B_{\text{total}}(1-\rho)$ (гл.~\ref{ch:walker}), который никогда
не ноль. Поэтому «момент озарения» (связь House-Уилсон и Cast Away-Уилсон) может
родиться и днём. Но показывать Рин \emph{сырой след} walker'ов нельзя — это шум: сотни
посещённых узлов, случайные касания, мусорные мосты. Между следом и репликой стоит
фильтр — \emph{thought-inbox}: туда попадают не все касания, а только отобранные
предложения-инсайты (сильная, но неочевидная со-активация далёких узлов):
\begin{lstlisting}[style=pseudo]
type: insight_candidate
content: "Уилсон у House и Уилсон в Cast Away могут делить нарративную функцию"
status: unverified_hypothesis
salience: 0.72
provenance: [node_a, node_b, path_summary]
source: background_walker
ttl: 2h
recommended_action: mention_as_hypothesis | verify_later | ignore
\end{lstlisting}
Context-packer иногда подмешивает самые «горячие» из них в промпт — но \emph{с явной
эпистемической меткой} (это ровно наш packer-слой меток, раздел~\ref{sec:packer}):
«\emph{возможная мысль, не факт: между A и B может быть связь; статус — не проверено;
использовать как догадку, не как установленную память}». Дальше цикл замыкается сам:
если контекст диалога позволяет, Рин озвучивает догадку — и в \emph{этот} момент она
легализуется, падает в event-log и становится узлом со статусом
\texttt{need\_to\_verify}; если не пригодилась — тихо затухает по TTL и убирается
сборщиком. Изоляция ролей сохранена: walker \emph{читает и предлагает}, LLM
\emph{осмысляет и озвучивает}, коммиттер \emph{пишет}.

\begin{pitfall}[Не писать активацию на каждый шаг walker'а]
«Walker дёшев» — верно относительно LLM, но не бесплатно, и узкое место обычно не в
математике, а в обвязке (индексы, сериализация, локи, provenance-traversal,
scope-фильтры, гл.~\ref{ch:experiment}). Поэтому обходчик \emph{не} пишет активацию в
хранилище на каждом шаге: он копит visits\,/\,деньты активации локально в памяти
процесса и \emph{батчем} сбрасывает в transient-activation store (по-прежнему max, не
сумма). И на сбросе различает \texttt{access\_reason}: \texttt{user\_recall} укрепляет
след $B_{\text{base}}$, а \texttt{scrub\_visit} — нет.
\end{pitfall}

\section{Таксономия ретрива: не всё — блуждание}
\label{sec:retrievaltax}

Walker как «почти всё» — переобобщение. Не всякая операция памяти естественно
является ассоциативным блужданием:
\begin{itemize}\itemsep1pt
\item \textbf{associative} (якоря $+$ walkers) — для живости и ассоциаций (гл.~\ref{ch:walker}).
\item \textbf{forensic} — точный поиск по метаданным/провенансу/времени: «покажи всё
про $X$», «узлы с источником $Y$», «что противоречит этому факту», «последняя версия
prompt-модуля». Это индексы и query-planner, \emph{не} блуждание.
\item \textbf{creative} — walkers $+$ fusion $+$ сон (гл.~\ref{ch:sleep}).
\end{itemize}
Forensic-ретрив — не приятное дополнение, а \emph{необходимость} для обслуживания:
аудит, миграции, точное удаление конкретных узлов. Walker полноты не гарантирует, а
скучный индекс — гарантирует. Для живости Рин нужны walkers; для корректности памяти
нужны скучные индексы.

\section{Context-packer: как граф превращается в промпт}
\label{sec:packer}

Это интерфейс графа и LLM на инференсе — что именно \emph{видит} модель, и
\textbf{последний эпистемический гейт перед речью}. Если packer отдаёт модели плоский
список содержаний без статусов, вся дисциплина графа (уверенность, гипотезы,
депрецирование) теряется на последнем метре. Поэтому в контекст попадают \emph{не
только contents, но и метки}: \texttt{observed}, \texttt{hypothesis},
\texttt{deprecated}, \texttt{contested}, \texttt{low\_conf}, \texttt{time\_validity},
и при необходимости \texttt{do\_not\_state\_as\_fact}. Плюс провенанс для важных
утверждений и mood/persona-обусловливание. Тогда модель знает, что подать как факт,
что — как «похоже», а что не выдавать за истину вовсе. Держать компактно: короткий
контекст $=$ маленький KV-кэш $=$ меньше VRAM (раздел~\ref{sec:costs}).

Это \texttt{prompt context} — \emph{временная сборка} (граница из
главы~\ref{ch:architecture}), а не хранилище: спаковали, отдали модели, выбросили.

Здесь же — естественное место для \emph{eligibility-гейта}: перед упаковкой отсеять
узлы, которые \emph{есть в памяти, но не годятся для этого контекста} (например, не
класть в промпт то, что не стоит показывать при данном собеседнике или выносить в
публичный контент). Это шире удаления и \emph{у́же} тяжёлой governance-бюрократии:
не «кто владеет и на каких условиях», а простой guardrail
\emph{retrieval $\to$ eligibility-гейт $\to$ packer $\to$ output}. Тяжёлую систему
доступа/согласия мы сознательно не строим (вектор проекта сместился), но сам крючок
стоит держать в виду.

\begin{openq}[Eligibility как отложенный guardrail]
Полноценный слой видимости/раскрытия (что можно показать здесь, класть в промпт,
использовать в публичном контенте, обобщать) — отложен. Пока — минимальный гейт
«память может существовать, но быть негодной для данного контекста», без бюрократии.
Разворачивать ли его — решается, когда вернётся потребность в family/публичном режиме.
\end{openq}

\begin{openq}[Формат упаковки]
Как именно показывать модели статус узла (уверенность, гипотеза vs факт, провенанс),
сколько узлов включать и как их ранжировать под текущий запрос — открыто, меряется на
стенде.
\end{openq}

\section{Оценки стоимости: память, время, узкие места}
\label{sec:costs}

Оценки порядковые (для локального железа класса RTX 4090 / 8--24\,ГБ VRAM),
но они сразу показывают, где узкие места, а где нет.

\subsection{Стоимость операций}

\begin{center}
\begin{tabular}{p{0.30\textwidth} p{0.14\textwidth} p{0.46\textwidth}}
\toprule
\textbf{Операция} & \textbf{Порядок} & \textbf{Комментарий} \\
\midrule
LLM-инференс (композиция, объяснение) & \textbf{0,5--5 с} & \textbf{узкое место}; $\sim$2\,с/запрос на 7--8B локально \\
векторный поиск якорей (HNSW, $10^6$) & $\sim$мс & не узкое место \\
обход/spreading (сотни-тысячи узлов) & мс--10\,мс & C-ядро igraph, дёшево \\
пропагация уверенности (лавина) & $\sim$мс & дёшево \\
детекторы телеметрии & O(1) инкрем. & дёшево \\
распад/ACT-R & дёшево \emph{лениво} & полный проход по $N$ узлам — дорого; делать лениво (при доступе) или во сне, не на каждый тик \\
переваривание учебника & \textbf{минуты-часы} & одноразово; дорого из-за LLM-извлечения фактов; только в \textsc{Idle}/\textsc{Sleeping} \\
\bottomrule
\end{tabular}
\end{center}

\subsection{Латентность ответа}

Критический путь: \texttt{user\_msg} $\to$ поиск якорей (мс) $\to$ обход (мс) $\to$
композиция ответа (\textbf{LLM, $\sim$1--3\,с}) $\to$ ответ. \textbf{Бутылочное горло —
LLM-вызов}; всё графовое к латентности почти не добавляет. Цель — держать ответ в
$\sim$1--3\,с, то есть минимизировать \emph{число} LLM-вызовов на переднем плане (в
идеале один: граф уже подал контекст, модель только оформляет).

\subsection{Память}

\begin{itemize}\itemsep1pt
\item \textbf{VRAM} — модель (фиксировано, узкое место по размеру модели: 7--8B в
квантизации $\sim$5--8\,ГБ). Отсюда весь смысл «факты в граф, а не в веса».
\item \textbf{RAM} — векторный индекс горячего множества; для $10^6$ эмбеддингов
плоско $\sim$1,5\,ГБ, с HNSW/квантизацией меньше; \emph{тиринг} держит это ограниченным.
\item \textbf{Диск} — сам граф: $\sim$2--3\,КБ/узел, $10^5$--$10^6$ узлов $=$
0,3--3\,ГБ (глава~\ref{ch:experiment} для оценки роста). Дёшево; учебник $\approx 10^4$
узлов $\approx$ 30\,МБ.
\end{itemize}

\subsection{Где узкие места, а где нет}

\begin{principle}[Всё сводится к LLM-слотам]
Дорог только GPU/LLM — и по латентности ответа, и по фоновой пропускной способности.
Всё графовое (хранилище, обход, статистика, распад, векторный поиск) — дёшево (мс) и
узким местом \emph{не} является. Поэтому вся инженерная задача рантайма сводится к
одному: \emph{как распределить дефицитные LLM-вызовы} между быстрым ответом (передний
план) и фоном (fusion/метакогниция/переваривание). Это и делает планировщик
(раздел~\ref{sec:scheduler}).
\end{principle}

\begin{openq}[Цели по латентности и throughput]
Конкретные цели (ответ $<$ N\,с; сколько фоновых fusion в минуту в простое; когда
переключать модель на меньшую ради скорости) — под стенд и под конкретное железо.
Динамический роутинг «мелкая модель на рутину, крупная на метакогницию» (как у Toy и
др.) — очевидный кандидат оптимизации.
\end{openq}

\chapter{Эксперимент: как это проверять}
\label{ch:experiment}

\section{Эрзац-граф лучше живой Рин}
\label{sec:groundtruth}

Тестировать архитектуру памяти на готовом агенте — долго и слепо: сначала собери всё,
а потом окажется, что с памятью куча неоптимальностей, найденных \emph{после} финальной
сборки. Нужен управляемый стенд — эрзац-граф с retrieval.

\begin{principle}[Ground truth важнее реализма]
Суррогатный граф — не бедный заменитель, а для отладки механик \emph{лучше} живой Рин,
потому что синтетика даёт \emph{ground truth}. С настоящей Рин ты не знаешь, что она
«должна» помнить, значит не измеришь, правильно ли отработали fusion, скраб, пропагация.
С подсаженной синтетикой истину задаёшь ты — и меряешь точность/полноту починки,
сходимость уверенности, нашёл ли скраб подсаженную ложь.
\end{principle}

Это тот же принцип, что контролируемая оценка в DTI: управляемый синтетический тест
изолирует эффект \emph{до} грязных реальных данных. Слабые места находятся \emph{до}
финальной сборки.

\section{Три яруса}

\paragraph{Ярус 1 — синтетическая топология (начинать отсюда).} Графы с известной
структурой, чтобы юнит-тестить и диагностику, и механики против теории: стохастическая
блочная модель (SBM — сообщества с регулируемым мостом, для over-squashing/проводимости),
Барабаши–Альберт (scale-free хабы — для «гравитационных ям» и концентрации $\pi$),
Уоттс–Строгац (small-world), деревья против экспандеров (крайности кривизны). Семантики
ноль, зато полная контролируемость: подсаживаешь ложное-правдоподобное убеждение и
противоречащее свидетельство в угол с малым $\pi$ и меряешь, находит ли телепорт-скраб.

\paragraph{Ярус 2 — семантический граф из готовых баз.} Чтобы у fusion/retrieval был
смысл. Разбор баз — раздел~\ref{sec:datasets}.

\paragraph{Ярус 3 — LLM-генерённый «поток жизни».} Полная петля с реалистичным
содержанием и контролируемой истиной: языковая модель генерит синтетическую
автобиографию/эпизоды/диалоги фейк-персоны с подсаженными фактами, противоречиями,
эволюцией убеждений. Прямой прообраз — \emph{Generative Agents} (Park и др.): поток
памяти $+$ retrieval (свежесть$\times$важность$\times$релевантность) $+$ рефлексия (это
и есть fusion). «Ползучую веру» подсаживаешь вручную и меряешь, лечится ли.

\section{Разбор баз знаний}
\label{sec:datasets}

\begin{center}
\begin{tabular}{p{0.16\textwidth} p{0.30\textwidth} p{0.42\textwidth}}
\toprule
\textbf{База} & \textbf{Что даёт} & \textbf{Под какую механику} \\
\midrule
ConceptNet & ассоциативно-таксономические связи, взвешенные & spreading activation,
ассоциативный fusion; веса $\approx$ мягкая уверенность \\
ATOMIC/10x & причинно-социальный смысл событий (намерение/реакция/эффект) &
эпизодика, сцепка с настроением, fusion-как-вывод \\
Wikidata & факты $+$ схема (время, references, ранг) & \emph{схему} красть под
темпоральную валидность, провенанс, ревизию; содержание — нет \\
\bottomrule
\end{tabular}
\end{center}

Не «какой из них», а «какой под какую механику». На старте — ConceptNet (чистый граф
под первую механику); Wikidata-\emph{схему} — когда дойдёшь до темпоральных слоёв и
ревизии. Порознь одного мало: ConceptNet без ATOMIC — мир без психологии, ATOMIC без
ConceptNet — реакции без предметного мира.

\section{Эмбриональный граф и работа стыка}

Объединить ConceptNet и ATOMIC в \emph{эмбриональный} граф и доращивать фьюжнами —
канонично: базовая когнитивная матрица, в которую новые факты встраиваются рёбрами.

\begin{pitfall}[«Форматы отличаются только на диске» — недооценка]
Внутри графовой базы это одна структура, и математика обхода от типа ребра не меняется.
Но настоящая работа не в формате, а в \emph{интерфейсе}: узлы ConceptNet — концепты,
узлы ATOMIC — события с плейсхолдерами \textsc{PersonX/Y}. Связать «разбить кружку» с
концептом «кружка» и заземлить \textsc{PersonX/Y} на реальных собеседников — это
entity/event linking (через Numberbatch, глава~\ref{ch:voice}), нетривиальный шаг, а
не объединение таблиц.
\end{pitfall}

\section{Батарея экспериментов}

Каждый тест — одна механика плюс метрика; всё гоняется уже на Ярусе 1.

\begin{center}
\begin{tabular}{p{0.05\textwidth} p{0.44\textwidth} p{0.42\textwidth}}
\toprule
\# & \textbf{Гипотеза} & \textbf{Установка и метрика} \\
\midrule
1 & $\pi$ концентрируется на хабах & BA-граф; распределение посещений против степени \\
2 & узкое место душит cross-cluster fusion & два SBM-сообщества; частота fusion против
ширины моста $m$ \\
3 & телепорт-скраб находит ложь в углу с малым $\pi$ & подсадка; precision/recall и
время-до-нахождения против $\alpha$ \\
4 & уверенность пересглаживается за много циклов & дисперсия $\conf$ во времени; тест
якоря-остатка \\
5 & руминация и её гашение & mood-congruent $+$ инерция; автокорреляция $\mood$ с
setpoint и без \\
\bottomrule
\end{tabular}
\end{center}

Скелет стенда (build $\to$ plant $\to$ operator $\to$ metric) — в
приложении~\ref{ch:algorithms}.

\section{Микроэксперименты: наблюдать физику потока}

Помимо «прошёл/не прошёл»-тестов есть класс более тонких наблюдений, вырастающих
прямо из главы~\ref{ch:topology}. Их ценность не в метрике «лучше/хуже», а в том,
что они делают \emph{видимой} скрытую физику графа.

\subsection{Траектории как проявитель реальной топологии}

Каждый обход оставляет след — кусочную кривую по узлам. Дневной retrieval даёт
короткие целевые дуги; сон — длинные телепорт-разорванные фрагменты. По отдельности
это шум. Но \emph{на большой статистике} следов их наложение даёт эмпирическое
\emph{распределение реальных маршрутов} по графу. И это распределение — прямой
проявитель топологии: где маршруты сгущаются, там хабы и «гравитационные ямы»; где
кривые почти не проходят — затенённые углы с малым $\pi$; где потоки протискиваются
через узкий пучок рёбер — там боттлнек и over-squashing.

\begin{principle}[Наблюдаемое $\ne$ предопределённое]
Ключевой смысл эксперимента: эмпирическое распределение переходов, снятое со следов,
будет \emph{отличаться} от «предопределённых» вероятностей, которые мы заложили в
методы вроде PageRank. Мы задаём переходы \emph{локально}, а наблюдаем \emph{глобальный}
эффект: стационарное $\pi$ — свойство всей топологии, а не наших локальных правил
(глава~\ref{ch:topology}). Расхождение между «встроенным приором» и «реально
пройденным» и есть отпечаток топологии, который иначе не увидеть.
\end{principle}

\emph{Установка.} На Ярусе~1 (SBM или BA) запускаем рой, копим следы, строим
гистограмму посещений узлов и эмпирическую матрицу переходов $\hat P$. Сравниваем:
(а)~гистограмму посещений с распределением степеней (проверка «гравитационных ям»);
(б)~$\hat P$ с заданным локальным правилом и с PageRank-приором (проверка тезиса
«наблюдаемое $\ne$ предопределённое»); (в)~плотность прохождения через мост-боттлнек с
теоретическим $\Reff$. \emph{Метрика} — расхождение распределений (например,
KL-дивергенция) и корреляция плотности следов с центральностью/$\Reff$.

\subsection{Короткие дуги против длинных треков}

Полезный контроль из главы~\ref{ch:topology}: короткие блуждания читают
\emph{локальную} геометрию (кривизну окрестности), длинные — глобальную $\pi$.
Нарезая следы на дуги разной длины и сравнивая их статистику, можно эмпирически
разделить локальный и глобальный вклад топологии. Телепорт при этом виден как
«разрывы» кривых, регуляризующие распределение (сдвигающие его к более равномерному).

\subsection{Скраб как измеритель покрытия}

Отдельный микроэксперимент — проверить \emph{эргодичность} телепорта напрямую:
подсадить метки в узлы с разным $\pi$ (в т.ч.\ в затенённые углы) и измерить время до
первого посещения каждой метки роем сна как функцию $\alpha$. Ожидание: без телепорта
($\alpha=0$) дальние-по-топологии метки не находятся вовсе; с ростом $\alpha$ время
покрытия падает и становится конечным для всех узлов. Это эмпирическое подтверждение
гарантии покрытия (глава~\ref{ch:sleep}).

\section{Инструмент}

\emph{igraph} (C-ядро, на порядки быстрее чистого Python-графа) — разумный выбор.
Атрибуты first-class и векторные (обогащать \texttt{confidence}, \texttt{valence},
\texttt{time} удобно). Есть SBM, Барабаши, Уоттс–Строгац, лапласиан,
персонализированный PageRank (телепорт), случайные блуждания, центральности. Две
осторожности: \emph{(i)}~эффективное сопротивление одной строкой нет — считаешь из
$L^{+}$ (диагональная формула), пара строк; \emph{(ii)}~igraph силён на \emph{статике},
а fusion/самоэволюция — \emph{частые мелкие мутации}, где он слабее (переиндексация при
добавлении узлов). Мутации батчить или держать «живой» граф в мутабельной структуре и
периодически снапшотить в igraph под тяжёлую спектральную диагностику.

\section{Батарея отказов памяти}
\label{sec:failuremodes}

Стенд нужен не только чтобы проверять, что механики \emph{работают}, но и чтобы
ловить, как они \emph{ломаются}. Память компаньона имеет свои патологии; их надо
уметь воспроизводить и мерить, а не ждать, пока всплывут на живой Рин. Ниже —
основные отказы и как их детектировать.

\begin{center}
\begin{tabular}{p{0.24\textwidth} p{0.34\textwidth} p{0.34\textwidth}}
\toprule
\textbf{Отказ} & \textbf{Что это} & \textbf{Как ловить} \\
\midrule
ложная память (false memory) & fusion родил гипотезу, она закрепилась как факт & подсадить недоопределённую конфигурацию; доля гипотез, ушедших в высокую уверенность без корроборации \\
мусорный fusion & память зарастает «интересными мыслями» & precision принятых fusion после отложенной ревизии (сколько реально полезно) \\
ползучая вера & посылка опровергнута, потомки живут & опровергнуть узел; доля потомков, продолжающих влиять спустя $N$ циклов \\
залежавшееся убеждение (stale) & мир изменился, узел нет & темпоральная подсадка; задержка до понижения устаревшего \\
переконсолидация & детали схлопнулись слишком рано/сильно & сравнить downstream-качество до/после консолидации (ошибка коммутации, гл.~\ref{ch:timescales}) \\
mood-biased recall & настроение искажает припоминание сверх меры & один запрос под разными $\mood$; расхождение окрестностей retrieval \\
руминация & петля настроения расходится & автокорреляция $\mood$ с/без медленных слоёв (гл.~\ref{ch:mood}) \\
затенённые углы & важное не находится из-за топологии & подсадка в регион с малым $\pi$; время до нахождения телепорт-скрабом \\
\bottomrule
\end{tabular}
\end{center}

Каждый отказ — это \emph{провал}, который стенд должен уметь \emph{вызвать по
команде}: подсадил сценарий $\to$ прогнал оператор $\to$ измерил, случился ли отказ и
починился ли он. Это и есть смысл ground truth (раздел~\ref{sec:groundtruth}).

\section{Таблица: утверждение / статус / как проверить}
\label{sec:claimtable}

Мануал держится на связке рабочих гипотез. Чтобы он был \emph{фальсифицируемым}, а не
красивым, соберём ключевые утверждения с их статусом (принцип-дизайна /
гипотеза / известная математика / \emph{эмпирически подтверждено}) и конкретным способом
проверки. Это — рабочий план стенда, а не украшение. \emph{Разделение источников}: мануал
держит \emph{утверждения} и способ проверки (вне времени), а детальные прогоны — таблицы,
по-ранные числа, setup — живут в репозитории (\texttt{experiments/} $+$ \texttt{reports/})
и версионируются с кодом; отдельной главы «наши эксперименты» в монографии нет намеренно,
чтобы текст не дублировал репо и не протухал. Указатели \texttt{reports/experiment\_N} в
колонке «как проверить» ведут к доказательству.

\begin{center}
\begin{tabular}{p{0.34\textwidth} p{0.12\textwidth} p{0.46\textwidth}}
\toprule
\textbf{Утверждение} & \textbf{Статус} & \textbf{Как проверить} \\
\midrule
телепорт улучшает покрытие скраба & гипотеза & подсадка противоречия в регион с малым $\pi$; время-до-нахождения vs $\alpha$ \\
mood-congruent recall поддерживает непрерывность персоны & гипотеза & A/B на окрестностях retrieval под возмущением $\mood$ \\
factor-узла достаточно до гиперрёбер & гипотеза & доля восстановленных n-арных выводов на подсадке (прилож.~\ref{ch:hypergraphs}) \\
$\pi$ концентрируется на хабах & известная мат.; \emph{подтв.} & распределение посещений vs степень на BA-графе (\texttt{reports/experiment\_1}) \\
$\text{rel}(B\mid q)$ на эмбеддингах — родственность, не причинная релевантность & эмпир. & distractor-vs-answer на golden-микромире; «почему»-запрос (\texttt{reports/experiment\_4}) \\
конвергенция на хабе пересиливает per-edge релевантность & эмпир. & caffeine/beverage: вентиль оценил ответ выше, но 3-источниковый хаб победил (\texttt{reports/experiment\_3--4}) \\
качество ретрива измеримо лишь на графе с ground truth & принцип & golden-микромир как фикстура, не механика (\texttt{reports/experiment\_3}) \\
longest-match даёт многословные якоря (\texttt{north\_dakota}) & эмпир. & подтверждено на микромире (\texttt{reports/experiment\_3}) \\
узкое место душит cross-cluster fusion & известная мат. & частота fusion vs ширина моста $m$ на SBM \\
consistent $\ne$ confirmed (гипотеза $\ne$ факт) & принцип & доля ложных гипотез, пойманных проверкой (пример House) \\
importance по членству не течёт на «помыть машину» & гипотеза & подсадка домена; проверить локальность importance (гл.~\ref{ch:timescales}) \\
эмбрион-подграф лучше слепого импорта & гипотеза & handcrafted/seeded подграф vs полный импорт: шум ретрива \\
«всё графовое дёшево» & гипотеза (осторожно!) & бенчмарк граф-операций на $10^5/10^6$ узлов в реальном прототипе \\
event-log позволяет сравнивать операторы & принцип & пересобрать проекцию из лога при двух алгоритмах консолидации, сравнить \\
\bottomrule
\end{tabular}
\end{center}

\begin{pitfall}[«Всё графовое дёшево» — проверять, а не верить]
Отдельно про строку выше. Тезис «бутылочное горло — LLM, графовое — миллисекунды»
верен \emph{в оптимизированной} реализации. Но в Python-прототипе с динамическим
мульти-реляционным графом, апдейтами HNSW, фильтрами по метаданным, обходом
провенанса и логированием графовая часть легко становится не «мс», а «почему оно
подвисло на странном месте». Не потому что математика плоха, а потому что инженерия
грязная. Конкретный виновник — обычно \emph{не} LLM и не математика, а \textbf{наша
обвязка}: криво написанные индексы, сериализация (JSON), локи, пересборка индексов,
фильтры по метаданным, обход провенанса и \emph{permission/scope-фильтр поверх каждого
запроса}. Тезис «граф дёшев» верен при батчинге, C-core и нормальных индексах — но не
автоматически. Вот что незаметно превращает 100k узлов в «почему оно подвисло». Поэтому
в прототипе граф-операции надо \emph{мерить}, а не постулировать их дешевизну
(гл.~\ref{ch:runtime}).
\end{pitfall}

\chapter{Заключение: единая картина}
\label{ch:conclusion}

\section{Один субстрат, много операторов}

Если из монографии стоит унести одну мысль, то вот она: память Рин — это
\emph{один типизированный темпоральный граф с метаданными}, а всё богатство
поведения — \emph{операторы} над ним. Retrieval, пропагация уверенности,
забывание и консолидация, ночные обходы, настроение и слой личности — не отдельные
подсистемы со своими моделями памяти, а разные способы читать и писать одни и те
же поля (\texttt{confidence}, \texttt{activation}, \texttt{valence},
\texttt{provenance}, \texttt{need\_to\_discover}, временной слой) на одном графе.

\section{Четыре кита связаны}

Отсюда — практический вывод о разработке. Четыре кита (правильный retrieval,
правильная пропагация уверенности, правильное забывание/консолидация, правильные
ночные обходы) сцеплены через общее состояние: забывание меняет топологию,
топология — retrieval, retrieval — что пропагируется, пропагация — что найдёт скраб.
Их нельзя тюнить порознь. Именно поэтому нужен стенд с ground truth
(глава~\ref{ch:experiment}): без него взаимодействие операторов не измерить, и
отладка «на живой Рин» найдёт проблемы слишком поздно.

\section{Сквозные принципы}

Несколько идей проходят через всю монографию и заслуживают повторения:

\begin{itemize}
\item \textbf{Гипотеза, а не факт.} Всё, что порождено выводом (fusion, реакции),
пишется провизорно, с низкой уверенностью и провенансом; управление — на границе
коммита. Это защищает от превращения обаятельной ассоциативности в фабрику ложных
убеждений.
\item \textbf{Внутрь, а не наружу.} Значимость (гейт консолидации характера)
берётся из само-оценивания, а не из внешнего одобрения; иначе строится угодник, что
противоречит ядру личности.
\item \textbf{Ингредиенты, а не реакции.} Граф хранит материал; реакция
композируется под настроение и мотивацию, а личность — смещение во время обхода, а
не таблица весов.
\item \textbf{Континуум, а не бинарность.} Врождённое/приобретённое,
семантическое/эпизодическое, стабильное/пластичное — это оси глубины консолидации,
а не выключатели.
\item \textbf{Топология незрима, но правит.} Любой локальный процесс на графе
наследует спектрально-топологический транспорт; телепорт-скраб — не украшение, а
необходимая эпистемическая гигиена с доказуемой гарантией покрытия.
\end{itemize}

\section{Единая спина двух проектов}

И последнее — то, ради чего стоило смотреть память рядом с DTI. Консолидация есть
огрубление графа; тот же вопрос «что выживает при переходе fine\,$\to$\,coarse»
стоит и там, и там. Over-squashing, спектр лапласиана, эффективное сопротивление,
персистентный лапласиан, обучаемые операторы передачи сообщений — инструменты,
работающие в обе стороны. Память Рин — «грязный», богатый испытательный стенд для
тех же операторов огрубления, где нет чистых меток; DTI — их «чистая» лаборатория с
измеримой истиной. Отрефлексировать это явно значит получить не два проекта, а один
с двумя приложениями.

\section{Что делать дальше}

Порядок разработки, вытекающий из монографии: сначала прибить \emph{схему
субстрата} (узел/ребро, метаданные, что значит коммит и кто его гейтит); затем
собрать Ярус-1 стенд и погнать батарею тестов, начиная с самого дешёвого
(концентрация $\pi$ на хабах — сразу показывает, что диагностика жива); затем
локальное правило $\Delta\text{conf}$ (знак, гейт источника, дисконт
независимости), поверх которого пороги и ветвящийся коэффициент навешиваются
тривиально; параллельно — построить $\varphi$-скорер как измерительный прибор для
отбора голоса и фильтрации выходов. Настроение, сон и слой личности — усилители
поверх субстрата, которому сначала надо стать доверенным.


\appendix
\chapter{Нотация}
\label{ch:notation}

Сводка обозначений, используемых по тексту.

\begin{longtable}{p{0.20\textwidth} p{0.74\textwidth}}
\toprule
\textbf{Символ} & \textbf{Смысл} \\
\midrule
\endhead
$G=(V,E)$ & граф памяти: узлы $V$, направленные типизированные знаковые рёбра $E$ \\
$v,\,u$ & узлы (воспоминания, факты, сущности, концепты) \\
$\Neigh(A)$ & окрестность узла $A$ (соседи, фронт распространения) \\
$\mathbf{x}_v\in\mathbb{R}^d$ & векторное представление содержания узла (эмбеддинг) \\
$c_v=\conf(v)\in[-1,1]$ & знаковая уверенность: $+$да, $-$нет, $0$ незнание \\
$a_v$ & транзиентная активация узла при обходе \\
$\text{importance}(v)$ & важность узла $\in[0,1]$: salience при кодировании + центральность; гейтит прунинг \\
$r_{uv}$ & тип отношения на ребре $u\to v$; знак $\pm$ \\
$w_{uv}$ & вес (эвиденциальная сила) ребра \\
$m_{u\to v}$ & сообщение, передаваемое по ребру: $m_{u\to v}=\phi_{r_{uv}}(h_u,h_v,e_{uv})$ \\
$\tau$ & постоянная времени слоя (интегратора) \\
$B_v$ & базовая активация ACT-R узла $v$ \\
$\pi$ & стационарное распределение случайного блуждания \\
$P$ & матрица переходов блуждающего агента; $P,Q$ — операторы проекции/поднятия при огрублении \\
$\alpha$ & доля телепортов (рестартов) в блуждании \\
$L,\,L^{+}$ & лапласиан графа и его псевдообратный \\
$\Reff(i,j)$ & эффективное сопротивление между узлами $i,j$ \\
$\sigma$ & ветвящийся коэффициент лавины пропагации \\
$\theta_{\text{emergent}},\theta_{\text{temp}}$ & верхний и нижний пороги распространения (гистерезис) \\
$\Delta\conf(B\mid A)$ & локальное изменение уверенности соседа $B$ под влиянием факта $A$ \\
$\mood_t\in\mathbb{R}^k$ & вектор настроения в момент $t$ (аффективное пространство) \\
$\varphi(\cdot)\in\mathbb{R}^{d}$ & скорер по осям характера \\
$\mathbf{t}$ & целевые координаты Рин в пространстве характера \\
$\mathrm{Beta}(\alpha,\beta)$ & апостериорное распределение успешности реакции \\
$\eta$ & темп обновления (learning rate) локального правила \\
\bottomrule
\end{longtable}

\chapter{Глоссарий внешних идей}
\label{ch:glossary}

Краткие описания инженерных и научных идей, на которые опирается монография, чтобы
читатель не был вынужден держать в голове внешний контекст. Порядок — тематический.

\section*{Память и активация}

\paragraph{Базовая активация ACT-R.} В когнитивной архитектуре ACT-R (Anderson)
сила следа в памяти есть функция \emph{частоты} и \emph{свежести} обращений со
степенным затуханием: $B_v=\ln\sum_k t_k^{-\alpha}$, где $t_k$ — давность $k$-го
обращения. Следствие: частые и \emph{разнесённые во времени} обращения укрепляют
след сильнее массовых (эффект распределённого повторения).

\paragraph{Spreading activation.} Модель ассоциативного припоминания
(Collins–Loftus): активация растекается от активированного узла по рёбрам, затухая
с расстоянием. Формальная реализация — случайные блуждания и персонализированный
PageRank.

\paragraph{Сети Хопфилда, обучение Хебба.} Ассоциативная память как граф со
взвешенными связями и динамикой активации, сходящейся к аттрактору (энергетический
минимум). Правило Хебба: связь между совместно активными узлами усиливается («что
вместе активируется, вместе укрепляется»). Наши fusion/консолидация — это
хеббовское обучение на графе.

\section*{Убеждения и истинность}

\paragraph{Система поддержки истинности (TMS/ATMS).} Механизм (Doyle;
de~Kleer), хранящий \emph{обоснования} выводов: для каждого производного убеждения
записано, из каких посылок оно получено. При ретракции посылки её обоснованные
следствия автоматически помечаются на пересмотр. Лечит «ползучую веру».

\paragraph{Недоопределённость (underdetermination).} Тезис Дюэма–Куайна: данные не
определяют теорию однозначно — одни и те же наблюдения совместимы с разными
гипотезами. Когнитивное проявление — преждевременный коллапс: ум схлопывает
суперпозицию совместимых гипотез в одну слишком рано.

\paragraph{Belief perseverance.} Устойчивость убеждения к опровержению его
основания: даже после опровержения посылки её следствия продолжают влиять.

\paragraph{Петлевое распространение (loopy belief propagation).} Распространение
уверенности по циклическому графу может привести к тому, что убеждение подтверждает
само себя по петле (эхо-камера), искусственно раздувая уверенность. Лечится
дисконтом по общей родословной свидетельств.

\section*{Динамика и критичность}

\paragraph{Самоорганизованная критичность (SOC).} Класс систем (Bak–Tang–Wiesenfeld;
образ — куча песка), самонастраивающихся к критической точке, где отклики следуют
степенному закону: почти все малы, изредка случаются крупные лавины. Нейронные
лавины (Beggs–Plenz) — проявление SOC в мозге.

\paragraph{Ветвящийся коэффициент $\sigma$.} Среднее число узлов, активируемых
одним активным. $\sigma<1$ — субкритично (затухает), $\sigma=1$ — критично,
$\sigma>1$ — суперкритично (runaway). Управляющий параметр лавины пропагации.

\paragraph{Гистерезис (триггер Шмитта).} Схема с двумя порогами (включения и
выключения), не дающая системе «дребезжать» у единственного порога.

\section*{Графы: спектр, поток, топология}

\paragraph{Стационарное распределение $\pi$.} Долгосрочная частота посещений узлов
случайным блужданием; глобальный объект (перронов собственный вектор матрицы
переходов), а не локальные вероятности.

\paragraph{Время перемешивания (mixing time).} Скорость сходимости блуждания к
$\pi$; управляется спектральным зазором матрицы переходов.

\paragraph{PageRank / персонализированный PageRank / RWR.} Случайное блуждание с
телепортом (рестартом): с вероятностью $\alpha$ прыжок в узел из распределения
рестарта. Телепорт делает блуждание эргодическим — гарантирует посещение всех
узлов.

\paragraph{Эффективное сопротивление $\Reff$.} Мера «электрической» удалённости
узлов, если рёбра — резисторы: $\Reff(i,j)=(e_i-e_j)^\top L^{+}(e_i-e_j)$. Большое
$\Reff$ между узлами ограничивает передачу влияния между ними.

\paragraph{Проводимость / константа Чигера.} Мера «узости» перешейков графа;
неравенство Чигера связывает её со спектральным зазором. Малая проводимость = узкое
место.

\paragraph{Over-squashing.} Удушение дальней связи: информация от экспоненциально
растущей окрестности сжимается в вектор фиксированной размерности, проходя через
узкое место графа; влияние дальних узлов затухает.

\paragraph{Over-smoothing.} Пересглаживание: при многократной диффузии
представления узлов сходятся к низкоразмерному консенсусу, теряя различимость.

\paragraph{Кривизна графа (Оливье, Форман).} Дискретные аналоги кривизны;
отрицательно искривлённые рёбра — мосты между сообществами, места over-squashing.
Полезны как локатор узких мест.

\paragraph{Анти-энтропийный скраб.} Приём из распределённых хранилищ (ZFS scrub;
Merkle-based anti-entropy gossip): фоновая случайная перепроверка всего хранилища,
поскольку неизвестно, какая запись повреждена.

\paragraph{Огрубление графа (graph coarsening), coarse-to-local.} Замена мелкого
графа крупным с сохранением релевантных свойств (спектра, разрезов, потоков).
Центральный вопрос — какой сигнал выживает при переходе от мелкого масштаба к
крупному; в памяти это консолидация.

\paragraph{Персистентный лапласиан, TDA.} Топологический анализ данных: как меняются
топологические признаки (числа Бетти, спектр лапласиана) при изменении масштаба
(фильтрации). Персистентный лапласиан — спектральное обобщение персистентной
гомологии.

\section*{Выразительность и обучение}

\paragraph{WL-тест, GIN.} Тест Вейсфейлера–Лемана — предел различающей способности
обычной передачи сообщений; GIN (Xu и др.) достигает этого предела инъективной
агрегацией. PNA — многостатистическая агрегация (сумма/среднее/макс/дисперсия).

\paragraph{Передача сообщений (GNN).} Параметрическая функция на графе, обучаемая
градиентным спуском: узел обновляется агрегированными сообщениями соседей. Наша
память — та же передача сообщений, но с локальными (не обучаемыми градиентом)
правилами.

\paragraph{LoRA.} Дешёвое дообучение: вместо всех весов обучается низкоранговая
добавка $\Delta W=BA$ (ранг $r\ll$ размерности), малая по памяти.

\paragraph{DPO (Direct Preference Optimization).} Дообучение под предпочтения без
отдельной модели вознаграждения: оптимизация на тройках (промпт, выбранный ответ,
отвергнутый ответ).

\section*{Эмоции и данные}

\paragraph{Теория оценивания (appraisal): OCC, Шерер.} Эмоция как результат оценки
события по измерениям: релевантность/новизна (ошибка предсказания), конгруэнтность
целям, валентность, интенсивность, атрибуция причины (agency).

\paragraph{Предиктивная обработка, ошибка предсказания.} Взгляд (Friston и др.), по
которому мозг минимизирует ошибку предсказания; величина ошибки — естественная мера
значимости события.

\paragraph{ConceptNet / Numberbatch.} Ассоциативно-таксономический граф здравого
смысла (взвешенные типизированные рёбра) и его родной эмбеддинг Numberbatch.

\paragraph{ATOMIC / ATOMIC10x.} Граф причинно-социального здравого смысла (если
субъект делает $Y$, то намерение/реакция/эффект). ATOMIC10x — крупная версия,
дистиллированная языковой моделью и отфильтрованная критиком.

\paragraph{Wikidata.} Фактологическая база с моделью утверждения: квалификаторы
(время), references (провенанс), ранг (недеструктивная ревизия). Ценна схемой, не
содержанием.

\paragraph{Generative Agents.} Симуляция агентов с потоком памяти, retrieval по
свежести$\times$важности$\times$релевантности и рефлексией (порождением обобщений);
прообраз субстрата.

\paragraph{Распределение Бета, сэмплирование Томпсона.} Сопряжённый апостериор для
доли успехов; удобно моделирует «успешность реакции» (успехи $\alpha$, провалы
$\beta$) и баланс исследования/эксплуатации.

\chapter{Сводка алгоритмов}
\label{ch:algorithms}

Псевдокод ключевых операторов, собранный в одном месте для удобства. Это скелеты,
а не финальные реализации: они фиксируют логику, обсуждённую в главах, но опускают
инженерные детали (батчинг мутаций, кэширование $\pi$, индексы эмбеддингов).

\section{Единый блуждающий агент (три политики)}

\begin{lstlisting}[style=pseudo]
def walk(G, start_nodes, mode, steps, mood, query=None):
    # mode задаёт политику через параметры: bias перехода, alpha (телепорт),
    # bias_uncertainty. Один движок — три режима (глава про обход).
    if mode == "retrieval":   alpha, bu = 0.0, 0.0
    if mode == "curiosity":   alpha, bu = 0.05, 1.0     # тянет к need_to_discover
    if mode == "sleep":       alpha, bu = 0.30, 0.0     # телепорт-скраб

    active = set(start_nodes); visited = Counter()
    for _ in range(steps):
        A = sample(active)
        if random() < alpha:                 # телепорт (рестарт)
            B = random_node(G)
        else:
            # $P(A\to B)\propto w_{AB}\cdot rel(B|query)\cdot moodcong(B|mood)\cdot unc(B)^{bu}$
            B = sample_neighbor(G, A, weights=transition_scores(
                    G, A, query, mood, bias_uncertainty=bu))
        visited[B] += 1
        maybe_fire_operations(G, B, mode)    # retrieval / fusion / inspect
        active = frontier(active, B)
    return visited
\end{lstlisting}

\section{Локальное правило и лавинная пропагация}

\begin{lstlisting}[style=pseudo]
def propagate(G, A, sigma_target=0.9):
    # Пропагация нового факта A как затухающей лавины с контролем ветвления.
    frontier = link_initial(G, A)            # A -> N(A), напр. последние top-k
    while frontier:
        next_frontier = []
        # нормировка под целевой ветвящийся коэффициент sigma < 1
        gain = branching_normalizer(frontier, sigma_target)
        for B in frontier:
            d = delta_conf(G, A, B) * gain    # см. ниже: знак, источник, независимость
            G.confidence[B] = clip(G.confidence[B] + d, -1, 1)  # знаковая уверенность
            th_e = theta_emergent(G, B)       # растёт с глубиной консолидации
            if abs(d) > th_e:                 # крупная переоценка -> в ширину
                next_frontier += neighbors(G, B)
            elif abs(d) > THETA_TEMP:         # промежуточная зона -> пометить
                mark(G, B, "need_to_revise")
            # else: глохнет здесь
        frontier = next_frontier

def delta_conf(G, A, B):
    sign      = edge_sign(G, A, B)            # supports (+) / contradicts (-)
    strength  = edge_weight(G, A, B)
    gate      = G.confidence[A] * source_trust(G, A)   # карантин нового факта
    indep     = independence_discount(G, A, B)         # анти-эхо по провенансу
    return ETA * sign * strength * gate * indep * G.confidence[A]
\end{lstlisting}

\section{Забывание, прунинг, консолидация}

\begin{lstlisting}[style=pseudo]
def nightly_maintenance(G):
    for v in G.nodes:
        G.strength[v] = act_r_base_level(G.access_log[v])   # частота + свежесть
    for v in low_activation(G):
        if importance(G, v) low and confidence(G, v) low:
            drop(G, v)                                       # редкое-и-неважное
    for cluster in redundant_episodic_clusters(G):
        merge_to_semantic(G, cluster)   # огрубление: эпизоды -> обобщение
    # эмбриональным узлам — меньшая скорость распада (континуум глубины)
\end{lstlisting}

\section{Телепорт-скраб (анти-энтропия)}

\begin{lstlisting}[style=pseudo]
def scrub(G, budget):
    # Случайная инспекция: гарантия покрытия за счёт эргодичности телепорта.
    for _ in range(budget):
        v = random_node(G)                    # рестарт куда угодно
        for w in local_walk(G, v, depth=k):
            if contradicts_any(G, w):         # столкновение -> ревизия
                demote(G, w); mark(G, justificands(G, w), "need_to_revise")
            if flagged(G, w, "need_to_discover"):
                maybe_fetch_external(G, w)     # любопытство/внешний вход
\end{lstlisting}

\section{Настроение и монитор руминации}

\begin{lstlisting}[style=pseudo]
def mood_update(mood_layers, event):
    # Стек протекающих интеграторов с разными tau; актуальное = суперпозиция.
    for L in mood_layers:                     # session, day, ..., life
        mood_L = (1 - 1/L.tau) * mood_L + (1/L.tau) * appraise(event)
    mood_now = sum(w_L * mood_L for L in mood_layers)
    if autocorr(history(mood_now)) > RUMINATION_TH:   # залипла валентность
        trigger(one_of=["reframe", "attention_shift", "invoke_sleep"])
    return mood_now
\end{lstlisting}

\section{Оценивание, значимость и консолидация реакции}

\begin{lstlisting}[style=pseudo]
def encode_event(G, event, mood_prev, mood_now):
    ap = appraisal(event)                     # surprise, goal_congruence, valence...
    salience = ap.surprise * ap.goal_relevance
    #  дешёвый прокси значимости, почти бесплатный:
    salience_cheap = abs(mood_now - mood_prev)
    consolidate(G, event, weight=salience)    # гейт значимости, НЕ одобрения
    if ap.surprise high: mark(G, event, "need_to_discover")

def update_reaction(G, situation, user, reaction, received_ok):
    e = reaction_edge(G, situation, user, reaction)   # условь по (ситуация, отношение)
    if received_ok: e.beta.alpha += 1
    else:           e.beta.beta  += 1                 # угасание, не второй коммит
    e.confidence = e.beta.alpha / (e.beta.alpha + e.beta.beta)
    if e.confidence > CONSOLIDATE_TH: promote_to_slow_layer(G, e)
\end{lstlisting}

\section{Скорер характера и отбор голоса}

\begin{lstlisting}[style=pseudo]
def phi(text, axes):                          # LLM-судья по рубрикам осей
    return [judge(text, rubric_a) for rubric_a in axes]   # вектор в R^d

def select_voice_corpus(candidates, target_t, axes):
    scored = [(c, phi(c, axes)) for c in candidates]
    # отбор/перевзвешивание, чтобы среднее выборки приблизилось к target_t
    return greedy_match_mean(scored, target_t)   # не усреднять муз, а фильтровать
\end{lstlisting}

\section{Стенд эксперимента}

\begin{lstlisting}[style=pseudo]
def run_experiment(build_graph, plant, operator, metric):
    G = build_graph()          # ярус 1: SBM/BA/WS + подсаженные сценарии
    truth = plant(G)           # ground truth: что должно произойти
    operator(G)                # walker / propagate / scrub / decay ...
    return metric(G, truth)    # precision/recall починки, сходимость, покрытие
\end{lstlisting}

\chapter{Будущее: стоит ли переходить на гиперграфы}
\label{ch:hypergraphs}

Это приложение — рассуждение о структурном направлении, а не готовая рекомендация к
внедрению. Вопрос: логично ли в будущем сделать память Рин \emph{гиперграфом}, а не
обычным графом.

\section{Проблема, которая это мотивирует}

Обычный граф связывает узлы \emph{попарно}. Но вывод часто \emph{n-арен}: он рождается
только при одновременном присутствии нескольких посылок, и ни одно подмножество не
достаточно. Канонический пример (глава~\ref{ch:fusion}):
\[
A=\text{«есть коты»},\quad B=\text{«живёт в Португалии»},\quad C=\text{«ездит на пары»},
\]
\[
X=f(A,B,C)=\text{«купи котам корм на выходных (в будни по пути с пар не успеть)»}.
\]
Каждая посылка ещё и подтягивает свой контекст (коты$\to$нужен корм; Португалия$\to$
магазины рано закрываются; универ$\to$поздний возврат; см.\ главу~\ref{ch:fusion}), и
только их совмещение даёт $X$. Ни $A\wedge B$, ни $B\wedge C$, ни $A\wedge C$ не дают
$X$. Попарные рёбра между посылками могут вообще отсутствовать: «Португалия» и «коты»
напрямую не связаны.

\section{Почему симплексы — плохой ответ}

Первое, что приходит в голову, — симплициальные комплексы: тройка $\{A,B,C\}$ как
2-симплекс (треугольник). Но у симплициального комплекса есть жёсткое требование
\emph{downward closure}: если есть 2-симплекс $\{A,B,C\}$, то обязаны быть и все его
грани — рёбра $A\!-\!B$, $A\!-\!C$, $B\!-\!C$. А в нашем случае этих попарных связей
\emph{нет и не должно быть}: «Португалия» и «коты» не связаны. Симплициальная память
вынудила бы нас проводить \emph{ложные} попарные рёбра ради формальной корректности
структуры. Это плохой обмен.

\section{Гиперребро — прямой ответ}

\emph{Гиперребро} связывает \emph{множество} узлов сразу: $e=\{A,B,C\}$, и оно
\emph{не требует} существования попарных рёбер между ними. Это ровно та семантика,
что нужна n-арному выводу: «эти три вместе порождают это», без обязательства, что
они связаны попарно.

\begin{intuition}[Гиперграф против графа]
В графе «дружба» — всегда пара. В гиперграфе можно сказать «эти пятеро были в одной
поездке» одним ребром, не утверждая, что каждый дружит с каждым. Для памяти это
разница между «связь между двумя воспоминаниями» и «условие из нескольких
воспоминаний, порождающее вывод».
\end{intuition}

\section{Ключевое наблюдение: мы уже наполовину гиперграф}

Важно не пропустить: \textbf{текущая схема Рин уже кодирует гиперрёбра} — просто не
называя их так. Когда fusion порождает узел $X$ с провенансом $\{A,B,C\}$
(глава~\ref{ch:fusion}), это \emph{factor-graph}-кодировка гиперребра
$\{A,B,C\}\to X$: сам вывод $X$ становится узлом, а его рёбра-обоснования указывают
на посылки. Попарных рёбер между $A,B,C$ при этом не требуется. Формально это просто
\textbf{двудольный граф} (узлы-посылки $\leftrightarrow$ узлы-факторы), и его
\emph{полностью достаточно} для n-арного вывода — никакой высшей размерности не нужно.

Значит «переход на гиперграфы» — это не «переписать всё», а более узкий вопрос:
\emph{делать ли гиперрёбра первоклассными объектами} (со своим обходом и спектральной
теорией), или достаточно узла-фактора.

\section{Первоклассные гиперрёбра: за и против}

\paragraph{За.} Появляются \emph{гиперграф-нативные} операции: обход/spreading
activation, уважающий \emph{совместность} срабатывания (гиперребро «зажигается»,
когда активны \emph{все} его узлы, а не любой один); гиперграфовые лапласианы и
спектральная диагностика (аналоги главы~\ref{ch:topology} для n-арной структуры). Это
может точнее моделировать «условие из нескольких посылок».

\paragraph{Против.} Алгоритмика гиперграфов менее зрелая и дороже графовой. На
практике гиперграф часто \emph{сводят} к графу — clique-expansion (заменить
гиперребро кликой, что как раз возвращает ложные попарные рёбра!) или
star-expansion (ввести узел-фактор — то есть ровно наша factor-кодировка). Есть
спектральная теория гиперграфов (нормализованный гиперграфовый лапласиан; тензорные
методы), но она сложнее и вычислительно тяжелее, а динамический мутирующий гиперграф
— тем более.

\section{Рекомендация}

\begin{principle}[Гиперграф — да, но постепенно]
Ближний горизонт: оставить \emph{factor-узлы} (вывод $X$ с провенансом-множеством).
Это уже даёт n-арную семантику без гиперграфовой машинерии и без ложных попарных
рёбер. Средний горизонт: сделать гиперрёбра первоклассными \emph{только там}, где
нужны гиперграф-нативные операции — совместное срабатывание при обходе и n-арная
спектральная диагностика. Симплициальные комплексы обходить стороной: их downward
closure противоречит разреженности реальных выводов.
\end{principle}

Заметим, что star-expansion (факторные узлы) и первоклассные гиперрёбра — это два
взгляда на одно и то же; переход между ними инкрементален, а не революционен.
Полное же объединение узлов, рёбер и выводов в объекты одного рода разной
размерности — это уже клеточные комплексы (приложение~\ref{ch:cellcomplex}).

\chapter{Будущее: память на клеточных комплексах}
\label{ch:cellcomplex}

\begin{pitfall}[Вердикт вперёд: это НЕ ближний шаг]
Чтобы никто не принял эту главу за роадмап: \textbf{на текущую задачу клеточные
комплексы — колоссальный over-engineering, и мы их сейчас не строим}. Спектр
Ходж-лапласиана $L_k=\partial_{k+1}\partial_{k+1}^\top+\partial_k^\top\partial_k$ на
\emph{динамически мутирующем} графе сожрёт вычислительные ресурсы. Factor-узлы
(приложение~\ref{ch:hypergraphs}) — по сути двудольный граф — \emph{полностью
достаточны} для n-арного вывода. К симплексам и высшим размерностям возвращаемся
только упёршись в реальный предел выразительности. Ниже — почему это красивая
\emph{северная звезда}, а не спринт.
\end{pitfall}

Это приложение — про наиболее амбициозное структурное направление. Оно серьёзнее
гиперграфов (приложение~\ref{ch:hypergraphs}) и относится к области R\&D, но обсудить
его стоит, потому что оно даёт \emph{единую} картину.

\section{Идея: всё — объекты одного рода, разных размерностей}

В обычном графе узлы и рёбра — сущности \emph{разной} природы. Идея клеточного
комплекса: сделать их сущностями \emph{одного} рода, но разной \emph{размерности}.
\emph{Клетки}: 0-клетки — узлы (воспоминания, концепты), 1-клетки — рёбра (связи),
2-клетки — выводы из нескольких посылок, и так далее. Все они — «клетки», просто
живущие на разных этажах.

Связывает этажи \emph{граничный оператор} $\partial$: у каждой $k$-клетки есть
\emph{граница} — набор $(k{-}1)$-клеток, на которых она «натянута». Граница 1-клетки
(ребра) — два её узла; граница 2-клетки (вывода) — цикл/цепочка из посылок.

\begin{intuition}[Вывод рождается, когда построена его граница]
Твоя формулировка точна: вывод \emph{рождается}, когда его граница построена на
посылках. В терминах комплекса это \emph{приклеивание} $k$-клетки по её границе:
как только в памяти есть подходящая цепочка посылок (граница), поверх неё можно
«натянуть» клетку-вывод. Узел-вывод — это заполненная область, чей контур —
посылки.
\end{intuition}

\section{Почему это решает возражение к симплексам}

Здесь ключевой технический факт. Симплициальные комплексы требуют downward closure
(все грани обязаны существовать) — и потому плохи для разреженных n-арных выводов
(приложение~\ref{ch:hypergraphs}). \emph{Клеточные} (CW) комплексы этого
\textbf{не} требуют: $k$-клетку можно приклеить к границе \emph{произвольным}
приклеивающим отображением, не заполняя все промежуточные симплексы. То есть
клеточные комплексы \emph{снимают} ровно то возражение, из-за которого симплексы не
подходят, — и при этом дают единую алгебраическую структуру. Это делает их более
принципиальным ответом на n-арность, чем и симплексы, и даже голые гиперрёбра.

\section{Ходж-лапласианы: пропагация между этажами}

Граничные операторы $\partial_k$ порождают \emph{Ходж-лапласианы} на каждой
размерности:
\[
L_k=\partial_{k+1}\partial_{k+1}^\top+\partial_k^\top\partial_k .
\]
Это обобщение графового лапласиана (случай $k=0$) на клетки высших размерностей. Оно
даёт \emph{принципиальную} диффузию не только между узлами, но и между рёбрами,
между выводами, а также \emph{между этажами} (вверх и вниз по $\partial$). Иначе
говоря, вся динамика уверенности и активации (главы~\ref{ch:propagation},
\ref{ch:topology}) получает единое спектральное описание, охватывающее и выводы как
полноправные объекты.

\section{Прямая связь с чтением по DTI}

Это не абстрактная экзотика: ровно эта математика лежит в статьях, которые уже есть в
DTI-волте проекта. Bodnar и др. «WL Go Cellular» и «WL Go Topological» — передача
сообщений на клеточных и симплициальных комплексах; sheaf-лапласианы (cellular sheaf)
— их обогащение. Таким образом «память на клеточных комплексах» — очередной мост
между двумя проектами автора: инструменты топологического и клеточного обучения из
магистерской работы напрямую применимы к структуре памяти Рин.

\section{Честная оценка: северная звезда, не ближний шаг}

\begin{principle}[Клеточные комплексы — направление, а не спринт]
Это самая красивая и самая объединяющая картина: один род объектов, разные
размерности; вывод — высшая клетка; пропагация — Ходж-лапласиан по этажам. Но
выигрыш над прагматичной factor-кодировкой (приложение~\ref{ch:hypergraphs}) пока
\emph{не доказан}, а цена высока: динамические клеточные комплексы с приклеиванием и
отклеиванием клеток на мутирующей памяти, эффективные Ходж-лапласианы в реальном
времени — всё это серьёзный R\&D. Разумный порядок: ближний горизонт — factor-узлы;
средний — первоклассные гиперрёбра при необходимости; дальний/R\&D — клеточные
комплексы ради единого алгебраического описания.
\end{principle}

\begin{workedbox}[Как это выглядело бы у Рин]
Узлы «Португалия», «коты», «универ» — 0-клетки. Вывод «купи корм по пути» — 2-клетка,
приклеенная к границе из этих трёх посылок (без требования, чтобы посылки были
связаны попарно). Позже, если посылка «есть коты» опровергается, граница 2-клетки
рушится — и вывод \emph{автоматически} теряет опору (ретракция как разрушение
границы). Ходж-лапласиан при этом единообразно описывает, как понижение уверенности
течёт с этажа посылок на этаж вывода. Красиво — и именно поэтому стоит держать это
как цель, к которой архитектура может дозреть.
\end{workedbox}

\chapter{Контракт Stage 1: минимальный срез, который строим первым}
\label{ch:stage1}

Монография зрелая, и улучшать её можно бесконечно. Но полезнее не ещё одна красивая
глава, а \emph{зафиксированный маленький вертикальный срез}, который можно собрать за
неделю-две и начать \emph{ломать} на стенде. Это приложение — контракт такой сборки:
не весь организм памяти, а сквозная линия от захвата до речи.

\section{Сквозной срез (pipeline v0)}

\begin{lstlisting}[style=pseudo]
event_log            # неизменяемый лог (гл. 17): append observed/committed/deleted
  -> typed nodes     # схема узла v0 (ниже), node_kind обязателен
  -> write/skip gate # epistemic + salience + anti-fusion; skip логируется (гл. 17)
  -> forensic index  # точный поиск по метаданным/провенансу/времени (гл. 17)
  -> anchor retrieval# вектор-якоря + ограниченный обход (гл. 6)
  -> context-packer v0 # с метками observed/hypothesis/low_conf/... (гл. 17)
  -> inspector CLI   # показать узлы/рёбра/поля/лог, вручную инспектировать
  -> synthetic tests # маленький граф с ground truth (гл. 18)
\end{lstlisting}

\section{Схема узла v0 (минимум полей)}

\texttt{id}, \texttt{content}, \texttt{embedding}, \texttt{node\_kind},
\texttt{confidence}$\in[-1,1]$, \texttt{provenance} (или support-skeleton),
\texttt{time}, \texttt{layer}. Остальные поля (\texttt{valence},
\texttt{importance}, \texttt{domains}, \texttt{need\_to\_*}) добавляются по мере
надобности, но \texttt{node\_kind} обязателен с самого начала — он задаёт смысл
\texttt{confidence} (гл.~\ref{ch:substrate}).

\section{Первые 6--8 сущностей (что реально кодим)}

\begin{enumerate}\itemsep1pt
\item \textbf{Event log} (append-only) $+$ пересборка графа-проекции из лога.
\item \textbf{Typed node store} (KV) $+$ edge store (adjacency).
\item \textbf{Write/skip gate} с логированием причины skip.
\item \textbf{Forensic index} (точный поиск/удаление по метаданным) — включая
работающий delete-cascade (гл.~\ref{ch:runtime}, раздел~\ref{sec:forget}).
\item \textbf{Anchor retrieval} (вектор-индекс $+$ ограниченный обход).
\item \textbf{Context-packer v0} (с эпистемическими метками).
\item \textbf{Inspector CLI} — «memory inspector»: смотреть узлы, рёбра, поля, лог.
\item \textbf{Synthetic graph tests} (SBM/BA $+$ подсаженные сценарии).
\end{enumerate}

\section{Что сознательно игнорируем на Stage 1}

Чтобы срез остался тонким: mood-динамика и её слои, личность/Rinness и drive-вектор,
сон/скраб и разрешение противоречий, тяжёлый fusion (только anti-fusion-скелет и
провизорные кандидаты), гиперрёбра и клеточные комплексы, governance-слой,
мета-обучение операторов. Это не отменяется — это \emph{откладывается}; красивые
главы ждут, пока нижний срез не станет доверенным.

Три конкретных инженерных упрощения на первую итерацию (полные механизмы описаны, но
дороги в реализации):
\begin{itemize}\itemsep2pt
\item \textbf{Удаление — только \texttt{tombstone}.} Crypto-shredding и
taint-propagation (гл.~\ref{ch:runtime}) безупречны, но каскад «яда» по кэшам, индексам
и бэкапам съест месяцы. В прототипе — tombstone $+$ пересборка проекции; криптостирание
позже.
\item \textbf{Коллизии консолидации — компоненты связности, не модулярность.} Полный
со-консолидационный граф $+$ community detection тяжёл. На старте: простые
\emph{connected components} по графу предложений с отсечением по жёсткому порогу
пересечения (LSH\,/\,MinHash). Модулярность — когда упрёмся.
\item \textbf{Семантический вентиль — только косинус кэш-эмбеддингов.}
\texttt{sense/type}-гейт против омонимии (гл.~\ref{ch:walker}) требует sense-слоя,
которого ещё нет. \emph{Implementation risk}: Stage 1 \emph{сознательно не обещает}
сильную защиту от омонимии — либо терпим протечки, либо руками делаем крошечную
sense-схему на 20--30 демо-кейсов. Лечим, когда эмпирически упрёмся.
\end{itemize}

\section{Критерии готовности v0 (как поймём, что работает)}

\begin{itemize}\itemsep1pt
\item запись и \emph{skip} работают, причина skip видна в логе;
\item \texttt{tombstone}-удаление сносит узел из проекции, это проверяемо inspector-ом
(taint-каскад — отдельная поздняя цель);
\item anchor-retrieval на синтетическом графе возвращает осмысленную окрестность;
\item реплей event-log \emph{перестраивает} тот же граф-проекцию (event sourcing);
\item хотя бы один тест из батареи отказов (гл.~\ref{ch:experiment}) — например,
$\pi$-концентрация на хабах или точное удаление — зелёный.
\end{itemize}

\begin{principle}[Первый успех — не «Рин помнит», а «физика графа наблюдаема»]
Формулировка критерия важнее списка. Если первым мерилом сделать «Рин умно пользуется
памятью в диалоге», сразу утонешь в LLM-шуме и не поймёшь, работает граф или врёт
модель. Поэтому Stage 1 — это \emph{memory inspector v0}, а не «память v0»: успех —
когда мы можем руками создать event, увидеть узел, увидеть retrieval, увидеть proposal,
увидеть skip\,/\,commit, увидеть support-skeleton. Наблюдаемость внутренней механики
раньше, чем «умное поведение».
\end{principle}

\begin{principle}[Дальше — строить и ломать, а не расширять]
На этом теоретический слой монографии \emph{заморожен}. Следующий лучший ход —
собрать этот срез, прогнать synthetic tests, и позволить реальности править
гипотезы. Каждая «красивая глава» включается только тогда, когда нижний срез её
выдержит.
\end{principle}

\chapter{Будущее: память как динамика (резервуар, ассоциативные аттракторы, кванты)}
\label{ch:dynamical}

\begin{pitfall}[Вердикт вперёд: это смена субстрата, а не смена топологии]
Приложения~\ref{ch:hypergraphs} и~\ref{ch:cellcomplex} спрашивали «какой \emph{формы} должен
быть граф». Это приложение спрашивает радикальнее: \emph{а нужен ли символический граф
вообще}, или память стоит держать как \emph{динамику} — состояние большой рекуррентной
системы, где факт не хранится объектом, а «протоптан» в связях. Идея красивая, глубоко
укоренённая в нейронауке, и \textbf{мы её сейчас не строим}. Но не потому, что она наивна —
наоборот, разберём её всерьёз. Причина в том, что для \emph{компаньона} с требованиями
редактируемости, забывания-по-команде, провенанса и поддержки истинности распределённая
память в весах ломает ровно те свойства, ради которых мы и завели типизированный граф. Ниже
— честный разбор: что здесь ценно (и что мы, оказывается, уже делаем), что дорого, и как
оставить дверь открытой, не переписывая мозг.
\end{pitfall}

\section{Постановка: память в весах, воспоминание как траектория}

Наш субстрат (гл.~\ref{ch:substrate}) — типизированные узлы и явные рёбра: у «узла Helena»
есть содержание, у ребра — тип и вес. Альтернатива, которую предлагает эта ветка, меняет
само определение:
\[
\text{память} = (G, W, \theta),\qquad
\text{воспоминание} = \text{траектория в пространстве состояний},\qquad
\text{обучение} = \Delta W .
\]
Здесь $G$ — топология (какие пути вообще есть), $W$ — веса связей (насколько они
протоптаны), а $\theta$ — параметры узлов-вычислителей. Узел \texttt{\#48372} не значит
«Лиссабон»; смысл распределён по тысячам узлов и существует только как \emph{паттерн
активации} $+$ \emph{связность} $+$ \emph{траектория}. Запрос не «ищет chunk», а
\emph{возмущает} систему: бросаем заряд, сеть эволюционирует, из неё выныривает активный
подграф, и уже его вербализует LLM.

\begin{intuition}[Бассейн, камни, датчики]
Классический образ резервуарных вычислений: бассейн воды, в который бросают камни (входной
поток). Каждый камень поднимает волны; следующий приходит, пока прежние не улеглись, — и
поверхность в момент $t$ несёт отпечаток нескольких последних бросков. Датчики по краям
считывают состояние, и простой линейный слой учится по ним отвечать. Сам бассейн обучать не
надо: его физика уже превращает историю входов в богатое состояние. Reservoir computing —
это \emph{случайное нелинейное динамическое отображение признаков для последовательностей},
аналог kernel trick, но во времени.
\end{intuition}

Важная оговорка сразу: \emph{классический} резервуар (Echo State Network, Liquid State
Machine) держит $W_{\text{res}}$ \emph{замороженным} и обучает только считывающий слой. Он
обладает \emph{затухающей} памятью (fading memory) на конечный горизонт — прекрасно для
временных рядов, но не для «узнала в 2026-м, вспомнила в 2031-м». Для памяти длиной в годы
нужен \emph{пластичный} резервуар ($W$ медленно меняется опытом) плюс механизм устойчивого
хранения из мира \emph{ассоциативных аттракторов} (Hopfield/CA3): воспоминание — это
\emph{бассейн притяжения}, частичный cue скатывается в него (pattern completion). То есть это
не «reservoir computing» в чистом виде, а гибрид на пересечении резервуара, ассоциативной
памяти, аттракторных сетей и пластичной динамики.

\section{Математическая ось: один лапласиан — три динамики}

Самое элегантное в этой ветке — то, что \emph{одна и та же} топология допускает три
принципиально разные динамики. Для графа с лапласианом $L=D-A$:
\begin{align}
\text{диффузия:}\quad & \dot x = -Lx & \Rightarrow\quad & x(t)=e^{-tL}x(0),\\
\text{волна:}\quad & \ddot x = -Lx & \Rightarrow\quad & x(t)=\cos(t\sqrt{L})\,x(0),\\
\text{квантовое блуждание:}\quad & i\dot\psi = L\psi & \Rightarrow\quad & \psi(t)=e^{-itL}\psi(0).
\end{align}
Один $L$ порождает $e^{-tL}$ (рассеивающую память), $\cos(t\sqrt L)$ (волновую, с
интерференцией и без затухания) и $e^{-itL}$ (унитарную квантовую). Это не словесная
аналогия, а линейная алгебра: «бросить заряд в узел и смотреть, куда потечёт» имеет буквальный
волновой и квантовый аналоги на той же сети.

\begin{intuition}[Мы уже живём в углу $e^{-tL}$]
Полезно осознать: наш ретрив (гл.~\ref{ch:walker}) — это \emph{и есть} диффузионный угол
этой оси. `diffuse()` считает усечённое $e^{-tL}$-подобное распространение из seed-якорей;
«вспышки, слабеющие к окраинам» — это спад $e^{-tL}$. То есть значительная часть красоты
этого приложения у нас уже реализована, просто в скромной, редактируемой форме. Волновой и
квантовый углы — это \emph{другие операторы над той же топологией}, и их честно можно
сравнить на одном графе (см.\ открытые вопросы в конце).
\end{intuition}

\section{Нейронаучная опора (принцип, а не крысиный мозг)}

Ветка опирается на реальные результаты, и стоит назвать их аккуратно. Обучение оставляет
устойчивые следы в самой сети (engram-клетки: клетки, участвовавшие в формировании
воспоминания, можно искусственно возбудить и вызвать его; хранение связывают с синаптическим
усилением engram-ансамблей). Гиппокамп работает не как «жёсткий диск», а как система
связывания/индексации и реактивации распределённых корковых представлений (hippocampal
indexing theory), причём со временем память дрейфует от гиппокампа к коре (systems
consolidation, replay). Complementary Learning Systems (CLS) прямо противопоставляет
\emph{быстро} обучающийся гиппокамп \emph{медленно} обучающейся коре — и это вычислительно
мотивированная защита от катастрофического забывания. Аттракторная память (Hopfield; CA3 как
autoassociative network) хранит паттерн как бассейн притяжения. Мы заимствуем \emph{принцип}
(быстрый индекс $+$ медленное распределённое хранилище $+$ восстановление по частичному cue),
не пытаясь эмулировать биологию до деталей.

\section{Предлагаемая гибридная архитектура}

Синтез, к которому сходится обсуждение, — не «резервуар вместо графа», а шестислойный объект
плюс архив сбоку:
\begin{enumerate}\itemsep2pt
\item \textbf{Event encoder.} Приходит не «факт», а \emph{событие} целиком — контент, говорящий,
тон, контекст, время, настроение — и кодируется в вектор $x_t$. Источник знания становится
\emph{частью паттерна}, а не строкой БД (значит, память сможет и \emph{ошибиться} в источнике
— по-человечески).
\item \textbf{Hippocampal pattern separator / index.} Из $x_t$ строится разреженный код
$h_t=\mathrm{kWTA}(Px_t)$ (активны немногие units из многих); похожие события разводятся в
непохожие коды (pattern separation). Это «умный стартер»: куда бросить заряд.
\item \textbf{Fast associative memory.} Быстрая Hebbian-запись «эти вещи произошли вместе»
($H\!\leftarrow\!H+\eta_h h_t h_t^\top$), связывающая индекс с распределённым состоянием.
\item \textbf{Slow plastic cortical reservoir.} Большая разреженная рекуррентная сеть
\[
r(t{+}1)=(1-\alpha)\,r(t)+\alpha\,\sigma\!\big(Wr(t)+Ux(t)+Bh(t)\big),
\]
где $W$ \emph{медленно} пластичен ($\eta_c\ll\eta_h$). Быстрый гиппокамп $+$ медленная кора —
это и есть CLS.
\item \textbf{Continuous activation field.} Живое поле $r(t)$ — «о чём Рин сейчас думает».
\item \textbf{Latent-to-LLM interface.} Мост от $r(t)$ к языковой модели.
\end{enumerate}
Плюс \textbf{replay/consolidation} в простое (ночью гиппокамп многократно проигрывает след,
медленно перестраивая $W$ — «переваривание дня» не через \texttt{summarize()}, а через
проигрывание), и \textbf{граф-архив} как «чёрный ящик самолёта» сбоку.

\begin{principle}[Граф не выбрасывается — он становится архивом и учителем]
Ключевой примирительный ход: наш явный граф перестаёт быть \emph{моделью} человеческой памяти
и становится (а) \emph{эпизодическим архивом} (что реально произошло, когда, откуда знаем) и
(б) \emph{учителем} для пластичной памяти (материал для replay-консолидации). Тогда
специализация чистая: \emph{граф} — «что действительно было записано», \emph{резервуар} — «что
Рин из этого сейчас помнит» (неточно, ассоциативно, с интерференцией). Это не дублирование, а
разделение ролей.
\end{principle}

\section{Настроение как нейромодулятор и «пирожки Алисы»}

Отдельная и, на мой взгляд, самая \emph{сразу применимая} идея ветки: mood перестаёт быть
вектором, который кладут в промпт, и становится \emph{нейромодулятором} — он меняет
\emph{геометрию} recall:
\[
W_{\text{effective}} = \gamma(m)\,W,\qquad
A_{\text{effective}} = A\odot M(\mathrm{mood}).
\]
Тот же cue при разном настроении уходит по разным траекториям («помнишь Тбилиси?» $\to$ переезд
$\to$ Ада; или $\to$ ОИЯИ $\to$ смена жизни). Инерция настроения тут — не недостаток, а
\emph{нужное} свойство: медленный low-pass-контроллер задаёт устойчивый ландшафт для быстрых
процессов и не «дёргает руль» на каждый токен.

\begin{intuition}[Гомеостаз как гормоны: заманчиво, но чинит болезнь, которую само же создаёт]
Предложение «мониторить энергию поля $E(t)=\frac1N\sum r_i^2$ и впрыскивать седатив (поднять
утечку $\lambda$, снизить крутизну $\beta$) при перевозбуждении или кофеин при апатии» — это
настоящая идея из neuromodulated networks, и она красива. Но будем честны: это \emph{лекарство
от нестабильности, которую вносит сам пластичный резервуар}. На нашем текущем субстрате
диффузия ограничена по построению (decay $+$ floor $+$ max\_steps $+$ max-CRDT активация,
гл.~\ref{ch:walker}) — «эпилептического шторма» просто негде взяться, и гормональный
контроллер не нужен. Более того, он добавляет ровно то, чего мы весь проект избегали:
свободные параметры ($E_{\min},E_{\max}$, коэффициенты впрыска), которые нечем калибровать без
downstream-сигнала. Вывод сбалансированный: \emph{mood-как-нейромодулятор} — берём хоть сейчас
(он и так входит в правило перехода через $\mathrm{mc}(B\mid\mood)$); \emph{гормональный
гомеостаз} — осмыслен только если уже ушли в распределённый пластичный резервуар, то есть это
инструмент \emph{той} ветки, а не нашей.
\end{intuition}

\section{Честный разбор: что покупаем и чем платим}

\paragraph{Что действительно ценно.} Ассоциативная \emph{реконструкция} вместо поиска chunk'а;
человекоподобная память (нечёткое время, ассоциативные ошибки источника, интерференция как
feature); mood как геометрия recall; CLS как защита от катастрофического забывания; и красивая
ось $e^{-tL}/\cos(t\sqrt L)/e^{-itL}$, дающая единый язык для diffusion/wave/quantum памяти.

\paragraph{Чем платим — и это надо назвать прямо.} Три технических риска и один
принципиальный.
\begin{itemize}\itemsep3pt
\item \textbf{Readout-бутылочное горлышко.} Перевод $r(t)\in\mathbb{R}^{10^5}$ обратно в
токены — самая хрупкая часть. Тупой декодер (MLP) может не извлечь смысл; умный — \emph{сам}
сделает всю работу и замаскирует бесполезность динамики (нужен экспериментальный контроль:
сначала нарочно \emph{тупой} $D$). А чтобы декодер обучить, нужны пары «состояние $\to$
правильный ответ» — которых для автобиографии \emph{нет}. Откуда берётся обучающий сигнал —
главный незакрытый вопрос. Элегантный вариант (резервуар модулирует hidden states LLM через
cross-attention/prefix) концептуально безупречен, но требует глубокого fine-tuning локальной
модели и почти исключает API-бэкенд; дешёвый вариант (декодировать в текст) — это, по сути,
RAG с лишними шагами, теряющий главную приманку «память остаётся неязыковой».
\item \textbf{Пластичность против гомеостаза (край хаоса).} Держать $\rho(W)\lesssim1$ на
\emph{непрерывно меняющемся} $W$ — нерешённая в вычислительной нейронауке задача: пластичность
воюет с устойчивостью. Плюс аттракторная память имеет малую ёмкость и \emph{ложные}
аттракторы; «плотные» версии Хопфилда чинят ёмкость, но математически превращаются в
attention — и тогда закономерен вопрос «а почему не взять память самого трансформера».
\item \textbf{Цена inference.} Прогонять $N$ шагов динамики по графу в сотни тысяч узлов на
\emph{каждый} запрос, да ещё фоном для replay — тяжело; нужна серьёзная оптимизация разреженных
умножений. Это не смертельно, но и не бесплатно (наш же принцип «граф дёшев только при чистой
обвязке»).
\item \textbf{Принципиальное: теряются свойства, ради которых мы строили граф.}
Редактируемость, забывание-по-команде (крипто-шред), провенанс, поточечная ревизия (TMS),
метки contested, семантика node\_kind — всё это распределённый резервуар \emph{уничтожает}.
Нельзя крипто-шреднуть факт из $W$; нельзя показать провенанс; нельзя гарантированно забыть
конкретное («Рин, забудь последние 10 минут» — у компаньона это \emph{UX/приватность}, а не
когнитивная поэзия). Возражение «человек тоже не забывает по команде» верно \emph{когнитивно},
но игнорирует, что у Рин есть \emph{юридическо-продуктовые} требования, которых у человека нет.
Архив-«чёрный ящик» спасает сырое удаление, но \emph{не} отменяет того, что резервуар уже
\emph{выучил} факт в $W$ — распределённо и не-разучиваемо.
\end{itemize}

\begin{principle}[Ключевая развилка: «не хранить строки» не значит «не иметь адресуемых узлов»]
Здесь легко проскочить подмену. Есть два разных смысла «памяти без фактов»:
\emph{(а)} не хранить факт как \emph{человекочитаемую строку} — узлы могут быть эмбеддингами;
это \emph{дёшево и совместимо} с нашим субстратом; \emph{(б)} вообще не иметь \emph{адресуемых
объектов памяти} — всё размазано в $W$; это и покупает человекоподобность, и \emph{ломает}
редактируемость/провенанс/забывание. Ценное — это (а) $+$ ассоциативная динамика; дорогая
переоценка — (б). А наш ретрив уже делает (а)$+$динамику: «возмути узел и смотри, что
загорится» — это диффузия по типизированному графу. Значит львиную долю пользы этой ветки мы
\emph{уже} имеем, не платя за (б). По-настоящему \emph{новых} адоптируемых кусков три:
пластичные рёбра (Hebbian-усиление от со-активации — по сути формализованная консолидация),
mood-как-нейромодулятор, и CLS-разделение быстрый/медленный (у нас грубо есть как слои
session$\to$life).
\end{principle}

\section{Квантовый слой: честно}

Соблазн прямой: ассоциативная память $=$ волновая динамика, кванты $=$ волновая механика,
значит кванты идеальны. Не совсем. \emph{Чистая} квантовая динамика унитарна
($|\psi(t)\rangle=U(t)|\psi(0)\rangle$): обратима, сохраняет норму, \emph{не забывает} и
\emph{не} схлопывается в аттрактор ($\|Ux-Uy\|=\|x-y\|$ — нет сжатия). А ассоциативной памяти
нужно ровно сжатие: много близких состояний $\to$ одно воспоминание. Поэтому квантовая
ассоциативная память неизбежно тащит \emph{диссипацию}, измерения, связь со средой,
open-system-динамику. Второй подвох: $2^n$ амплитуд $\ne$ $2^n$ \emph{доступных} чисел —
измерение выдаёт ограниченную классическую информацию, а классический вход ещё надо
\emph{загрузить} в квантовое состояние. Поэтому экспоненциальное гильбертово пространство не
даёт автоматически экспоненциальной полезной памяти.

\begin{openq}[Оценка кванта: исследовательски высоко, практически сейчас низко]
Область quantum reservoir computing и quantum associative memory реально существует и активно
развивается (спиновые сети, фотоника на сотнях мод, quantum Hopfield на малых NISQ-устройствах).
\emph{Дисклеймер о верифицируемости}: конкретные свежие эксперименты и цифры «превосходства»,
всплывшие в обсуждении, лежат за горизонтом надёжной проверки этого текста — их стоит держать
как \emph{иллюстрацию направления}, а не как установленный quantum advantage; сам advantage для
памяти пока не доказан. Разумная ближняя форма — \emph{не} «память Рин $=$ квантовый компьютер»,
а гибрид: классический гиппокамп $\to$ \emph{маленький} квантовый резервуар как чрезвычайно
богатый нелинейный feature-generator $\to$ классическое пластичное хранилище. Квант тут — не
хранилище на двадцать лет жизни, а мощный возбудитель одного оператора.
\end{openq}

\section{Вердикт и как оставить дверь открытой}

Для Рин с её требованиями (редактировать, забывать, объяснять источник, поддерживать
истинность) \textbf{символический типизированный граф остаётся хребтом}, а динамический слой —
это \emph{дополняющая} ассоциативная надстройка («что я сейчас помню») и отношения
\emph{учитель\,--\,ученик} (граф $=$ архив/учитель). Радикальную версию «убрать семантический
граф» для наших требований отклоняем — с конкретными причинами (потеря забывания/провенанса/
TMS), а не из робости. Но инженерно правильно \textbf{держать динамический слой за чистым
интерфейсом}: тогда классический решатель $R_{\text{classical}}$ можно однажды заменить на
$R_{\text{oscillator}}/R_{\text{photonic}}/R_{\text{quantum}}$, не переписывая остальное.

\begin{openq}[Фальсифицируемые эксперименты, если идти в эту ветку]
Три чистых вопроса, которые можно поставить на стенде, не ломая текущую систему:
\emph{(1)} сравнить $e^{-tL}$, $\cos(t\sqrt L)$ и $e^{-itL}$ на \emph{одной} topologии
memory-графа — какой профиль лучше поднимает правильную окрестность (diffusion vs wave vs
quantum-walk memory); \emph{(2)} даёт ли \emph{пластичная} надстройка над типизированным графом
(Hebbian-усиление со-активированных рёбер) измеримо лучший recall, чем статические веса, — и не
ценой ли катастрофической интерференции; \emph{(3)} readout-контроль: посадить нарочно тупой
декодер и проверить, есть ли вообще содержательная информация в динамике, прежде чем строить
сложный мост к LLM. Все три — на golden-микромире (гл.~\ref{ch:experiment}), где есть ground
truth.
\end{openq}


\backmatter
\chapter{Литература}
\label{ch:references}

Ниже — работы, на которые опирается монография, в согласованном формате. Для
препринтов и внутренних материалов проекта (библиографии волтов RIN и RIN–DTI) в
конце записи в скобках указан локатор во внутренней библиографии; это указание на
местоположение источника, а не замена выходных данных.

\begin{thebibliography}{99}
\small

\bibitem{anderson2004} J.~R. Anderson, D.~Bothell, M.~D. Byrne и др.
An Integrated Theory of the Mind. \emph{Psychological Review}, 111(4), 2004.
(Базовая активация ACT-R; см. также Anderson \& Schooler, 1991.)

\bibitem{collins1975} A.~M. Collins, E.~F. Loftus.
A Spreading-Activation Theory of Semantic Processing.
\emph{Psychological Review}, 82(6), 1975.

\bibitem{hebb1949} D.~O. Hebb. \emph{The Organization of Behavior}. Wiley, 1949.

\bibitem{hopfield1982} J.~J. Hopfield.
Neural Networks and Physical Systems with Emergent Collective Computational
Abilities. \emph{PNAS}, 79(8), 1982.

\bibitem{park2023} J.~S. Park, J.~C. O'Brien, C.~J. Cai и др.
Generative Agents: Interactive Simulacra of Human Behavior. \emph{UIST}, 2023.

\bibitem{toy2024} J.~Toy, J.~MacAdam, P.~Tabor.
Metacognition is all you need? Using Introspection in Generative Agents to Improve
Goal-directed Behavior. arXiv:2401.10910, 2024.

\bibitem{kahneman2011} D.~Kahneman. \emph{Thinking, Fast and Slow}. Macmillan, 2011.
(System~1 / System~2.)

\bibitem{nisbett1977} R.~E. Nisbett, T.~D. Wilson.
Telling More Than We Can Know: Verbal Reports on Mental Processes.
\emph{Psychological Review}, 84(3), 1977. (Конфабуляция при самообъяснении.)

\bibitem{doyle1979} J.~Doyle. A Truth Maintenance System.
\emph{Artificial Intelligence}, 12(3), 1979.

\bibitem{dekleer1986} J.~de~Kleer. An Assumption-Based TMS.
\emph{Artificial Intelligence}, 28(2), 1986.

\bibitem{quine1951} W.~V.~O. Quine. Two Dogmas of Empiricism.
\emph{The Philosophical Review}, 60(1), 1951. (Тезис недоопределённости
Дюэма–Куайна.)

\bibitem{ross1975} L.~Ross, M.~R. Lepper, M.~Hubbard.
Perseverance in Self-Perception and Social Perception.
\emph{J. Personality and Social Psychology}, 32(5), 1975. (Belief perseverance.)

\bibitem{bak1987} P.~Bak, C.~Tang, K.~Wiesenfeld.
Self-Organized Criticality. \emph{Physical Review Letters}, 59(4), 1987.

\bibitem{beggs2003} J.~M. Beggs, D.~Plenz.
Neuronal Avalanches in Neocortical Circuits. \emph{J. Neuroscience}, 23(35), 2003.

\bibitem{page1999} L.~Page, S.~Brin, R.~Motwani, T.~Winograd.
The PageRank Citation Ranking: Bringing Order to the Web.
Stanford InfoLab Technical Report, 1999.

\bibitem{tong2006} H.~Tong, C.~Faloutsos, J.-Y. Pan.
Fast Random Walk with Restart and Its Applications. \emph{ICDM}, 2006.

\bibitem{spielman2019} D.~A. Spielman.
\emph{Spectral and Algebraic Graph Theory} (лекционные заметки), Yale University.

\bibitem{ortega2018} A.~Ortega, P.~Frossard, J.~Kova\v{c}evi\'c и др.
Graph Signal Processing: Overview, Challenges, and Applications.
\emph{Proceedings of the IEEE}, 106(5), 2018.

\bibitem{topping2022} J.~Topping, F.~Di~Giovanni, B.~P. Chamberlain и др.
Understanding Over-Squashing and Bottlenecks on Graphs via Curvature.
\emph{ICLR}, 2022. (Волт RIN–DTI, поз.~148.)

\bibitem{black2023} M.~Black, Z.~Wan, A.~Nayyeri, Y.~Wang.
Understanding Oversquashing in GNNs through the Lens of Effective Resistance.
\emph{ICML}, 2023. (Волт RIN–DTI, поз.~183.)

\bibitem{rusch2023} T.~K. Rusch, M.~M. Bronstein, S.~Mishra.
A Survey on Oversmoothing in Graph Neural Networks. arXiv:2303.10993, 2023.
(Волт RIN–DTI, поз.~170.)

\bibitem{bodnar2021cw} C.~Bodnar, F.~Frasca, N.~Otter и др.
Weisfeiler and Lehman Go Cellular: CW Networks. \emph{NeurIPS}, 2021.
(Волт RIN–DTI, поз.~185.)

\bibitem{bodnar2021mpsn} C.~Bodnar, F.~Frasca, Y.~Wang и др.
Weisfeiler and Lehman Go Topological: Message Passing Simplicial Networks.
\emph{ICML}, 2021. (Волт RIN–DTI, поз.~186.)

\bibitem{bodnar2022sheaf} C.~Bodnar, F.~Di~Giovanni, B.~P. Chamberlain и др.
Neural Sheaf Diffusion. \emph{NeurIPS}, 2022.

\bibitem{wei2020} R.~Wang, D.~D. Nguyen, G.-W. Wei.
Persistent Spectral Graph. \emph{Int. J. Numerical Methods in Biomedical
Engineering}, 2020. (TDA/персистентный лапласиан для биомолекул; ср. TopologyNet,
Cang \& Wei, 2017. Волт RIN–DTI, поз.~100, 103, 104, 144.)

\bibitem{occ1988} A.~Ortony, G.~L. Clore, A.~Collins.
\emph{The Cognitive Structure of Emotions}. Cambridge University Press, 1988.
(OCC-модель оценивания; ср. компонентная модель K.~Scherer.)

\bibitem{friston2010} K.~Friston. The Free-Energy Principle: A Unified Brain
Theory? \emph{Nature Reviews Neuroscience}, 11(2), 2010. (Ошибка предсказания как
сигнал значимости.)

\bibitem{speer2017} R.~Speer, J.~Chin, C.~Havasi.
ConceptNet 5.5: An Open Multilingual Graph of General Knowledge. \emph{AAAI},
2017. (Включая эмбеддинги Numberbatch.)

\bibitem{sap2019} M.~Sap, R.~Le~Bras, E.~Allaway и др.
ATOMIC: An Atlas of Machine Commonsense for If-Then Reasoning. \emph{AAAI}, 2019.

\bibitem{hwang2021} J.~D. Hwang, C.~Bhagavatula, R.~Le~Bras и др.
(Comet-)ATOMIC 2020: On Symbolic and Neural Commonsense Knowledge Graphs.
\emph{AAAI}, 2021.

\bibitem{west2022} P.~West, C.~Bhagavatula, J.~Hessel и др.
Symbolic Knowledge Distillation: from General Language Models to Commonsense
Models. \emph{NAACL}, 2022. (ATOMIC10x.)

\bibitem{vrandecic2014} D.~Vrande\v{c}i\'c, M.~Kr\"otzsch.
Wikidata: A Free Collaborative Knowledgebase. \emph{Communications of the ACM},
57(10), 2014.

\bibitem{xu2019} K.~Xu, W.~Hu, J.~Leskovec, S.~Jegelka.
How Powerful are Graph Neural Networks? \emph{ICLR}, 2019. (GIN; волт RIN–DTI,
поз.~68.)

\bibitem{corso2020} G.~Corso, L.~Cavalleri, D.~Beaini и др.
Principal Neighbourhood Aggregation for Graph Nets. \emph{NeurIPS}, 2020.
(Волт RIN–DTI, поз.~178.)

\bibitem{hu2021lora} E.~J. Hu, Y.~Shen, P.~Wallis и др.
LoRA: Low-Rank Adaptation of Large Language Models. arXiv:2106.09685, 2021.

\bibitem{rafailov2023} R.~Rafailov, A.~Sharma, E.~Mitchell и др.
Direct Preference Optimization. \emph{NeurIPS}, 2023.

\bibitem{zhou2006hg} D.~Zhou, J.~Huang, B.~Sch\"olkopf.
Learning with Hypergraphs: Clustering, Classification, and Embedding.
\emph{NeurIPS}, 2006. (Гиперграфовый лапласиан.)

\bibitem{lim2023hodge} L.-H. Lim. Hodge Laplacians on Graphs.
\emph{SIAM Review}, 62(3), 2020. (Ходж-лапласианы и высшие размерности.)

\bibitem{helamem2026} J.~Zhu, J.~Li, C.~Zhang и др.
HeLa-Mem: Hebbian Learning and Associative Memory for LLM Agents. Препринт, 2026.
(Внутренняя библиография волта RIN, поз.~53.)

\bibitem{rinvault} Библиография памяти LLM-агентов из волта RIN: APEX-MEM, TRACE,
SYNAPSE, GAAMA, MRMS, Zep/Graphiti, Persistent Sycophancy, LifeSide, LoCoMo и др.
(Внутренние материалы проекта; полные выходные данные — в
\texttt{99 Sources/1.~Bibliography (RIN)}.)

\bibitem{rinproject} Материалы проекта RIN: модульные системные промпты
(\texttt{prompts/global\_core/*}), persona-тесты и логи Stage~0
(\texttt{logs/rin\_stage0.jsonl}), Ринциклопедии
(\texttt{docs/rincyclopedia\_*}).

\end{thebibliography}


\end{document}
